%%%%%%%%%%%%%%%%%%%%%%%%%%%%%%%%%%%%%%%%%%%%%%%%%%%%%%%%%%%%%%%%%%%%%%%%%%%%%%%
%% TKS v7.3 QUANTUM TRACK — Quantum Noetics
%% Agent: tks-quantum
%%
%% This file contains the complete quantum-mechanical formalization of TKS,
%% including Hilbert space structure, operator theory, and measurement.
%%
%% DEPENDENCIES:
%%   - v7.0 Quantum Algebra (basic quantum structures)
%%   - v7.2 Noetic Hilbert Spaces (inner products, bounded operators)
%%   - v7.2 Spectral Theory (spectral decomposition)
%%   - v7.2 Quantum Fractals (preliminary)
%%
%% STATUS: SKELETON — To be filled by tks-quantum agent
%%%%%%%%%%%%%%%%%%%%%%%%%%%%%%%%%%%%%%%%%%%%%%%%%%%%%%%%%%%%%%%%%%%%%%%%%%%%%%%

%% ============================================================================
%% CHAPTER 0: HILBERT SPACE SEMANTICS FOR IDEAS
%% ============================================================================
\chapter{Hilbert Space Semantics for Ideas}
\label{ch:hilbert-semantics-ideas}

\begin{tcolorbox}[colback=blue!5!white,colframe=blue!75!black,title=\vnewthree]
This chapter establishes the foundational quantum-semantic framework for TKS by constructing a Hilbert space $\Hspace_{\Ideas}$ of Ideas, embedding classical Idea states into this space, and interpreting the inner product structure in terms of semantic similarity and interference. This provides the basis for all subsequent quantum noetic constructions in v7.3.
\end{tcolorbox}

%% ----------------------------------------------------------------------------
\section{The Base Hilbert Space of Ideas}
\label{sec:base-hilbert-ideas}

We begin by constructing the abstract Hilbert space $\Hspace_{\Ideas}$ that serves as the quantum-semantic arena for Ideas. This construction generalizes the finite-dimensional $\Hspace_\nu$ of v7.2 to accommodate arbitrary Idea spaces while maintaining compatibility with the classical TKS ontology.

\begin{definition}[Idea Configuration Space]
\label{def:idea-config-space}
Let $\Ideas$ denote the classical set of Ideas as established in v6.1. The \textbf{Idea configuration space} is the complex vector space:
\[
V_{\Ideas} = \mathrm{span}_{\mathbb{C}}\{\delta_I \mid I \in \Ideas\}
\]
where $\delta_I$ denotes the formal basis vector associated with Idea $I$. For finite $|\Ideas| = N$, we have $V_{\Ideas} \cong \mathbb{C}^N$.
\end{definition}

\begin{definition}[Base Hilbert Space of Ideas]
\label{def:base-hilbert-ideas}
The \textbf{base Hilbert space of Ideas} $\Hspace_{\Ideas}$ is the completion of $V_{\Ideas}$ with respect to the inner product:
\[
\langle \delta_I, \delta_J \rangle_{\Hspace} = \delta_{IJ}
\]
extended sesquilinearly (linear in the second argument, conjugate-linear in the first). For the standard TKS setting with $\Ideas = \{Wn \mid W \in \Worlds, n \in \Elements\}$:
\[
\Hspace_{\Ideas} = \ell^2(\Ideas) \cong \mathbb{C}^{40}
\]
recovering the noetic Hilbert space $\Hspace_\nu$ of v7.2.
\end{definition}

\begin{notation}[Dirac Notation]
\label{not:dirac}
Following quantum-mechanical convention, we write:
\begin{align}
\qket{I} &\equiv \delta_I \quad \text{(ket: column vector)} \\
\qbra{I} &\equiv \delta_I^\dagger \quad \text{(bra: row vector)} \\
\qbraket{I}{J} &\equiv \langle \delta_I, \delta_J \rangle_{\Hspace} = \delta_{IJ}
\end{align}
The set $\{\qket{I}\}_{I \in \Ideas}$ forms an orthonormal basis for $\Hspace_{\Ideas}$.
\end{notation}

\begin{definition}[Quantum Idea State]
\label{def:quantum-idea-state}
A \textbf{quantum Idea state} is a unit vector $\qket{\psi} \in \Hspace_{\Ideas}$ with $\|\qket{\psi}\|^2 = 1$. In the Idea basis:
\[
\qket{\psi} = \sum_{I \in \Ideas} c_I \qket{I}, \quad \text{where } \sum_{I \in \Ideas} |c_I|^2 = 1
\]
The complex coefficients $c_I \in \mathbb{C}$ are called \textbf{Idea amplitudes}. The \textbf{projective Idea space} is:
\[
\mathbb{P}\Hspace_{\Ideas} = (\Hspace_{\Ideas} \setminus \{0\}) / \mathbb{C}^*
\]
where physical states correspond to rays (equivalence classes under scalar multiplication).
\end{definition}

\begin{remark}[Superposition Principle]
\label{rem:superposition}
A quantum Idea state $\qket{\psi}$ with multiple nonzero amplitudes $c_I$ represents a coherent superposition of Ideas. Unlike classical probability distributions (which also sum to 1), the amplitudes $c_I$ are complex-valued, enabling interference effects that distinguish quantum from classical semantics.
\end{remark}

%% ----------------------------------------------------------------------------
\section{Embedding of Classical Idea States}
\label{sec:classical-embedding}

The quantum formalism must be a refinement, not a replacement, of classical TKS semantics. We now establish how classical Idea states embed into $\Hspace_{\Ideas}$.

\begin{definition}[Classical Idea Embedding]
\label{def:classical-embedding}
The \textbf{classical embedding map} $\iota : \Ideas \hookrightarrow \Hspace_{\Ideas}$ is defined by:
\[
\iota(I) = \qket{I}
\]
This embeds each classical Idea $I$ as the corresponding basis state in $\Hspace_{\Ideas}$.
\end{definition}

\begin{proposition}[Embedding Properties]
\label{prop:embedding-props}
The classical embedding satisfies:
\begin{enumerate}
  \item \textbf{Injectivity}: $\iota(I) = \iota(J) \implies I = J$
  \item \textbf{Orthonormality Preservation}: $\qbraket{\iota(I)}{\iota(J)} = \delta_{IJ}$
  \item \textbf{Norm Preservation}: $\|\iota(I)\| = 1$ for all $I \in \Ideas$
  \item \textbf{Span}: $\mathrm{span}_{\mathbb{C}}(\mathrm{Im}(\iota)) = \Hspace_{\Ideas}$
\end{enumerate}
\end{proposition}

\begin{proof}
Properties (1)--(3) follow directly from the orthonormality of the basis $\{\qket{I}\}$. Property (4) holds by construction of $\Hspace_{\Ideas}$ as the span-completion of these basis vectors.
\end{proof}

\begin{definition}[Classical-Quantum Correspondence]
\label{def:cq-correspondence}
The \textbf{classical-quantum correspondence} for Ideas is the commutative diagram:
\[
\begin{tikzcd}
\Ideas \arrow[r, "\iota"] \arrow[d, "\sem{\cdot}_{\mathrm{cl}}"'] & \Hspace_{\Ideas} \arrow[d, "\sem{\cdot}_{\mathrm{qu}}"] \\
\sem{\Ideas}_{\mathrm{cl}} \arrow[r, "\tilde{\iota}"'] & \sem{\Hspace_{\Ideas}}_{\mathrm{qu}}
\end{tikzcd}
\]
where $\sem{\cdot}_{\mathrm{cl}}$ denotes classical denotational semantics (v7.1) and $\sem{\cdot}_{\mathrm{qu}}$ denotes quantum semantics. The bottom map $\tilde{\iota}$ extends the embedding to semantic domains.
\end{definition}

\begin{theorem}[Classical Limit Recovery]
\label{thm:quantum-classical-limit}
For any classical observable $O : \Ideas \to \mathbb{R}$ and its quantum extension $\qop{O} : \Hspace_{\Ideas} \to \Hspace_{\Ideas}$, the expectation value on an embedded classical state recovers the classical value:
\[
\qbra{I} \qop{O} \qket{I} = O(I)
\]
for all $I \in \Ideas$.
\end{theorem}

\begin{proof}
Define $\qop{O} = \sum_{I \in \Ideas} O(I) \qket{I}\qbra{I}$ (the diagonal extension of $O$ to an operator). Then:
\[
\qbra{I} \qop{O} \qket{I} = \sum_{J \in \Ideas} O(J) \qbra{I} \qket{J}\qbra{J} \qket{I} = \sum_{J} O(J) \delta_{IJ}^2 = O(I)
\]
\end{proof}

\begin{definition}[World-Stratified Embedding]
\label{def:world-stratified}
For the standard TKS Idea space $\Ideas = \{Wn \mid W \in \Worlds, n \in \{1,\ldots,10\}\}$, the embedding respects the world structure:
\[
\iota_{W} : \Ideas_W \hookrightarrow \Hspace_W \subset \Hspace_{\Ideas}
\]
where $\Hspace_W = \mathrm{span}\{\qket{Wn} \mid n = 1,\ldots,10\}$ and:
\[
\Hspace_{\Ideas} = \bigoplus_{W \in \Worlds} \Hspace_W = \Hspace_A \oplus \Hspace_B \oplus \Hspace_C \oplus \Hspace_D
\]
This decomposition is orthogonal: $\Hspace_W \perp \Hspace_{W'}$ for $W \neq W'$.
\end{definition}

%% ----------------------------------------------------------------------------
\section{Inner Product as Similarity and Interference}
\label{sec:inner-product-interpretation}

The inner product structure of $\Hspace_{\Ideas}$ admits a rich semantic interpretation that goes beyond classical set-theoretic notions of similarity.

\begin{definition}[Idea Overlap]
\label{def:idea-overlap}
For quantum Idea states $\qket{\psi}, \qket{\phi} \in \Hspace_{\Ideas}$, their \textbf{overlap amplitude} is:
\[
\alpha(\psi, \phi) = \qbraket{\psi}{\phi} \in \mathbb{C}
\]
The \textbf{transition probability} (or \textbf{similarity measure}) is:
\[
P(\psi \to \phi) = |\qbraket{\psi}{\phi}|^2 \in [0,1]
\]
\end{definition}

\begin{proposition}[Similarity Properties]
\label{prop:similarity-props}
The transition probability $P(\psi \to \phi)$ satisfies:
\begin{enumerate}
  \item \textbf{Reflexivity}: $P(\psi \to \psi) = 1$
  \item \textbf{Symmetry}: $P(\psi \to \phi) = P(\phi \to \psi)$
  \item \textbf{Boundedness}: $0 \leq P(\psi \to \phi) \leq 1$
  \item \textbf{Orthogonality}: $P(\psi \to \phi) = 0$ iff $\qket{\psi} \perp \qket{\phi}$
  \item \textbf{Classical Discreteness}: For embedded classical Ideas, $P(I \to J) = \delta_{IJ}$
\end{enumerate}
\end{proposition}

\begin{proof}
Properties (1)--(4) follow from properties of the inner product and Cauchy--Schwarz inequality. Property (5) is the orthonormality of the classical basis.
\end{proof}

\begin{definition}[Semantic Interference]
\label{def:semantic-interference}
Consider a quantum Idea state $\qket{\psi} = \sum_I c_I \qket{I}$ and a measurement in a different basis $\{\qket{\phi_j}\}_j$. The probability of outcome $j$ is:
\[
P(j | \psi) = |\qbraket{\phi_j}{\psi}|^2 = \left| \sum_I c_I \qbraket{\phi_j}{I} \right|^2
\]
This exhibits \textbf{interference}: the cross-terms $c_I^* c_J \qbra{I}\phi_j\rangle \qbraket{\phi_j}{J}$ for $I \neq J$ can be positive or negative, leading to constructive or destructive interference of Idea amplitudes.
\end{definition}

\begin{example}[Superposition and Interference]
\label{ex:superposition-interference}
Consider the equal superposition of two Ideas $I_1, I_2 \in \Ideas$:
\[
\qket{\psi_+} = \frac{1}{\sqrt{2}}(\qket{I_1} + \qket{I_2}), \quad
\qket{\psi_-} = \frac{1}{\sqrt{2}}(\qket{I_1} - \qket{I_2})
\]
These are orthogonal: $\qbraket{\psi_+}{\psi_-} = \frac{1}{2}(1 - 1) = 0$, despite both having the same marginal distribution over $\{I_1, I_2\}$. The relative phase distinguishes them quantum-mechanically.
\end{example}

\begin{definition}[Noetic Fidelity]
\label{def:noetic-fidelity}
The \textbf{noetic fidelity} between quantum Idea states is:
\[
F(\psi, \phi) = |\qbraket{\psi}{\phi}|
\]
For mixed states (density operators), this generalizes to:
\[
F(\rho, \sigma) = \Trace\sqrt{\sqrt{\rho}\sigma\sqrt{\rho}}
\]
Fidelity measures how well one Idea state can simulate another under quantum operations.
\end{definition}

\begin{theorem}[Fidelity and Distinguishability]
\label{thm:fidelity-distinguish}
Two quantum Idea states $\qket{\psi}$ and $\qket{\phi}$ are:
\begin{enumerate}
  \item \textbf{Perfectly distinguishable} (by a single measurement) iff $F(\psi,\phi) = 0$
  \item \textbf{Identical} (up to global phase) iff $F(\psi,\phi) = 1$
  \item \textbf{Partially distinguishable} for $0 < F(\psi,\phi) < 1$
\end{enumerate}
The optimal probability of correctly distinguishing $\qket{\psi}$ from $\qket{\phi}$ (given with equal prior) is:
\[
P_{\mathrm{correct}} = \frac{1}{2}\left(1 + \sqrt{1 - F(\psi,\phi)^2}\right)
\]
\end{theorem}

\begin{proof}
This follows from the Helstrom bound in quantum state discrimination theory. For pure states, optimal discrimination uses the POVM defined by the eigenprojectors of $\qket{\psi}\qbra{\psi} - \qket{\phi}\qbra{\phi}$.
\end{proof}

\begin{definition}[Idea Distance]
\label{def:idea-distance}
The \textbf{Bures distance} between quantum Idea states is:
\[
d_B(\psi, \phi) = \sqrt{2(1 - F(\psi, \phi))} = \sqrt{2(1 - |\qbraket{\psi}{\phi}|)}
\]
This is a proper metric on $\mathbb{P}\Hspace_{\Ideas}$ satisfying the triangle inequality.
\end{definition}

\begin{remark}[Geometric Interpretation]
\label{rem:geometric}
The projective space $\mathbb{P}\Hspace_{\Ideas}$ equipped with the Fubini--Study metric has a natural Riemannian structure. The geodesic distance between states $\qket{\psi}$ and $\qket{\phi}$ is:
\[
d_{FS}(\psi, \phi) = \arccos|\qbraket{\psi}{\phi}|
\]
This geometric structure allows us to speak of ``nearby'' Ideas, continuous paths of Idea transformation, and the curvature of Idea space.
\end{remark}

%% ----------------------------------------------------------------------------
\section{Noetic Operators on the Idea Hilbert Space}
\label{sec:noetic-ops-idea-hilbert}

Having established $\Hspace_{\Ideas}$, we now connect classical noetic operators to bounded linear operators, preparing for the full quantum noetic theory of Chapter~\ref{ch:quantum-noetic-operators}.

\begin{definition}[Noetic Operator Lifting]
\label{def:noetic-lifting}
For each classical noetic $\noe{k} : \Ideas \to \Ideas$, define the \textbf{lifted noetic operator} $\qop{\nu}_k : \Hspace_{\Ideas} \to \Hspace_{\Ideas}$ by:
\[
\qop{\nu}_k = \sum_{I \in \Ideas} \qket{\mathfrak{n}_k(I)}\qbra{I}
\]
where $\mathfrak{n}_k(I)$ denotes the classical action of $\noe{k}$ on Idea $I$.
\end{definition}

\begin{theorem}[Noetic Operators are Bounded]
\label{thm:noetic-bounded}
For each noetic $\noe{k}$, the lifted operator $\qop{\nu}_k \in \BoundedOp(\Hspace_{\Ideas})$ is a bounded linear operator with:
\[
\|\qop{\nu}_k\|_{\mathrm{op}} \leq 1
\]
Moreover, if $\mathfrak{n}_k : \Ideas \to \Ideas$ is a bijection, then $\qop{\nu}_k$ is unitary: $\qop{\nu}_k^\dagger \qop{\nu}_k = \qop{\nu}_k \qop{\nu}_k^\dagger = \mathbf{I}$.
\end{theorem}

\begin{proof}
\textbf{Boundedness}: Let $\qket{\psi} = \sum_I c_I \qket{I} \in \Hspace_{\Ideas}$ with $\|\psi\|^2 = \sum_I |c_I|^2 = 1$. Then:
\begin{align}
\qop{\nu}_k \qket{\psi} &= \sum_I c_I \qket{\mathfrak{n}_k(I)}
\end{align}
We compute the norm:
\[
\|\qop{\nu}_k \qket{\psi}\|^2 = \sum_J \left| \sum_{I : \mathfrak{n}_k(I) = J} c_I \right|^2
\]
By the triangle inequality and Cauchy--Schwarz:
\[
\left| \sum_{I : \mathfrak{n}_k(I) = J} c_I \right|^2 \leq \left(\sum_{I : \mathfrak{n}_k(I) = J} |c_I|\right)^2 \leq m_J \sum_{I : \mathfrak{n}_k(I) = J} |c_I|^2
\]
where $m_J = |\mathfrak{n}_k^{-1}(J)|$ is the multiplicity. Summing over $J$:
\[
\|\qop{\nu}_k \qket{\psi}\|^2 \leq \max_J(m_J) \cdot \|\psi\|^2
\]
Since $\Ideas$ is finite (40 elements in standard TKS), $m_J \leq |\Ideas|$, so $\|\qop{\nu}_k\|_{\mathrm{op}} < \infty$. For permutation noetics where each $m_J = 1$, we get $\|\qop{\nu}_k\|_{\mathrm{op}} = 1$.

\textbf{Unitarity for bijections}: If $\mathfrak{n}_k$ is a bijection, then:
\[
\qop{\nu}_k^\dagger = \sum_I \qket{I}\qbra{\mathfrak{n}_k(I)} = \sum_J \qket{\mathfrak{n}_k^{-1}(J)}\qbra{J}
\]
Computing:
\[
\qop{\nu}_k^\dagger \qop{\nu}_k = \sum_{I,J} \qket{I}\qbraket{\mathfrak{n}_k(I)}{\mathfrak{n}_k(J)}\qbra{J} = \sum_I \qket{I}\qbra{I} = \mathbf{I}
\]
Similarly for $\qop{\nu}_k \qop{\nu}_k^\dagger = \mathbf{I}$.
\end{proof}

\begin{corollary}[Noetic Algebra Embedding]
\label{cor:noetic-algebra-embedding}
The classical noetic algebra $\Noetics$ embeds into the C*-algebra $\BoundedOp(\Hspace_{\Ideas})$ via the lifting map $\noe{k} \mapsto \qop{\nu}_k$. This embedding preserves:
\begin{enumerate}
  \item Composition: $\qop{\nu}_{k \circ j} = \qop{\nu}_k \circ \qop{\nu}_j$
  \item Identity: $\qop{\nu}_0 = \mathbf{I}$ (where $\noe{0}$ is the identity noetic)
  \item Algebraic relations: Any identity satisfied by noetics in the classical algebra is preserved
\end{enumerate}
\end{corollary}

\begin{proof}
(1) follows from:
\[
\qop{\nu}_k \qop{\nu}_j = \sum_{I,J} \qket{\mathfrak{n}_k(J)}\qbraket{J}{\mathfrak{n}_j(I)}\qbra{I} = \sum_I \qket{\mathfrak{n}_k(\mathfrak{n}_j(I))}\qbra{I} = \qop{\nu}_{k \circ j}
\]
(2) and (3) follow by direct computation.
\end{proof}

\begin{remark}[Compatibility with v7.2]
\label{rem:v72-compat}
The construction here is fully compatible with the 40-dimensional noetic Hilbert space $\Hspace_\nu$ of v7.2 (Definition~\ref{def:noetic-hilbert} in v7.2). Setting $\Ideas = \{Wn \mid W \in \{A,B,C,D\}, n \in \{1,\ldots,10\}\}$ recovers $\Hspace_{\Ideas} = \Hspace_\nu$ and $\qop{\nu}_k$ as defined in v7.2 Definition~\ref{def:noetic-operator-quantum}.
\end{remark}

%% ----------------------------------------------------------------------------
%% HANDOFF NOTE
%% ----------------------------------------------------------------------------
\begin{tcolorbox}[colback=green!5!white,colframe=green!50!black,title=\textbf{Handoff: Next Steps},breakable]
This completes the foundational Hilbert space semantics for Ideas. The next section (Chapter~\ref{ch:quantum-noetic-operators}) should:
\begin{itemize}
  \item Develop the full operator algebra structure of noetic operators (C*-algebra properties, GNS representation)
  \item Analyze commutation relations $[\qop{\nu}_i, \qop{\nu}_j]$ and their physical/semantic interpretation
  \item Characterize which noetics are Hermitian (observables) vs. unitary (reversible transformations)
  \item Establish the spectral theory connecting eigenvalues to semantic fixed points
\end{itemize}
\end{tcolorbox}


%% ============================================================================
%% CHAPTER 1: QUANTUM NOETIC OPERATORS
%% ============================================================================
\chapter{Quantum Noetic Operators}
\label{ch:quantum-noetic-operators}

\begin{tcolorbox}[colback=blue!5!white,colframe=blue!75!black,title=\vnewthree]
This chapter provides the complete theory of noetic operators as bounded linear operators on the noetic Hilbert space, extending v7.2 foundations.
\end{tcolorbox}

%% ----------------------------------------------------------------------------
\section{Review: Noetic Hilbert Spaces from v7.2}
\label{sec:review-hilbert}

We briefly recall the essential structures established in v7.2 and Chapter~\ref{ch:hilbert-semantics-ideas}, which provide the foundation for the operator-algebraic development.

\begin{definition}[Noetic Hilbert Space (v7.2 Recap)]
\label{def:noetic-hilbert-recap}
The \textbf{noetic Hilbert space} $\Hspace_{\Ideas}$ is the 40-dimensional complex Hilbert space:
\[
\Hspace_{\Ideas} = \mathbb{C}^{40} = \mathrm{span}_{\mathbb{C}}\{\qket{Wn} \mid W \in \{A,B,C,D\}, n \in \{1,\ldots,10\}\}
\]
with inner product $\qbraket{Wn}{W'n'} = \delta_{WW'}\delta_{nn'}$. The orthogonal world decomposition is:
\[
\Hspace_{\Ideas} = \Hspace_A \oplus \Hspace_B \oplus \Hspace_C \oplus \Hspace_D
\]
where each $\Hspace_W$ is 10-dimensional.
\end{definition}

\begin{remark}[Completeness]
\label{rem:completeness}
As a finite-dimensional Hilbert space, $\Hspace_{\Ideas}$ is automatically complete. This ensures that limits of Cauchy sequences of operators exist, which is essential for the C*-algebra structure developed below.
\end{remark}

%% ----------------------------------------------------------------------------
\section{The C*-Algebra of Noetic Operators}
\label{sec:cstar-algebra}

We now develop the C*-algebraic framework for noetic operators, providing the operator-algebraic foundation for quantum TKS semantics.

\begin{definition}[Noetic Operator Set]
\label{def:noetic-op-set}
The \textbf{noetic operator set} is:
\[
\mathcal{S}_\nu = \{\qop{\nu}_0, \qop{\nu}_1, \ldots, \qop{\nu}_9\} \subset \BoundedOp(\Hspace_{\Ideas})
\]
where each $\qop{\nu}_k$ is the lifted noetic operator from Definition~\ref{def:noetic-lifting}.
\end{definition}

\begin{definition}[Noetic C*-Algebra]
\label{def:noetic-cstar}
The \textbf{noetic C*-algebra} $\mathcal{A}_\nu$ is the smallest C*-subalgebra of $\BoundedOp(\Hspace_{\Ideas})$ containing $\mathcal{S}_\nu$:
\[
\mathcal{A}_\nu = C^*(\mathcal{S}_\nu) = \overline{\mathrm{Alg}(\mathcal{S}_\nu \cup \mathcal{S}_\nu^*)}^{\|\cdot\|}
\]
where:
\begin{itemize}
  \item $\mathcal{S}_\nu^* = \{\qop{\nu}_k^\dagger \mid k = 0,\ldots,9\}$ is the set of adjoints
  \item $\mathrm{Alg}(\cdot)$ denotes the algebra generated (polynomials in the operators)
  \item The overline denotes norm closure
\end{itemize}
\end{definition}

\begin{proposition}[C*-Algebra Properties]
\label{prop:cstar-props}
The noetic C*-algebra $\mathcal{A}_\nu$ satisfies:
\begin{enumerate}
  \item \textbf{Unital}: $\mathbf{I} = \qop{\nu}_0 \in \mathcal{A}_\nu$
  \item \textbf{*-closed}: $A \in \mathcal{A}_\nu \implies A^\dagger \in \mathcal{A}_\nu$
  \item \textbf{Norm-closed}: If $\{A_n\} \subset \mathcal{A}_\nu$ and $\|A_n - A\| \to 0$, then $A \in \mathcal{A}_\nu$
  \item \textbf{C*-identity}: $\|A^\dagger A\| = \|A\|^2$ for all $A \in \mathcal{A}_\nu$
\end{enumerate}
\end{proposition}

\begin{proof}
Properties (1)--(3) hold by construction. Property (4) is inherited from $\BoundedOp(\Hspace_{\Ideas})$, which is a C*-algebra.
\end{proof}

\begin{definition}[Involution Structure]
\label{def:involution}
The \textbf{involution} (adjoint map) $* : \mathcal{A}_\nu \to \mathcal{A}_\nu$ is defined by $A^* = A^\dagger$. For the generators:
\[
\qop{\nu}_k^* = \qop{\nu}_k^\dagger = \sum_{I \in \Ideas} \qket{I}\qbra{\mathfrak{n}_k(I)}
\]
The involution satisfies:
\begin{enumerate}
  \item $(A^*)^* = A$
  \item $(AB)^* = B^* A^*$
  \item $(\alpha A + \beta B)^* = \bar{\alpha} A^* + \bar{\beta} B^*$
  \item $\|A^* A\| = \|A\|^2$
\end{enumerate}
\end{definition}

\begin{theorem}[Finite Generation]
\label{thm:finite-generation}
The noetic C*-algebra is finitely generated:
\[
\mathcal{A}_\nu = C^*(\qop{\nu}_1, \ldots, \qop{\nu}_9)
\]
(noting $\qop{\nu}_0 = \mathbf{I}$ is automatic in unital algebras). Moreover:
\[
\dim(\mathcal{A}_\nu) \leq 40^2 = 1600
\]
with equality iff $\mathcal{A}_\nu = \BoundedOp(\Hspace_{\Ideas}) \cong M_{40}(\mathbb{C})$.
\end{theorem}

\begin{proof}
The bound follows from $\mathcal{A}_\nu \subseteq \BoundedOp(\Hspace_{\Ideas}) \cong M_{40}(\mathbb{C})$, which has dimension $40^2$. The actual dimension depends on linear independence of products of noetics.
\end{proof}

\begin{definition}[Algebraic vs. Topological Generators]
\label{def:generators}
\textbf{Algebraic generators}: $\mathcal{A}_\nu^{\mathrm{alg}} = \mathrm{Alg}(\mathcal{S}_\nu \cup \mathcal{S}_\nu^*)$ consists of finite polynomials:
\[
A = \sum_{\text{finite}} c_{i_1 \cdots i_m j_1 \cdots j_n} \qop{\nu}_{i_1} \cdots \qop{\nu}_{i_m} \qop{\nu}_{j_1}^\dagger \cdots \qop{\nu}_{j_n}^\dagger
\]

\textbf{Topological closure}: On finite-dimensional $\Hspace_{\Ideas}$, $\mathcal{A}_\nu^{\mathrm{alg}} = \mathcal{A}_\nu$ (no completion needed).
\end{definition}

\begin{proposition}[Relationship to Full Operator Algebra]
\label{prop:full-algebra-relation}
The inclusion $\mathcal{A}_\nu \hookrightarrow \BoundedOp(\Hspace_{\Ideas})$ is:
\begin{enumerate}
  \item \textbf{Proper} if some operators in $\BoundedOp(\Hspace_{\Ideas})$ cannot be expressed as polynomials in noetics
  \item \textbf{Surjective} ($\mathcal{A}_\nu = \BoundedOp(\Hspace_{\Ideas})$) iff the noetics generate all of $M_{40}(\mathbb{C})$
\end{enumerate}
The surjectivity condition is: $\mathcal{S}_\nu$ has no common invariant subspace other than $\{0\}$ and $\Hspace_{\Ideas}$.
\end{proposition}

\begin{theorem}[Noetic Algebra Structure]
\label{thm:noetic-algebra-structure}
The noetic C*-algebra decomposes as:
\[
\mathcal{A}_\nu \cong \bigoplus_{i=1}^{r} M_{n_i}(\mathbb{C})
\]
for some integers $n_i$ with $\sum_i n_i^2 = \dim(\mathcal{A}_\nu)$. Each factor corresponds to an irreducible representation of the noetic algebra.
\end{theorem}

\begin{proof}
This follows from the Artin--Wedderburn theorem: every finite-dimensional semisimple algebra over $\mathbb{C}$ is a direct sum of matrix algebras. The noetic C*-algebra is semisimple (as a finite-dimensional C*-algebra).
\end{proof}

%% ----------------------------------------------------------------------------
\section{Noetic Commutators and Non-Commutativity}
\label{sec:noetic-commutators}

The commutation relations between noetic operators encode fundamental semantic relationships about the order-dependence of Idea transformations.

\begin{definition}[Noetic Commutator]
\label{def:noetic-commutator}
The \textbf{commutator} of noetic operators $\qop{\nu}_i$ and $\qop{\nu}_j$ is:
\[
[\qop{\nu}_i, \qop{\nu}_j] = \qop{\nu}_i \qop{\nu}_j - \qop{\nu}_j \qop{\nu}_i
\]
We say $\qop{\nu}_i$ and $\qop{\nu}_j$ \textbf{commute} if $[\qop{\nu}_i, \qop{\nu}_j] = 0$.
\end{definition}

\begin{proposition}[Commutator Matrix Elements]
\label{prop:commutator-elements}
The matrix elements of $[\qop{\nu}_i, \qop{\nu}_j]$ in the Idea basis are:
\[
\qbra{I}[\qop{\nu}_i, \qop{\nu}_j]\qket{J} = \delta_{I, \mathfrak{n}_i(\mathfrak{n}_j(J))} - \delta_{I, \mathfrak{n}_j(\mathfrak{n}_i(J))}
\]
Thus $[\qop{\nu}_i, \qop{\nu}_j] = 0$ iff $\mathfrak{n}_i \circ \mathfrak{n}_j = \mathfrak{n}_j \circ \mathfrak{n}_i$ as functions on $\Ideas$.
\end{proposition}

\begin{proof}
We compute:
\begin{align}
\qop{\nu}_i \qop{\nu}_j \qket{J} &= \qop{\nu}_i \qket{\mathfrak{n}_j(J)} = \qket{\mathfrak{n}_i(\mathfrak{n}_j(J))} \\
\qop{\nu}_j \qop{\nu}_i \qket{J} &= \qop{\nu}_j \qket{\mathfrak{n}_i(J)} = \qket{\mathfrak{n}_j(\mathfrak{n}_i(J))}
\end{align}
The difference vanishes on all basis states iff the classical compositions commute.
\end{proof}

\begin{definition}[Commuting Noetic Classes]
\label{def:commuting-classes}
Define the equivalence relation on $\{0,\ldots,9\}$:
\[
i \sim j \iff [\qop{\nu}_i, \qop{\nu}_k] = [\qop{\nu}_j, \qop{\nu}_k] \text{ for all } k
\]
The \textbf{commutant} of $\qop{\nu}_k$ is:
\[
\mathcal{C}(\qop{\nu}_k) = \{A \in \mathcal{A}_\nu \mid [A, \qop{\nu}_k] = 0\}
\]
\end{definition}

\begin{theorem}[Identity Noetic Commutativity]
\label{thm:identity-commutes}
The identity noetic $\qop{\nu}_0 = \mathbf{I}$ commutes with all operators:
\[
[\qop{\nu}_0, \qop{\nu}_k] = 0 \quad \text{for all } k \in \{0,\ldots,9\}
\]
\end{theorem}

\begin{proof}
$[\mathbf{I}, A] = \mathbf{I} A - A \mathbf{I} = A - A = 0$ for any operator $A$.
\end{proof}

\begin{proposition}[Idempotent Noetic Self-Commutation]
\label{prop:idem-self-commute}
For idempotent noetics ($\mathfrak{n}_k \circ \mathfrak{n}_k = \mathfrak{n}_k$), the operator $\qop{\nu}_k$ satisfies:
\[
[\qop{\nu}_k, \qop{\nu}_k] = 0 \quad \text{(trivially)}
\]
and more importantly, $\qop{\nu}_k^2 = \qop{\nu}_k$ (projector property from Theorem~\ref{thm:idem-projective} of v7.2).
\end{proposition}

\begin{definition}[Commutator Norm]
\label{def:commutator-norm}
The \textbf{commutator norm} measures the degree of non-commutativity:
\[
\kappa(i,j) = \|[\qop{\nu}_i, \qop{\nu}_j]\|_{\mathrm{op}}
\]
We have $\kappa(i,j) = 0$ iff noetics $i$ and $j$ commute.
\end{definition}

\begin{theorem}[Non-Commutativity and Semantic Order-Dependence]
\label{thm:noncomm-semantic}
For noetic operators $\qop{\nu}_i, \qop{\nu}_j$ with $[\qop{\nu}_i, \qop{\nu}_j] \neq 0$:
\begin{enumerate}
  \item There exists at least one Idea $I \in \Ideas$ such that:
  \[
  \mathfrak{n}_i(\mathfrak{n}_j(I)) \neq \mathfrak{n}_j(\mathfrak{n}_i(I))
  \]
  \item The order of applying noetic transformations matters for the final Idea state
  \item For quantum superposition states $\qket{\psi}$, the expectation values differ:
  \[
  \qbra{\psi}\qop{\nu}_i \qop{\nu}_j\qket{\psi} \neq \qbra{\psi}\qop{\nu}_j \qop{\nu}_i\qket{\psi}
  \]
  in general
\end{enumerate}
\end{theorem}

\begin{proof}
(1) follows from Proposition~\ref{prop:commutator-elements}: if $[\qop{\nu}_i, \qop{\nu}_j] \neq 0$, some matrix element is nonzero, so the classical compositions differ on some $I$.

(2) is a restatement of (1) in semantic terms.

(3) follows from the non-vanishing of $[\qop{\nu}_i, \qop{\nu}_j]$; choose $\qket{\psi}$ with nonzero projection onto the subspace where the commutator acts nontrivially.
\end{proof}

\begin{remark}[Physical Interpretation of Non-Commutativity]
\label{rem:noncomm-interp}
Non-commuting noetics represent \textbf{incompatible transformations}: applying them in different orders yields different results. This is the quantum-semantic analogue of non-commuting observables in physics, where simultaneous sharp measurement is impossible. In TKS:
\begin{itemize}
  \item \textbf{Commuting noetics}: The transformations can be ``performed simultaneously'' without order-dependence; they have a common eigenbasis
  \item \textbf{Non-commuting noetics}: Represent fundamentally different modes of Idea transformation that interfere with each other
\end{itemize}
\end{remark}

\begin{definition}[Noetic Lie Algebra]
\label{def:noetic-lie}
The \textbf{noetic Lie algebra} $\mathfrak{g}_\nu$ is the real vector space spanned by anti-Hermitian combinations:
\[
\mathfrak{g}_\nu = \mathrm{span}_{\mathbb{R}}\{i(\qop{\nu}_k - \qop{\nu}_k^\dagger), \qop{\nu}_k + \qop{\nu}_k^\dagger \mid k = 1,\ldots,9\}
\]
equipped with the Lie bracket $[A, B] = AB - BA$.
\end{definition}

%% ----------------------------------------------------------------------------
\section{Unitarity and Hermiticity}
\label{sec:unitarity-hermiticity}

The distinction between unitary and Hermitian noetic operators corresponds to the difference between reversible transformations and observable quantities.

\begin{definition}[Hermitian Noetic Operators]
\label{def:hermitian-noetics}
A noetic operator $\qop{\nu}_k$ is \textbf{Hermitian} (self-adjoint) if:
\[
\qop{\nu}_k^\dagger = \qop{\nu}_k
\]
Equivalently, $\qbra{I}\qop{\nu}_k\qket{J} = \overline{\qbra{J}\qop{\nu}_k\qket{I}}$ for all $I, J \in \Ideas$.
\end{definition}

\begin{proposition}[Hermiticity Condition]
\label{prop:hermiticity-cond}
The lifted noetic operator $\qop{\nu}_k$ is Hermitian iff for all $I, J \in \Ideas$:
\[
\mathfrak{n}_k(I) = J \iff \mathfrak{n}_k(J) = I
\]
That is, the classical noetic $\mathfrak{n}_k$ is a \textbf{symmetric relation} (equivalently, an involution up to fixed points).
\end{proposition}

\begin{proof}
We have:
\[
\qop{\nu}_k^\dagger = \sum_I \qket{I}\qbra{\mathfrak{n}_k(I)}
\]
For Hermiticity, we need $\qop{\nu}_k^\dagger = \qop{\nu}_k$, i.e.:
\[
\sum_I \qket{I}\qbra{\mathfrak{n}_k(I)} = \sum_J \qket{\mathfrak{n}_k(J)}\qbra{J}
\]
Comparing coefficients: $\qket{I}\qbra{\mathfrak{n}_k(I)} = \qket{\mathfrak{n}_k(I)}\qbra{I}$ for all $I$, which requires $\mathfrak{n}_k(I) = I$ (fixed point) or $\mathfrak{n}_k(\mathfrak{n}_k(I)) = I$ (involution).
\end{proof}

\begin{definition}[Unitary Noetic Operators]
\label{def:unitary-noetics}
A noetic operator $\qop{\nu}_k$ is \textbf{unitary} if:
\[
\qop{\nu}_k^\dagger \qop{\nu}_k = \qop{\nu}_k \qop{\nu}_k^\dagger = \mathbf{I}
\]
\end{definition}

\begin{theorem}[Unitarity Characterization]
\label{thm:unitarity-char}
The lifted noetic operator $\qop{\nu}_k$ is unitary iff the classical noetic $\mathfrak{n}_k : \Ideas \to \Ideas$ is a \textbf{bijection} (permutation of the 40 Ideas).
\end{theorem}

\begin{proof}
$(\Rightarrow)$ If $\qop{\nu}_k$ is unitary, it preserves inner products: $\qbraket{\qop{\nu}_k\psi}{\qop{\nu}_k\phi} = \qbraket{\psi}{\phi}$. For basis states: $\qbraket{\mathfrak{n}_k(I)}{\mathfrak{n}_k(J)} = \delta_{IJ}$. This requires $\mathfrak{n}_k$ to be injective (hence bijective on finite $\Ideas$).

$(\Leftarrow)$ Proven in Theorem~\ref{thm:noetic-bounded}: bijective $\mathfrak{n}_k$ yields unitary $\qop{\nu}_k$.
\end{proof}

\begin{corollary}[Observable Noetics]
\label{cor:observable-noetics}
A noetic operator $\qop{\nu}_k$ represents a \textbf{quantum observable} (measurable quantity) iff it is Hermitian. In this case:
\begin{enumerate}
  \item All eigenvalues of $\qop{\nu}_k$ are real
  \item Eigenstates form an orthonormal basis
  \item Measurement yields eigenvalues with probabilities given by Born rule
\end{enumerate}
\end{corollary}

\begin{definition}[Normal Noetic Operators]
\label{def:normal-noetics}
A noetic operator $\qop{\nu}_k$ is \textbf{normal} if:
\[
\qop{\nu}_k^\dagger \qop{\nu}_k = \qop{\nu}_k \qop{\nu}_k^\dagger
\]
Both Hermitian and unitary operators are normal. Normal operators have particularly nice spectral properties.
\end{definition}

%% ============================================================================
%% SECTION: SPECTRAL THEORY OF NOETIC OPERATORS
%% ============================================================================
\section{Spectral Theory of Noetic Operators}
\label{sec:spectral-theory-ops}

We now develop the spectral theory for noetic operators, connecting eigenvalues to semantic fixed points and establishing the spectral decomposition that underlies quantum measurement in TKS.

\begin{definition}[Spectrum of Noetic Operator]
\label{def:spectrum-noetic}
The \textbf{spectrum} of $\qop{\nu}_k$ is:
\[
\sigma(\qop{\nu}_k) = \{\lambda \in \mathbb{C} \mid \qop{\nu}_k - \lambda\mathbf{I} \text{ is not invertible}\}
\]
On finite-dimensional $\Hspace_{\Ideas}$, this equals the set of eigenvalues.
\end{definition}

\begin{proposition}[Spectral Bounds]
\label{prop:spectral-bounds}
For any noetic operator $\qop{\nu}_k$:
\begin{enumerate}
  \item $|\sigma(\qop{\nu}_k)| \leq 40$ (at most 40 distinct eigenvalues)
  \item $\sigma(\qop{\nu}_k) \subseteq \{\lambda \in \mathbb{C} \mid |\lambda| \leq \|\qop{\nu}_k\|_{\mathrm{op}}\}$
  \item For permutation noetics (unitary $\qop{\nu}_k$): $\sigma(\qop{\nu}_k) \subseteq \{e^{i\theta} \mid \theta \in [0, 2\pi)\}$
  \item For Hermitian noetics: $\sigma(\qop{\nu}_k) \subseteq \mathbb{R}$
\end{enumerate}
\end{proposition}

\begin{proof}
(1) follows from $\dim(\Hspace_{\Ideas}) = 40$. (2) is the spectral radius formula. (3) follows from unitarity. (4) follows from Hermiticity.
\end{proof}

\begin{theorem}[Fixed Points and Unit Eigenvalue]
\label{thm:fixed-point-eigenvalue}
Let $\mathrm{Fix}(\mathfrak{n}_k) = \{I \in \Ideas \mid \mathfrak{n}_k(I) = I\}$ be the classical fixed point set. Then:
\begin{enumerate}
  \item $1 \in \sigma(\qop{\nu}_k)$ iff $\mathrm{Fix}(\mathfrak{n}_k) \neq \varnothing$
  \item The eigenspace $E_1(\qop{\nu}_k) = \ker(\qop{\nu}_k - \mathbf{I})$ satisfies:
  \[
  E_1(\qop{\nu}_k) \supseteq \mathrm{span}\{\qket{I} \mid I \in \mathrm{Fix}(\mathfrak{n}_k)\}
  \]
  \item The multiplicity of eigenvalue 1 is at least $|\mathrm{Fix}(\mathfrak{n}_k)|$
\end{enumerate}
\end{theorem}

\begin{proof}
For any fixed point $I \in \mathrm{Fix}(\mathfrak{n}_k)$:
\[
\qop{\nu}_k \qket{I} = \qket{\mathfrak{n}_k(I)} = \qket{I}
\]
so $\qket{I}$ is an eigenvector with eigenvalue 1. This establishes (1)--(3).
\end{proof}

\begin{definition}[Cycle Structure and Spectrum]
\label{def:cycle-spectrum}
For a permutation noetic $\mathfrak{n}_k$ (bijection), decompose $\Ideas$ into disjoint cycles:
\[
\Ideas = C_1 \sqcup C_2 \sqcup \cdots \sqcup C_m
\]
where $C_j$ is a cycle of length $\ell_j$. The cycle structure determines the spectrum.
\end{definition}

\begin{theorem}[Spectrum of Permutation Noetics]
\label{thm:spectrum-permutation}
For a permutation noetic $\mathfrak{n}_k$ with cycle structure $(\ell_1, \ldots, \ell_m)$:
\[
\sigma(\qop{\nu}_k) = \bigcup_{j=1}^{m} \{e^{2\pi i r / \ell_j} \mid r = 0, 1, \ldots, \ell_j - 1\}
\]
The multiplicity of $e^{2\pi i r / \ell}$ equals the number of cycles of length divisible by $\ell / \gcd(\ell, r)$.
\end{theorem}

\begin{proof}
Each cycle $C_j$ of length $\ell_j$ contributes a cyclic permutation matrix to $\qop{\nu}_k$ (in the appropriate block). The eigenvalues of an $\ell$-cycle permutation matrix are the $\ell$-th roots of unity $e^{2\pi i r / \ell}$ for $r = 0, \ldots, \ell-1$.
\end{proof}

\begin{corollary}[Fixed Points from Cycle Length]
\label{cor:fixed-cycles}
The number of fixed points $|\mathrm{Fix}(\mathfrak{n}_k)|$ equals the number of 1-cycles (cycles of length 1) in the cycle decomposition, which equals the multiplicity of eigenvalue 1.
\end{corollary}

\begin{definition}[Spectral Projectors]
\label{def:spectral-projectors}
For eigenvalue $\lambda \in \sigma(\qop{\nu}_k)$, the \textbf{spectral projector} $\Projector_\lambda^{(k)}$ projects onto the $\lambda$-eigenspace:
\[
\Projector_\lambda^{(k)} : \Hspace_{\Ideas} \to E_\lambda(\qop{\nu}_k) = \ker(\qop{\nu}_k - \lambda\mathbf{I})
\]
These satisfy:
\begin{enumerate}
  \item $(\Projector_\lambda^{(k)})^2 = \Projector_\lambda^{(k)}$ (idempotent)
  \item $\Projector_\lambda^{(k)} \Projector_{\mu}^{(k)} = 0$ for $\lambda \neq \mu$ (orthogonal for normal $\qop{\nu}_k$)
  \item $\sum_{\lambda \in \sigma(\qop{\nu}_k)} \Projector_\lambda^{(k)} = \mathbf{I}$ (completeness)
\end{enumerate}
\end{definition}

\begin{theorem}[Spectral Decomposition of Noetic Operators]
\label{thm:spectral-decomposition}
For any normal noetic operator $\qop{\nu}_k$ (including Hermitian and unitary noetics):
\[
\qop{\nu}_k = \sum_{\lambda \in \sigma(\qop{\nu}_k)} \lambda \, \Projector_\lambda^{(k)}
\]
This is the \textbf{spectral decomposition} of $\qop{\nu}_k$.
\end{theorem}

\begin{proof}
This is the spectral theorem for normal operators on finite-dimensional Hilbert space. The spectral projectors are:
\[
\Projector_\lambda^{(k)} = \prod_{\mu \neq \lambda} \frac{\qop{\nu}_k - \mu\mathbf{I}}{\lambda - \mu}
\]
\end{proof}

\begin{theorem}[Spectral-Dynamic Correspondence]
\label{thm:spectral-dynamic}
\textbf{(Main Theorem: Connecting Spectral Properties to Classical TKS Dynamics)}

Let $\qop{\nu}_k$ be a noetic operator and consider the classical iterative dynamics $x \mapsto \mathfrak{n}_k(x)$ on $\Ideas$. Then:
\begin{enumerate}
  \item \textbf{Fixed Points}: $I$ is a classical fixed point ($\mathfrak{n}_k(I) = I$) iff $\qket{I}$ is an eigenstate with eigenvalue 1.

  \item \textbf{Periodic Orbits}: $I$ lies on a $p$-cycle ($\mathfrak{n}_k^p(I) = I$, $p$ minimal) iff $\qket{I}$ contributes to eigenstates with eigenvalues $e^{2\pi i r / p}$, $r = 0, \ldots, p-1$.

  \item \textbf{Eventual Fixed Points}: If $\mathfrak{n}_k^m(I) \in \mathrm{Fix}(\mathfrak{n}_k)$ for some $m > 0$, then repeated application of $\qop{\nu}_k$ to $\qket{I}$ converges:
  \[
  \lim_{n \to \infty} \qop{\nu}_k^n \qket{I} = \qket{\mathfrak{n}_k^\omega(I)}
  \]
  where $\mathfrak{n}_k^\omega(I)$ is the transfinite fixed point from v7.2 Theorem~\ref{thm:transfinite-fixed}.

  \item \textbf{Attractor Correspondence}: The quantum fixed-point projector:
  \[
  \Projector_{\mathrm{fix}}^{(k)} = \lim_{n \to \infty} \frac{1}{n} \sum_{j=0}^{n-1} \qop{\nu}_k^j
  \]
  projects onto the span of classical attractors of $\mathfrak{n}_k$.
\end{enumerate}
\end{theorem}

\begin{proof}
\textbf{(1)} Proven in Theorem~\ref{thm:fixed-point-eigenvalue}.

\textbf{(2)} If $I$ lies on a $p$-cycle $\{I, \mathfrak{n}_k(I), \ldots, \mathfrak{n}_k^{p-1}(I)\}$, consider the cyclic superpositions:
\[
\qket{\phi_r} = \frac{1}{\sqrt{p}} \sum_{j=0}^{p-1} e^{-2\pi i r j / p} \qket{\mathfrak{n}_k^j(I)}
\]
These satisfy $\qop{\nu}_k \qket{\phi_r} = e^{2\pi i r / p} \qket{\phi_r}$.

\textbf{(3)} For non-permutation noetics where orbits eventually reach fixed points, write $\qop{\nu}_k = \sum_\lambda \lambda \Projector_\lambda$. For $|\lambda| < 1$, the term $\lambda^n \Projector_\lambda \to 0$. For $|\lambda| = 1$ but $\lambda \neq 1$, these contribute oscillatory but bounded terms. Only the $\lambda = 1$ component survives in the limit, yielding the fixed-point subspace.

\textbf{(4)} The Cesaro mean $\frac{1}{n} \sum_{j=0}^{n-1} \qop{\nu}_k^j$ averages out oscillatory components (eigenvalues $\lambda \neq 1$ with $|\lambda| = 1$ satisfy $\frac{1}{n}\sum_{j=0}^{n-1} \lambda^j \to 0$), leaving only the fixed-point projector.
\end{proof}

\begin{corollary}[Idempotent Noetics are Instant Projectors]
\label{cor:idem-instant}
For idempotent noetics ($\mathfrak{n}_k^2 = \mathfrak{n}_k$), the spectral structure simplifies:
\[
\sigma(\qop{\nu}_k) \subseteq \{0, 1\}
\]
and $\qop{\nu}_k$ is already a projector: $\qop{\nu}_k = \Projector_1^{(k)}$.
\end{corollary}

\begin{proof}
Idempotence implies $\qop{\nu}_k^2 = \qop{\nu}_k$, so $\qop{\nu}_k(\qop{\nu}_k - \mathbf{I}) = 0$. Every eigenvalue satisfies $\lambda(\lambda - 1) = 0$, hence $\lambda \in \{0, 1\}$.
\end{proof}

%% ----------------------------------------------------------------------------
%% HANDOFF NOTE
%% ----------------------------------------------------------------------------
\begin{tcolorbox}[colback=green!5!white,colframe=green!50!black,title=\textbf{Handoff: Next Steps},breakable]
This completes the C*-algebraic structure and core spectral theory for noetic operators in Chapter 1. The subsequent chapters should:
\begin{itemize}
  \item \textbf{Chapter 2 (Spectral Theory of Noetics)}: Develop detailed spectrum computations for each specific noetic $\qop{\nu}_1, \ldots, \qop{\nu}_9$; establish joint spectral theory for commuting families; analyze spectral gaps and their semantic interpretation
  \item \textbf{Chapter 3 (Density Matrix Formalism)}: Extend to mixed states $\rho$; develop von Neumann entropy; connect to statistical mixtures of noetic transformations
  \item \textbf{Chapter 5 (Entanglement)}: Build on tensor products $\Hspace_{\Ideas} \otimes \Hspace_{\Ideas}$ for multi-Idea systems (deferred from this chunk as instructed)
\end{itemize}
Key results established:
\begin{enumerate}
  \item Noetic C*-algebra $\mathcal{A}_\nu = C^*(\qop{\nu}_1, \ldots, \qop{\nu}_9)$ (Def.~\ref{def:noetic-cstar})
  \item Commutator analysis: $[\qop{\nu}_i, \qop{\nu}_j] = 0$ iff classical compositions commute (Prop.~\ref{prop:commutator-elements})
  \item Spectral-Dynamic Correspondence Theorem~\ref{thm:spectral-dynamic}: eigenvalues encode fixed points and cycle structure
\end{enumerate}
\end{tcolorbox}

%% ============================================================================
%% CHAPTER 2: SPECTRAL THEORY OF NOETICS
%% ============================================================================
\chapter{Spectral Theory of Noetics}
\label{ch:spectral-theory}

\begin{tcolorbox}[colback=blue!5!white,colframe=blue!75!black,title=\vnewthree]
This chapter develops complete spectral theory for noetic operators, extending v7.2 spectral decomposition with detailed eigenvalue analysis.
\end{tcolorbox}

%% ----------------------------------------------------------------------------
\section{Spectrum of Each Noetic}
\label{sec:spectrum-each-noetic}
%% Content: Compute sigma(hat{nu}_k) for k = 0, ..., 9
%% Agent: tks-quantum
%% Handoff: Point spectrum, continuous spectrum, residual spectrum

%% ----------------------------------------------------------------------------
\section{Spectral Decomposition}
\label{sec:spectral-decomposition}
%% Content: hat{nu}_k = integral lambda dE_lambda
%% Agent: tks-quantum
%% Handoff: Projection-valued measure, spectral theorem application

%% ----------------------------------------------------------------------------
\section{Functional Calculus}
\label{sec:functional-calculus}
%% Content: f(hat{nu}_k) for measurable f
%% Agent: tks-quantum
%% Handoff: Polynomial, continuous, Borel functional calculus

%% ----------------------------------------------------------------------------
\section{Joint Spectra}
\label{sec:joint-spectra}
%% Content: Simultaneous diagonalization of commuting noetics
%% Agent: tks-quantum
%% Handoff: Joint spectral measure, complete sets of commuting operators

%% ----------------------------------------------------------------------------
\section{Spectral Gaps}
\label{sec:spectral-gaps}
%% Content: Gaps in noetic spectra and their interpretation
%% Agent: tks-quantum
%% Handoff: Gap = stable region, no transition

%% ============================================================================
%% CHAPTER 3: DENSITY MATRIX FORMALISM
%% ============================================================================
\chapter{Density Matrix Formalism}
\label{ch:density-matrix}

\begin{tcolorbox}[colback=blue!5!white,colframe=blue!75!black,title=\vnewthree]
This chapter develops the density matrix formalism for mixed states in TKS, enabling description of statistical ensembles and open systems.
\end{tcolorbox}

%% ----------------------------------------------------------------------------
\section{Pure vs Mixed Noetic States}
\label{sec:pure-vs-mixed}

This section introduces the density operator formalism, which generalizes pure quantum states to include statistical mixtures arising from incomplete knowledge of the noetic state.

\begin{definition}[Density Operator]
\label{def:density-operator}
A \textbf{density operator} (or \textbf{density matrix}) on $\Hspace_{\Ideas}$ is a linear operator $\rho : \Hspace_{\Ideas} \to \Hspace_{\Ideas}$ satisfying:
\begin{enumerate}
  \item \textbf{Hermiticity}: $\rho^\dagger = \rho$
  \item \textbf{Positive semi-definiteness}: $\qbra{\psi}\rho\qket{\psi} \geq 0$ for all $\qket{\psi} \in \Hspace_{\Ideas}$
  \item \textbf{Trace normalization}: $\Trace(\rho) = 1$
\end{enumerate}
The set of all density operators on $\Hspace_{\Ideas}$ is denoted $\mathcal{D}(\Hspace_{\Ideas})$.
\end{definition}

\begin{proposition}[Convexity of Density Operators]
\label{prop:density-convex}
The set $\mathcal{D}(\Hspace_{\Ideas})$ is convex: if $\rho_1, \rho_2 \in \mathcal{D}(\Hspace_{\Ideas})$ and $0 \leq p \leq 1$, then:
\[
\rho = p \rho_1 + (1-p) \rho_2 \in \mathcal{D}(\Hspace_{\Ideas})
\]
\end{proposition}

\begin{proof}
Hermiticity: $(p\rho_1 + (1-p)\rho_2)^\dagger = p\rho_1^\dagger + (1-p)\rho_2^\dagger = p\rho_1 + (1-p)\rho_2$.

Positive semi-definiteness: $\qbra{\psi}\rho\qket{\psi} = p\qbra{\psi}\rho_1\qket{\psi} + (1-p)\qbra{\psi}\rho_2\qket{\psi} \geq 0$.

Trace: $\Trace(\rho) = p\Trace(\rho_1) + (1-p)\Trace(\rho_2) = p + (1-p) = 1$.
\end{proof}

\begin{definition}[Pure State Density Operator]
\label{def:pure-state-density}
A \textbf{pure state} $\qket{\psi} \in \Hspace_{\Ideas}$ (with $\|\psi\| = 1$) corresponds to the \textbf{rank-1 projector}:
\[
\rho_\psi = \qket{\psi}\qbra{\psi}
\]
This satisfies:
\begin{enumerate}
  \item $\rho_\psi^2 = \rho_\psi$ (idempotent)
  \item $\mathrm{rank}(\rho_\psi) = 1$
  \item $\Trace(\rho_\psi^2) = 1$ (maximal purity)
\end{enumerate}
\end{definition}

\begin{definition}[Mixed State]
\label{def:mixed-state}
A \textbf{mixed state} is a density operator that is not pure, i.e., cannot be written as $\qket{\psi}\qbra{\psi}$ for any single $\qket{\psi}$. Equivalently:
\begin{enumerate}
  \item $\rho^2 \neq \rho$
  \item $\mathrm{rank}(\rho) > 1$
  \item $\Trace(\rho^2) < 1$
\end{enumerate}
\end{definition}

\begin{definition}[Statistical Mixture]
\label{def:statistical-mixture}
A \textbf{statistical mixture} (or \textbf{ensemble}) is a density operator of the form:
\[
\rho = \sum_{i=1}^{n} p_i \qket{\psi_i}\qbra{\psi_i}
\]
where:
\begin{itemize}
  \item $\{p_i\}$ are classical probabilities: $p_i \geq 0$ and $\sum_i p_i = 1$
  \item $\{\qket{\psi_i}\}$ are normalized (but not necessarily orthogonal) states
\end{itemize}
This represents classical uncertainty about which pure state $\qket{\psi_i}$ the system is in.
\end{definition}

\begin{theorem}[Spectral Decomposition of Density Operators]
\label{thm:density-spectral}
Every density operator $\rho \in \mathcal{D}(\Hspace_{\Ideas})$ has a unique spectral decomposition:
\[
\rho = \sum_{j=1}^{r} \lambda_j \qket{\phi_j}\qbra{\phi_j}
\]
where:
\begin{itemize}
  \item $\{\qket{\phi_j}\}$ are orthonormal eigenvectors
  \item $\lambda_j \geq 0$ are eigenvalues with $\sum_j \lambda_j = 1$
  \item $r = \mathrm{rank}(\rho) \leq 40$
\end{itemize}
\end{theorem}

\begin{proof}
By the spectral theorem for Hermitian operators. Positive semi-definiteness ensures $\lambda_j \geq 0$; trace normalization ensures $\sum_j \lambda_j = 1$.
\end{proof}

\begin{remark}[Superposition vs Mixture: Physical Interpretation]
\label{rem:superposition-vs-mixture}
The distinction between pure and mixed states is fundamental:

\textbf{Superposition} (pure state): $\qket{\psi} = \alpha\qket{A1} + \beta\qket{B1}$ with $|\alpha|^2 + |\beta|^2 = 1$
\begin{itemize}
  \item The system genuinely ``is'' in both states simultaneously
  \item Interference effects are present: $|\alpha + \beta|^2 \neq |\alpha|^2 + |\beta|^2$ in general
  \item Complete knowledge of the quantum state
  \item Density matrix: $\rho = \qket{\psi}\qbra{\psi}$ has off-diagonal terms $\alpha\bar{\beta}\qket{A1}\qbra{B1}$
\end{itemize}

\textbf{Mixture} (mixed state): $\rho = p\qket{A1}\qbra{A1} + (1-p)\qket{B1}\qbra{B1}$
\begin{itemize}
  \item The system is in state $\qket{A1}$ with probability $p$ or $\qket{B1}$ with probability $1-p$
  \item No interference: classical probabilistic uncertainty
  \item Incomplete knowledge about which pure state applies
  \item Density matrix is diagonal in the $\{\qket{A1}, \qket{B1}\}$ basis
\end{itemize}

In TKS semantics: a superposition represents genuine quantum indefiniteness of an Idea, while a mixture represents classical ignorance about which Idea is present.
\end{remark}

\begin{definition}[Idea Basis Density Matrix]
\label{def:idea-basis-density}
In the Idea basis $\{\qket{Wn}\}$, any density operator has the matrix representation:
\[
\rho = \sum_{W,n} \sum_{W',n'} \rho_{Wn,W'n'} \qket{Wn}\qbra{W'n'}
\]
where:
\begin{itemize}
  \item Diagonal elements $\rho_{Wn,Wn} = \qbra{Wn}\rho\qket{Wn} \geq 0$ are probabilities of finding Idea $Wn$
  \item Off-diagonal elements $\rho_{Wn,W'n'}$ ($Wn \neq W'n'$) are \textbf{coherences} encoding quantum superposition
\end{itemize}
\end{definition}

\begin{definition}[Classical Noetic State]
\label{def:classical-noetic-state}
A \textbf{classical noetic state} is a density operator diagonal in the Idea basis:
\[
\rho_{\mathrm{cl}} = \sum_{W,n} p_{Wn} \qket{Wn}\qbra{Wn}
\]
where $\{p_{Wn}\}$ is a classical probability distribution over $\Ideas$. This corresponds to the v7.2 probability measures on $\Ideas$ (Definition~\ref{def:uniform-measure}).
\end{definition}

\begin{proposition}[Uniform Noetic State]
\label{prop:uniform-state}
The \textbf{maximally mixed state} on $\Hspace_{\Ideas}$:
\[
\rho_{\mathrm{mix}} = \frac{1}{40} \mathbf{I} = \frac{1}{40} \sum_{W,n} \qket{Wn}\qbra{Wn}
\]
represents complete ignorance about the Idea. This is the quantum analogue of the uniform measure $\mu_U$ from v7.2.
\end{proposition}

%% ----------------------------------------------------------------------------
\section{Density Matrix Properties}
\label{sec:density-matrix-properties}

\begin{definition}[Purity]
\label{def:purity}
The \textbf{purity} of a density operator $\rho$ is:
\[
\gamma(\rho) = \Trace(\rho^2)
\]
Properties:
\begin{enumerate}
  \item $\frac{1}{40} \leq \gamma(\rho) \leq 1$
  \item $\gamma(\rho) = 1$ iff $\rho$ is pure
  \item $\gamma(\rho) = \frac{1}{40}$ iff $\rho = \rho_{\mathrm{mix}}$ (maximally mixed)
\end{enumerate}
\end{definition}

\begin{proof}
By spectral decomposition $\rho = \sum_j \lambda_j \qket{\phi_j}\qbra{\phi_j}$:
\[
\gamma(\rho) = \Trace(\rho^2) = \sum_j \lambda_j^2
\]
Since $\sum_j \lambda_j = 1$ and $\lambda_j \geq 0$, the sum of squares is maximized when one $\lambda_j = 1$ (pure state, $\gamma = 1$) and minimized when all $\lambda_j = \frac{1}{40}$ (maximally mixed, $\gamma = \frac{1}{40}$).
\end{proof}

\begin{definition}[Linear Entropy]
\label{def:linear-entropy}
The \textbf{linear entropy} provides a computationally simpler measure of mixedness:
\[
S_L(\rho) = 1 - \Trace(\rho^2) = 1 - \gamma(\rho)
\]
This satisfies $0 \leq S_L(\rho) \leq \frac{39}{40}$, with $S_L = 0$ for pure states.
\end{definition}

\begin{definition}[Von Neumann Entropy]
\label{def:von-neumann-entropy}
The \textbf{von Neumann entropy} of a density operator $\rho$ is:
\[
S(\rho) = -\Trace(\rho \log \rho)
\]
where $\log$ denotes the natural logarithm and we use the convention $0 \log 0 = 0$. By spectral decomposition:
\[
S(\rho) = -\sum_j \lambda_j \log \lambda_j
\]
\end{definition}

\begin{proposition}[Von Neumann Entropy Properties]
\label{prop:entropy-properties}
The von Neumann entropy satisfies:
\begin{enumerate}
  \item \textbf{Non-negativity}: $S(\rho) \geq 0$
  \item \textbf{Purity criterion}: $S(\rho) = 0$ iff $\rho$ is pure
  \item \textbf{Maximum entropy}: $S(\rho) \leq \log(40)$, with equality iff $\rho = \rho_{\mathrm{mix}}$
  \item \textbf{Concavity}: $S(p\rho_1 + (1-p)\rho_2) \geq p S(\rho_1) + (1-p) S(\rho_2)$
  \item \textbf{Unitary invariance}: $S(U\rho U^\dagger) = S(\rho)$ for unitary $U$
\end{enumerate}
\end{proposition}

\begin{proof}
These are standard properties of von Neumann entropy:
\begin{enumerate}
  \item Since $0 \leq \lambda_j \leq 1$, we have $\lambda_j \log \lambda_j \leq 0$, so $S \geq 0$.
  \item $S = 0$ iff all but one $\lambda_j = 0$, i.e., rank-1 (pure).
  \item Maximum of $-\sum_j \lambda_j \log \lambda_j$ subject to $\sum_j \lambda_j = 1$ is achieved when all $\lambda_j = \frac{1}{40}$, giving $S = \log(40)$.
  \item Standard concavity proof using operator concavity of $-x\log x$.
  \item Unitary conjugation preserves eigenvalues.
\end{enumerate}
\end{proof}

\begin{theorem}[Entropy-Information Correspondence]
\label{thm:entropy-info-correspondence}
\textbf{(Main Theorem: Connecting Quantum Entropy to Classical TKS Information)}

Let $\rho$ be a density operator on $\Hspace_{\Ideas}$ and let $\{p_{Wn}\} = \{\qbra{Wn}\rho\qket{Wn}\}$ be the diagonal (classical) probability distribution over Ideas. Define the \textbf{classical Shannon entropy}:
\[
H_{\mathrm{cl}}(\rho) = -\sum_{W,n} p_{Wn} \log p_{Wn}
\]
Then:
\begin{enumerate}
  \item \textbf{Entropy inequality}: $S(\rho) \leq H_{\mathrm{cl}}(\rho)$, with equality iff $\rho$ is diagonal in the Idea basis (classical noetic state).

  \item \textbf{Coherence contribution}: The difference $H_{\mathrm{cl}}(\rho) - S(\rho) \geq 0$ quantifies the ``quantum coherence'' of $\rho$---information stored in off-diagonal elements.

  \item \textbf{Measurement interpretation}: $H_{\mathrm{cl}}(\rho)$ is the entropy of outcomes when measuring $\rho$ in the Idea basis. The inequality says measurement cannot decrease uncertainty.

  \item \textbf{Connection to v7.2 measure theory}: For a classical noetic state $\rho_{\mathrm{cl}}$ corresponding to probability measure $\mu$ on $\Ideas$:
  \[
  S(\rho_{\mathrm{cl}}) = H_{\mathrm{cl}}(\rho_{\mathrm{cl}}) = -\int_{\Ideas} \log\left(\frac{d\mu}{d\mu_U}\right) d\mu
  \]
  the standard entropy of $\mu$ relative to the counting measure.
\end{enumerate}
\end{theorem}

\begin{proof}
\textbf{(1)} Let $\rho = \sum_j \lambda_j \qket{\phi_j}\qbra{\phi_j}$ be the spectral decomposition and let $U$ be the unitary mapping $\qket{\phi_j} \mapsto \qket{e_j}$ for some ordering of basis states. Define $\sigma = U\rho U^\dagger$, which has the same eigenvalues as $\rho$.

The diagonal elements of $\rho$ in the Idea basis are $p_{Wn} = \sum_j \lambda_j |\langle Wn | \phi_j \rangle|^2$. By Schur-concavity of entropy and the doubly stochastic nature of the matrix $|\langle Wn | \phi_j \rangle|^2$:
\[
H_{\mathrm{cl}}(\rho) = H(\{p_{Wn}\}) \geq H(\{\lambda_j\}) = S(\rho)
\]
Equality holds iff $|\langle Wn | \phi_j \rangle|^2 \in \{0, 1\}$, i.e., each $\qket{\phi_j}$ is an Idea basis state.

\textbf{(2)} Follows from (1).

\textbf{(3)} Measurement in basis $\{\qket{Wn}\}$ produces outcome $Wn$ with probability $p_{Wn}$. The post-measurement state is $\qket{Wn}\qbra{Wn}$ with probability $p_{Wn}$, giving mixture $\sum_{W,n} p_{Wn} \qket{Wn}\qbra{Wn}$ with entropy $H_{\mathrm{cl}}(\rho)$.

\textbf{(4)} For classical $\rho_{\mathrm{cl}} = \sum_{W,n} p_{Wn} \qket{Wn}\qbra{Wn}$, the spectral decomposition is already in the Idea basis with eigenvalues $p_{Wn}$, so $S(\rho_{\mathrm{cl}}) = -\sum_{W,n} p_{Wn} \log p_{Wn} = H_{\mathrm{cl}}(\rho_{\mathrm{cl}})$. The integral form is the standard entropy of a discrete probability measure.
\end{proof}

%% ----------------------------------------------------------------------------
\section{Time Evolution of Density Matrices}
\label{sec:density-evolution}

We now develop the quantum dynamics of density matrices under noetic evolution.

\begin{definition}[Discrete Noetic Evolution]
\label{def:discrete-density-evolution}
For a noetic operator $\qop{\nu}_k$, the \textbf{discrete evolution} of a density matrix is:
\[
\rho \mapsto \qop{\nu}_k \rho \qop{\nu}_k^\dagger
\]
For unitary noetics (permutation $\mathfrak{n}_k$), this is the standard unitary evolution preserving trace and positivity.
\end{definition}

\begin{proposition}[Trace Preservation under Unitary Evolution]
\label{prop:trace-preserve-unitary}
If $\qop{\nu}_k$ is unitary (i.e., $\mathfrak{n}_k$ is a bijection), then:
\[
\Trace(\qop{\nu}_k \rho \qop{\nu}_k^\dagger) = \Trace(\rho)
\]
and the evolution preserves the set of density operators: $\qop{\nu}_k \rho \qop{\nu}_k^\dagger \in \mathcal{D}(\Hspace_{\Ideas})$.
\end{proposition}

\begin{proof}
$\Trace(\qop{\nu}_k \rho \qop{\nu}_k^\dagger) = \Trace(\qop{\nu}_k^\dagger \qop{\nu}_k \rho) = \Trace(\rho)$ by cyclicity of trace and unitarity.
\end{proof}

\begin{definition}[Noetic Hamiltonian]
\label{def:noetic-hamiltonian}
The \textbf{noetic Hamiltonian} associated with noetic $k$ is defined (for Hermitian $\qop{\nu}_k$) as:
\[
\qop{H}_k = \hbar \omega_k \qop{\nu}_k
\]
where $\omega_k$ is a characteristic frequency. For non-Hermitian noetics, we use the Hermitian combination:
\[
\qop{H}_k = \frac{\hbar \omega_k}{2}(\qop{\nu}_k + \qop{\nu}_k^\dagger)
\]
\end{definition}

\begin{definition}[Von Neumann Equation]
\label{def:von-neumann-equation}
For a noetic Hamiltonian $\qop{H}_k$, the \textbf{von Neumann equation} governs continuous-time evolution:
\[
\frac{d\rho}{dt} = -\frac{i}{\hbar}[\qop{H}_k, \rho] = -i\omega_k[\qop{\nu}_k, \rho]
\]
where $[\qop{A}, \qop{B}] = \qop{A}\qop{B} - \qop{B}\qop{A}$ is the commutator.
\end{definition}

\begin{theorem}[Solution of Von Neumann Equation]
\label{thm:von-neumann-solution}
The solution to the von Neumann equation with initial condition $\rho(0) = \rho_0$ is:
\[
\rho(t) = \qop{U}_k(t) \rho_0 \qop{U}_k(t)^\dagger
\]
where $\qop{U}_k(t) = e^{-i\omega_k t \qop{\nu}_k}$ is the time evolution operator (assuming Hermitian $\qop{\nu}_k$).
\end{theorem}

\begin{proof}
Differentiating:
\begin{align}
\frac{d\rho}{dt} &= \frac{d}{dt}\left(\qop{U}_k \rho_0 \qop{U}_k^\dagger\right) \\
&= \left(\frac{d\qop{U}_k}{dt}\right) \rho_0 \qop{U}_k^\dagger + \qop{U}_k \rho_0 \left(\frac{d\qop{U}_k^\dagger}{dt}\right) \\
&= (-i\omega_k \qop{\nu}_k \qop{U}_k) \rho_0 \qop{U}_k^\dagger + \qop{U}_k \rho_0 (i\omega_k \qop{U}_k^\dagger \qop{\nu}_k) \\
&= -i\omega_k[\qop{\nu}_k, \qop{U}_k \rho_0 \qop{U}_k^\dagger] = -i\omega_k[\qop{\nu}_k, \rho]
\end{align}
\end{proof}

\begin{definition}[Liouvillian Superoperator]
\label{def:liouvillian}
The \textbf{Liouvillian} (or \textbf{Liouville superoperator}) for noetic $k$ is:
\[
\mathcal{L}_k(\rho) = -i\omega_k[\qop{\nu}_k, \rho]
\]
This is a linear map on operators (a ``superoperator''). The von Neumann equation becomes:
\[
\frac{d\rho}{dt} = \mathcal{L}_k(\rho)
\]
\end{definition}

\begin{proposition}[Relationship to Classical Iteration]
\label{prop:classical-iteration-relation}
For a classical noetic state $\rho_{\mathrm{cl}} = \sum_{W,n} p_{Wn} \qket{Wn}\qbra{Wn}$ and a permutation noetic $\qop{\nu}_k$:
\[
\qop{\nu}_k \rho_{\mathrm{cl}} \qop{\nu}_k^\dagger = \sum_{W,n} p_{Wn} \qket{\mathfrak{n}_k(Wn)}\qbra{\mathfrak{n}_k(Wn)}
\]
This is the quantum version of the pushforward: if $\mu$ is the probability measure on $\Ideas$ corresponding to $\rho_{\mathrm{cl}}$, then the evolved state corresponds to $(\mathfrak{n}_k)_* \mu$.
\end{proposition}

\begin{proof}
Direct calculation:
\[
\qop{\nu}_k \qket{Wn}\qbra{Wn} \qop{\nu}_k^\dagger = \qket{\mathfrak{n}_k(Wn)}\qbra{\mathfrak{n}_k(Wn)}
\]
so the coefficients are permuted according to $\mathfrak{n}_k^{-1}$.
\end{proof}

\begin{definition}[Steady-State Density Matrix]
\label{def:steady-state}
A \textbf{steady state} (or \textbf{stationary state}) for noetic $k$ is a density operator $\rho_{\mathrm{ss}}$ satisfying:
\[
[\qop{\nu}_k, \rho_{\mathrm{ss}}] = 0
\]
Equivalently, $\mathcal{L}_k(\rho_{\mathrm{ss}}) = 0$, so $\rho_{\mathrm{ss}}$ is time-independent under von Neumann evolution.
\end{definition}

\begin{theorem}[Characterization of Steady States]
\label{thm:steady-state-char}
A density operator $\rho$ is a steady state for noetic $k$ iff it is block-diagonal in the eigenspaces of $\qop{\nu}_k$:
\[
\rho_{\mathrm{ss}} = \sum_{\lambda \in \sigma(\qop{\nu}_k)} \rho_\lambda
\]
where $\rho_\lambda$ is supported on the $\lambda$-eigenspace $E_\lambda(\qop{\nu}_k)$.
\end{theorem}

\begin{proof}
$[\qop{\nu}_k, \rho] = 0$ iff $\rho$ commutes with $\qop{\nu}_k$. For normal $\qop{\nu}_k$, this happens iff $\rho$ is diagonal in the same eigenbasis, i.e., block-diagonal in eigenspaces.
\end{proof}

\begin{corollary}[Fixed-Point Steady States]
\label{cor:fixed-point-steady}
The projector onto the fixed-point eigenspace:
\[
\rho_{\mathrm{fix}}^{(k)} = \frac{1}{|\mathrm{Fix}(\mathfrak{n}_k)|} \sum_{I \in \mathrm{Fix}(\mathfrak{n}_k)} \qket{I}\qbra{I}
\]
is a steady state for noetic $k$. This connects to the transfinite fixed-point convergence of v7.2 (Theorem~\ref{thm:transfinite-fixed}).
\end{corollary}

%% ----------------------------------------------------------------------------
\section{Noetic Mixtures}
\label{sec:noetic-mixtures}

Mixed states arise naturally in TKS when there is uncertainty about which noetic transformation has been applied.

\begin{definition}[Noetic Mixture]
\label{def:noetic-mixture}
A \textbf{noetic mixture} is a density operator representing classical uncertainty about which noetic has been applied:
\[
\rho = \sum_{k=0}^{9} p_k \, \qop{\nu}_k \rho_0 \qop{\nu}_k^\dagger
\]
where $\{p_k\}$ is a probability distribution over noetics and $\rho_0$ is the initial state.
\end{definition}

\begin{example}[Uncertain Element Application]
\label{ex:uncertain-element}
If we know an Element was applied but not which one, and Elements correspond to noetics $\{\nu_1, \nu_2, \ldots, \nu_m\}$ with equal probability:
\[
\rho = \frac{1}{m} \sum_{j=1}^{m} \qop{\nu}_j \rho_0 \qop{\nu}_j^\dagger
\]
This ``averages out'' the effect of the unknown Element.
\end{example}

\begin{proposition}[Noetic Mixture as Quantum Channel]
\label{prop:mixture-channel}
The map $\mathcal{E}_{\mathrm{mix}} : \rho \mapsto \sum_k p_k \qop{\nu}_k \rho \qop{\nu}_k^\dagger$ is a \textbf{quantum channel} (completely positive trace-preserving map) when each $\qop{\nu}_k$ is unitary. This is a \textbf{random unitary channel}.
\end{proposition}

\begin{definition}[Noetic Dephasing]
\label{def:noetic-dephasing}
\textbf{Noetic dephasing} occurs when the system interacts with an environment that ``measures'' the Idea without our knowledge. The resulting channel:
\[
\mathcal{E}_{\mathrm{deph}}(\rho) = \sum_{W,n} \qket{Wn}\qbra{Wn} \rho \qket{Wn}\qbra{Wn}
\]
destroys all off-diagonal coherences, leaving only the classical probability distribution:
\[
\mathcal{E}_{\mathrm{deph}}(\rho) = \sum_{W,n} \qbra{Wn}\rho\qket{Wn} \cdot \qket{Wn}\qbra{Wn}
\]
\end{definition}

\begin{proposition}[Dephasing Increases Entropy]
\label{prop:dephasing-entropy}
For any density operator $\rho$:
\[
S(\mathcal{E}_{\mathrm{deph}}(\rho)) \geq S(\rho)
\]
with equality iff $\rho$ is already classical (diagonal in the Idea basis).
\end{proposition}

\begin{proof}
By Theorem~\ref{thm:entropy-info-correspondence}, $S(\mathcal{E}_{\mathrm{deph}}(\rho)) = H_{\mathrm{cl}}(\rho) \geq S(\rho)$.
\end{proof}

%% ----------------------------------------------------------------------------
\section{Partial Trace and Reduced States}
\label{sec:partial-trace}

The partial trace operation allows us to describe subsystems of a larger noetic system, particularly when we focus on individual Worlds within the full Idea space.

\begin{definition}[World Decomposition of Density Operators]
\label{def:world-decomp-density}
Using the world decomposition $\Hspace_{\Ideas} = \Hspace_A \oplus \Hspace_B \oplus \Hspace_C \oplus \Hspace_D$, a density operator can be written in block form:
\[
\rho = \begin{pmatrix}
\rho_{AA} & \rho_{AB} & \rho_{AC} & \rho_{AD} \\
\rho_{BA} & \rho_{BB} & \rho_{BC} & \rho_{BD} \\
\rho_{CA} & \rho_{CB} & \rho_{CC} & \rho_{CD} \\
\rho_{DA} & \rho_{DB} & \rho_{DC} & \rho_{DD}
\end{pmatrix}
\]
where $\rho_{WW'} : \Hspace_{W'} \to \Hspace_W$ are the block operators.
\end{definition}

\begin{definition}[World Projector]
\label{def:world-projector}
The \textbf{projector onto World $W$} is:
\[
\Projector_W = \sum_{n=1}^{10} \qket{Wn}\qbra{Wn}
\]
These satisfy $\Projector_W^2 = \Projector_W$, $\Projector_W^\dagger = \Projector_W$, and $\sum_{W \in \{A,B,C,D\}} \Projector_W = \mathbf{I}$.
\end{definition}

\begin{definition}[World-Restricted Density Operator]
\label{def:world-restricted}
The \textbf{restriction of $\rho$ to World $W$} is:
\[
\rho|_W = \Projector_W \rho \Projector_W
\]
This is a positive semi-definite operator on $\Hspace_W$ with $\Trace(\rho|_W) = p_W$, the probability of finding the system in World $W$.
\end{definition}

\begin{definition}[Normalized World State]
\label{def:normalized-world-state}
If $p_W = \Trace(\rho|_W) > 0$, the \textbf{normalized density operator for World $W$} is:
\[
\rho_W = \frac{\rho|_W}{p_W} = \frac{\Projector_W \rho \Projector_W}{\Trace(\Projector_W \rho)}
\]
This represents the state conditioned on being in World $W$.
\end{definition}

\begin{proposition}[World Probability Distribution]
\label{prop:world-probs}
The probability of finding the system in World $W$ is:
\[
p_W = \Trace(\Projector_W \rho) = \sum_{n=1}^{10} \qbra{Wn}\rho\qket{Wn}
\]
These satisfy $\sum_W p_W = 1$.
\end{proposition}

\begin{definition}[Partial Trace over Worlds]
\label{def:partial-trace-worlds}
For a tensor product structure $\Hspace_{\Ideas} \cong \mathbb{C}^4 \otimes \mathbb{C}^{10}$ (World $\otimes$ Element-within-World), define:

\textbf{Trace over Elements} (keeping World information):
\[
\Trace_{\mathrm{Elem}}(\rho) = \sum_{n=1}^{10} (\mathbf{I}_W \otimes \qbra{n}) \rho (\mathbf{I}_W \otimes \qket{n})
\]
This is a $4 \times 4$ matrix on the World space.

\textbf{Trace over Worlds} (keeping Element information):
\[
\Trace_{\mathrm{World}}(\rho) = \sum_{W \in \{A,B,C,D\}} (\qbra{W} \otimes \mathbf{I}_{10}) \rho (\qket{W} \otimes \mathbf{I}_{10})
\]
This is a $10 \times 10$ matrix on the Element space.
\end{definition}

\begin{proposition}[Reduced World State]
\label{prop:reduced-world-state}
The \textbf{reduced density matrix on Worlds} is:
\[
\rho_{\mathrm{World}} = \Trace_{\mathrm{Elem}}(\rho) = \sum_{W,W'} \left(\sum_{n=1}^{10} \rho_{WnW'n}\right) \qket{W}\qbra{W'}
\]
where $\rho_{WnW'n'} = \qbra{Wn}\rho\qket{W'n'}$. The diagonal elements are $(\rho_{\mathrm{World}})_{WW} = p_W$.
\end{proposition}

\begin{theorem}[Information Loss under Partial Trace]
\label{thm:info-loss-trace}
For any density operator $\rho$ on $\Hspace_{\Ideas}$:
\[
S(\rho) \leq S(\rho_{\mathrm{World}}) + S(\rho_{\mathrm{Elem}})
\]
Equality holds iff $\rho$ is a product state $\rho = \rho_{\mathrm{World}} \otimes \rho_{\mathrm{Elem}}$ (no World-Element correlations).
\end{theorem}

\begin{proof}
This is the subadditivity of von Neumann entropy for tensor product decompositions. The difference $S(\rho_{\mathrm{World}}) + S(\rho_{\mathrm{Elem}}) - S(\rho)$ is the quantum mutual information, which is non-negative and zero iff the state factorizes.
\end{proof}

\begin{definition}[World Entropy]
\label{def:world-entropy}
The \textbf{World entropy} is the von Neumann entropy of the reduced World state:
\[
S_W(\rho) = S(\rho_{\mathrm{World}}) = -\Trace(\rho_{\mathrm{World}} \log \rho_{\mathrm{World}})
\]
This measures uncertainty about which World the Idea belongs to. Maximum value is $\log(4) \approx 1.39$ (uniform over Worlds).
\end{definition}

\begin{definition}[Conditional Element Entropy]
\label{def:conditional-entropy}
The \textbf{average Element entropy given World} is:
\[
S_{E|W}(\rho) = \sum_W p_W S(\rho_W^{\mathrm{Elem}})
\]
where $\rho_W^{\mathrm{Elem}}$ is the normalized Element distribution within World $W$. This measures uncertainty about the Element within each World.
\end{definition}

\begin{theorem}[Entropy Decomposition]
\label{thm:entropy-decomp}
For a classical noetic state $\rho_{\mathrm{cl}}$ (diagonal in Idea basis):
\[
S(\rho_{\mathrm{cl}}) = S_W(\rho_{\mathrm{cl}}) + S_{E|W}(\rho_{\mathrm{cl}})
\]
This is the chain rule for Shannon entropy: total uncertainty = World uncertainty + average Element uncertainty given World.
\end{theorem}

\begin{proof}
For classical states, von Neumann entropy equals Shannon entropy, and the chain rule $H(X,Y) = H(X) + H(Y|X)$ applies directly with $X = $ World and $Y = $ Element.
\end{proof}

%% ----------------------------------------------------------------------------
%% HANDOFF NOTE
%% ----------------------------------------------------------------------------
\begin{tcolorbox}[colback=green!5!white,colframe=green!50!black,title=\textbf{Handoff: Next Steps},breakable]
This completes the density matrix formalism for Chapter 3. The subsequent chapters should:
\begin{itemize}
  \item \textbf{Chapter 4 (Quantum Channels)}: Develop completely positive maps, Kraus representation, Lindblad master equation for dissipative noetic dynamics
  \item \textbf{Chapter 5 (Entanglement)}: Use reduced density matrices for entanglement entropy $S(\rho_A) = S(\Trace_B(\rho_{AB}))$; develop entanglement measures beyond von Neumann entropy (concurrence, negativity)
  \item \textbf{Chapter 6 (Measurement)}: Connect density matrix collapse to POVM formalism; measurement-induced state updates
\end{itemize}

Key results established in this chapter:
\begin{enumerate}
  \item \textbf{Density operator formalism}: Pure states $\rho_\psi = \qket{\psi}\qbra{\psi}$, mixed states, convexity (Def.~\ref{def:density-operator}--\ref{def:statistical-mixture})
  \item \textbf{Von Neumann entropy}: $S(\rho) = -\Trace(\rho \log \rho)$ with properties (Def.~\ref{def:von-neumann-entropy}, Prop.~\ref{prop:entropy-properties})
  \item \textbf{Entropy-Information Correspondence Theorem~\ref{thm:entropy-info-correspondence}}: $S(\rho) \leq H_{\mathrm{cl}}(\rho)$ connecting quantum and classical information measures
  \item \textbf{Von Neumann equation}: $\frac{d\rho}{dt} = -i\omega_k[\qop{\nu}_k, \rho]$ for continuous dynamics (Def.~\ref{def:von-neumann-equation})
  \item \textbf{Steady-state characterization}: Block-diagonal in noetic eigenspaces (Thm.~\ref{thm:steady-state-char})
  \item \textbf{Partial trace and reduced states}: World-wise reduction, information loss (Thm.~\ref{thm:info-loss-trace})
  \item \textbf{Entropy decomposition}: $S(\rho) = S_W(\rho) + S_{E|W}(\rho)$ for classical states (Thm.~\ref{thm:entropy-decomp})
\end{enumerate}
\end{tcolorbox}

%% ============================================================================
%% CHAPTER 4: QUANTUM CHANNELS
%% ============================================================================
\chapter{Quantum Channels and Noetic Dynamics}
\label{ch:quantum-channels}

\begin{tcolorbox}[colback=blue!5!white,colframe=blue!75!black,title=\vnewthree]
This chapter develops quantum channel theory for TKS, describing the most general evolution of noetic states including noise and decoherence.
\end{tcolorbox}

%% ----------------------------------------------------------------------------
\section{Completely Positive Maps}
\label{sec:cp-maps}

This section develops the mathematical framework for quantum channels---the most general physically realizable transformations of quantum states. We begin with the abstract theory and then specialize to TKS.

\subsection{Positive and Completely Positive Maps}

\begin{definition}[Linear Map on Operators]
\label{def:linear-map-operators}
A \textbf{linear map} (or \textbf{superoperator}) on $\BoundedOp(\Hspace_{\Ideas})$ is a function:
\[
\mathcal{E} : \BoundedOp(\Hspace_{\Ideas}) \to \BoundedOp(\Hspace_{\Ideas})
\]
satisfying $\mathcal{E}(\alpha A + \beta B) = \alpha \mathcal{E}(A) + \beta \mathcal{E}(B)$ for all $A, B \in \BoundedOp(\Hspace_{\Ideas})$ and $\alpha, \beta \in \mathbb{C}$.
\end{definition}

\begin{definition}[Positive Map]
\label{def:positive-map}
A linear map $\mathcal{E} : \BoundedOp(\Hspace_{\Ideas}) \to \BoundedOp(\Hspace_{\Ideas})$ is \textbf{positive} if:
\[
A \geq 0 \implies \mathcal{E}(A) \geq 0
\]
That is, $\mathcal{E}$ maps positive semi-definite operators to positive semi-definite operators.
\end{definition}

\begin{remark}[Positivity is Necessary but Not Sufficient]
\label{rem:positivity-not-sufficient}
Positivity ensures that density matrices (which are positive with trace 1) are mapped to positive operators. However, positivity alone is \emph{not} sufficient for a map to represent a valid physical operation. The issue arises when the system is part of a larger composite system.
\end{remark}

\begin{definition}[Completely Positive Map]
\label{def:cp-map}
A linear map $\mathcal{E} : \BoundedOp(\Hspace_{\Ideas}) \to \BoundedOp(\Hspace_{\Ideas})$ is \textbf{completely positive (CP)} if for every $n \geq 1$, the extended map:
\[
\mathcal{E} \otimes \mathrm{id}_n : \BoundedOp(\Hspace_{\Ideas} \otimes \mathbb{C}^n) \to \BoundedOp(\Hspace_{\Ideas} \otimes \mathbb{C}^n)
\]
is positive. Here $\mathrm{id}_n$ is the identity map on $\BoundedOp(\mathbb{C}^n)$.

Equivalently, $\mathcal{E}$ is CP iff $(\mathcal{E} \otimes \mathrm{id}_n)(\sigma) \geq 0$ for all positive $\sigma \in \BoundedOp(\Hspace_{\Ideas} \otimes \mathbb{C}^n)$ and all $n$.
\end{definition}

\begin{proposition}[CP Implies Positive]
\label{prop:cp-implies-positive}
Every completely positive map is positive (take $n = 1$), but the converse is false.
\end{proposition}

\begin{example}[Transpose is Positive but Not CP]
\label{ex:transpose-not-cp}
The \textbf{transpose map} $T : A \mapsto A^T$ (in the Idea basis) is positive but \emph{not} completely positive. To see this, consider the $2 \times 2$ case: the maximally entangled state $\qket{\Phi^+} = \frac{1}{\sqrt{2}}(\qket{00} + \qket{11})$ has density matrix $\rho = \qket{\Phi^+}\qbra{\Phi^+}$. Then $(T \otimes \mathrm{id})(\rho)$ has a negative eigenvalue $-\frac{1}{2}$, so $T$ is not CP.

In TKS, this shows that ``transposing'' an Idea state (reversing coherences) is not a physically realizable operation.
\end{example}

\begin{definition}[Trace-Preserving Map]
\label{def:trace-preserving}
A linear map $\mathcal{E}$ is \textbf{trace-preserving (TP)} if:
\[
\Trace(\mathcal{E}(A)) = \Trace(A) \quad \text{for all } A \in \BoundedOp(\Hspace_{\Ideas})
\]
\end{definition}

\begin{definition}[Quantum Channel (CPTP Map)]
\label{def:quantum-channel}
A \textbf{quantum channel} is a linear map $\mathcal{E} : \BoundedOp(\Hspace_{\Ideas}) \to \BoundedOp(\Hspace_{\Ideas})$ that is:
\begin{enumerate}
  \item \textbf{Completely positive (CP)}: $(\mathcal{E} \otimes \mathrm{id}_n)(\sigma) \geq 0$ for all positive $\sigma$ and all $n$
  \item \textbf{Trace-preserving (TP)}: $\Trace(\mathcal{E}(\rho)) = \Trace(\rho)$ for all $\rho$
\end{enumerate}
We denote the set of all quantum channels on $\Hspace_{\Ideas}$ by $\mathrm{CPTP}(\Hspace_{\Ideas})$.
\end{definition}

\begin{proposition}[Channels Map Density Operators to Density Operators]
\label{prop:channel-preserves-density}
If $\mathcal{E}$ is a quantum channel and $\rho \in \mathcal{D}(\Hspace_{\Ideas})$ is a density operator, then $\mathcal{E}(\rho) \in \mathcal{D}(\Hspace_{\Ideas})$.
\end{proposition}

\begin{proof}
By complete positivity (with $n=1$), $\mathcal{E}(\rho) \geq 0$ since $\rho \geq 0$. By trace preservation, $\Trace(\mathcal{E}(\rho)) = \Trace(\rho) = 1$. Finally, $\mathcal{E}(\rho)$ is Hermitian because $\mathcal{E}$ preserves Hermiticity (a consequence of CP and TP).
\end{proof}

\subsection{Choi-Jamio\l{}kowski Isomorphism}

The Choi-Jamio\l{}kowski isomorphism provides a powerful correspondence between channels and states, enabling us to study channels via their ``Choi matrices.''

\begin{definition}[Maximally Entangled State]
\label{def:max-entangled-channel}
The \textbf{unnormalized maximally entangled state} on $\Hspace_{\Ideas} \otimes \Hspace_{\Ideas}$ is:
\[
\qket{\Omega} = \sum_{I \in \Ideas} \qket{I} \otimes \qket{I} = \sum_{W,n} \qket{Wn} \otimes \qket{Wn}
\]
This satisfies $\qbra{\Omega}\Omega\rangle = 40$.
\end{definition}

\begin{definition}[Choi Matrix]
\label{def:choi-matrix}
For a linear map $\mathcal{E} : \BoundedOp(\Hspace_{\Ideas}) \to \BoundedOp(\Hspace_{\Ideas})$, the \textbf{Choi matrix} (or \textbf{Choi-Jamio\l{}kowski matrix}) is:
\[
J_{\mathcal{E}} = (\mathcal{E} \otimes \mathrm{id})(\qket{\Omega}\qbra{\Omega}) = \sum_{I, J \in \Ideas} \mathcal{E}(\qket{I}\qbra{J}) \otimes \qket{I}\qbra{J}
\]
This is an operator on $\Hspace_{\Ideas} \otimes \Hspace_{\Ideas}$.
\end{definition}

\begin{theorem}[Choi-Jamio\l{}kowski Isomorphism]
\label{thm:choi-jamiolkowski}
The map $\mathcal{E} \mapsto J_{\mathcal{E}}$ is a linear isomorphism between:
\begin{itemize}
  \item Linear maps $\mathcal{E} : \BoundedOp(\Hspace_{\Ideas}) \to \BoundedOp(\Hspace_{\Ideas})$
  \item Operators $J \in \BoundedOp(\Hspace_{\Ideas} \otimes \Hspace_{\Ideas})$
\end{itemize}
Moreover:
\begin{enumerate}
  \item $\mathcal{E}$ is \textbf{completely positive} iff $J_{\mathcal{E}} \geq 0$ (positive semi-definite)
  \item $\mathcal{E}$ is \textbf{trace-preserving} iff $\Trace_1(J_{\mathcal{E}}) = \mathbf{I}$ (partial trace over first system is identity)
  \item $\mathcal{E}$ is a \textbf{quantum channel} iff $J_{\mathcal{E}} \geq 0$ and $\Trace_1(J_{\mathcal{E}}) = \mathbf{I}$
\end{enumerate}
\end{theorem}

\begin{proof}
\textbf{Isomorphism}: The vector space of linear maps has dimension $40^4$, as does the space of operators on $\Hspace_{\Ideas} \otimes \Hspace_{\Ideas}$. The map $\mathcal{E} \mapsto J_{\mathcal{E}}$ is linear and injective (since $\mathcal{E}$ can be recovered from $J_{\mathcal{E}}$), hence an isomorphism.

\textbf{(1) CP iff $J_{\mathcal{E}} \geq 0$}: ($\Rightarrow$) If $\mathcal{E}$ is CP, then $(\mathcal{E} \otimes \mathrm{id})(\qket{\Omega}\qbra{\Omega}) \geq 0$ since $\qket{\Omega}\qbra{\Omega} \geq 0$. ($\Leftarrow$) If $J_{\mathcal{E}} \geq 0$, one can show via the Kraus representation (Theorem~\ref{thm:kraus}) that $\mathcal{E}$ is CP.

\textbf{(2) TP iff $\Trace_1(J_{\mathcal{E}}) = \mathbf{I}$}: Direct calculation shows $\Trace_1(J_{\mathcal{E}}) = \sum_I \mathcal{E}(\qket{I}\qbra{I})$. This equals $\mathbf{I}$ iff $\Trace(\mathcal{E}(A)) = \Trace(A)$ for all $A$.
\end{proof}

\begin{corollary}[Choi Rank and Channel Properties]
\label{cor:choi-rank}
For a quantum channel $\mathcal{E}$:
\begin{enumerate}
  \item The rank of $J_{\mathcal{E}}$ is called the \textbf{Kraus rank} of $\mathcal{E}$
  \item $\mathcal{E}$ is a unitary channel iff $\mathrm{rank}(J_{\mathcal{E}}) = 1$
  \item The maximum Kraus rank for channels on $\Hspace_{\Ideas}$ is $40^2 = 1600$
\end{enumerate}
\end{corollary}

\begin{definition}[Choi Matrix for Noetic Operator]
\label{def:choi-noetic}
For a unitary noetic operator $\qop{\nu}_k$ (corresponding to a bijective permutation $\mathfrak{n}_k$), the associated unitary channel $\mathcal{U}_k(\rho) = \qop{\nu}_k \rho \qop{\nu}_k^\dagger$ has Choi matrix:
\[
J_{\mathcal{U}_k} = (\qop{\nu}_k \otimes \mathbf{I}) \qket{\Omega}\qbra{\Omega} (\qop{\nu}_k^\dagger \otimes \mathbf{I}) = \sum_{I,J} \qket{\mathfrak{n}_k(I)}\qbra{\mathfrak{n}_k(J)} \otimes \qket{I}\qbra{J}
\]
This has rank 1 (it is a pure state up to normalization).
\end{definition}

%% ----------------------------------------------------------------------------
\section{Kraus Representation}
\label{sec:kraus-representation}

The Kraus representation theorem provides the most useful characterization of quantum channels for computational purposes.

\begin{theorem}[Kraus Representation Theorem]
\label{thm:kraus}
A linear map $\mathcal{E} : \BoundedOp(\Hspace_{\Ideas}) \to \BoundedOp(\Hspace_{\Ideas})$ is completely positive if and only if there exist operators $\{K_i\}_{i=1}^{r} \subset \BoundedOp(\Hspace_{\Ideas})$ such that:
\[
\mathcal{E}(\rho) = \sum_{i=1}^{r} K_i \rho K_i^\dagger
\]
The operators $\{K_i\}$ are called \textbf{Kraus operators} (or \textbf{operation elements}).

Moreover:
\begin{enumerate}
  \item $\mathcal{E}$ is \textbf{trace-preserving} iff $\sum_{i=1}^{r} K_i^\dagger K_i = \mathbf{I}$
  \item $\mathcal{E}$ is \textbf{trace-nonincreasing} iff $\sum_{i=1}^{r} K_i^\dagger K_i \leq \mathbf{I}$
  \item The minimum number of Kraus operators needed equals $\mathrm{rank}(J_{\mathcal{E}})$
\end{enumerate}
\end{theorem}

\begin{proof}
($\Leftarrow$) Given Kraus operators, we verify CP: For any positive $\sigma \in \BoundedOp(\Hspace_{\Ideas} \otimes \mathbb{C}^n)$:
\[
(\mathcal{E} \otimes \mathrm{id})(\sigma) = \sum_i (K_i \otimes \mathbf{I}_n) \sigma (K_i^\dagger \otimes \mathbf{I}_n) \geq 0
\]
since each term $(K_i \otimes \mathbf{I}_n) \sigma (K_i \otimes \mathbf{I}_n)^\dagger$ is positive.

($\Rightarrow$) If $\mathcal{E}$ is CP, then $J_{\mathcal{E}} \geq 0$ by Theorem~\ref{thm:choi-jamiolkowski}. Let $J_{\mathcal{E}} = \sum_i \qket{\psi_i}\qbra{\psi_i}$ be the spectral decomposition. Each $\qket{\psi_i} \in \Hspace_{\Ideas} \otimes \Hspace_{\Ideas}$ defines a Kraus operator $K_i$ via the correspondence $\qket{\psi_i} \leftrightarrow K_i$ (viewing $\qket{\psi_i}$ as a matrix).

For trace preservation: $\Trace(\mathcal{E}(\rho)) = \Trace(\sum_i K_i \rho K_i^\dagger) = \Trace(\rho \sum_i K_i^\dagger K_i)$. This equals $\Trace(\rho)$ for all $\rho$ iff $\sum_i K_i^\dagger K_i = \mathbf{I}$.
\end{proof}

\begin{remark}[Non-Uniqueness of Kraus Representation]
\label{rem:kraus-nonunique}
The Kraus representation is not unique. If $\{K_i\}_{i=1}^{r}$ is a Kraus representation of $\mathcal{E}$, then so is $\{K'_j\}_{j=1}^{s}$ where:
\[
K'_j = \sum_{i=1}^{r} U_{ji} K_i
\]
for any isometry $U : \mathbb{C}^r \to \mathbb{C}^s$ (i.e., $U^\dagger U = \mathbf{I}_r$).
\end{remark}

\begin{definition}[Unitary Channel]
\label{def:unitary-channel}
A \textbf{unitary channel} has a single Kraus operator that is unitary:
\[
\mathcal{U}(\rho) = U \rho U^\dagger \quad \text{where } U^\dagger U = U U^\dagger = \mathbf{I}
\]
This is the only type of channel that is reversible (has an inverse that is also a channel).
\end{definition}

\begin{proposition}[Noetic Operators as Kraus Operators]
\label{prop:noetic-kraus}
For each noetic $\noe{k}$ with lifted operator $\qop{\nu}_k$:
\begin{enumerate}
  \item If $\mathfrak{n}_k$ is a bijection, then $\{\qop{\nu}_k\}$ is a single-element Kraus representation of a unitary channel
  \item If $\mathfrak{n}_k$ is not injective, then $\qop{\nu}_k$ is not unitary, and $\qop{\nu}_k^\dagger \qop{\nu}_k \neq \mathbf{I}$
\end{enumerate}
\end{proposition}

\begin{proof}
(1) For bijective $\mathfrak{n}_k$: $\qop{\nu}_k^\dagger \qop{\nu}_k = \sum_I \qket{I}\qbra{\mathfrak{n}_k(I)} \cdot \sum_J \qket{\mathfrak{n}_k(J)}\qbra{J} = \sum_I \qket{I}\qbra{I} = \mathbf{I}$.

(2) If $\mathfrak{n}_k(I_1) = \mathfrak{n}_k(I_2) = J$ for $I_1 \neq I_2$, then the $J$-th column of $\qop{\nu}_k$ has two non-zero entries, making it non-unitary.
\end{proof}

\begin{definition}[Random Noetic Channel]
\label{def:random-noetic-channel}
The \textbf{random noetic channel} with probability distribution $\{p_k\}_{k=0}^{9}$ over noetics is:
\[
\mathcal{E}_{\mathrm{rand}}(\rho) = \sum_{k=0}^{9} p_k \, \qop{\nu}_k \rho \qop{\nu}_k^\dagger
\]
This has Kraus operators $\{K_k = \sqrt{p_k} \, \qop{\nu}_k\}_{k=0}^{9}$ satisfying:
\[
\sum_{k=0}^{9} K_k^\dagger K_k = \sum_{k=0}^{9} p_k \, \qop{\nu}_k^\dagger \qop{\nu}_k = \mathbf{I}
\]
(when all noetics are bijective).
\end{definition}

\begin{example}[Uniform Noetic Channel]
\label{ex:uniform-noetic}
The \textbf{uniform noetic channel} applies each noetic with equal probability:
\[
\mathcal{E}_{\mathrm{unif}}(\rho) = \frac{1}{10} \sum_{k=0}^{9} \qop{\nu}_k \rho \qop{\nu}_k^\dagger
\]
This represents maximal uncertainty about which noetic transformation occurred.
\end{example}

\begin{definition}[ACBE Channel (from v7.2)]
\label{def:acbe-channel-v73}
The \textbf{ACBE quantum channel} implements the Spiritual-to-Physical descent:
\[
\mathcal{E}_{\mathrm{ACBE}}(\rho) = \sum_{n=1}^{10} K_n \rho K_n^\dagger
\]
with Kraus operators $K_n = \qket{Dn}\qbra{An}$. These satisfy:
\[
\sum_{n=1}^{10} K_n^\dagger K_n = \sum_{n=1}^{10} \qket{An}\qbra{An} = \Projector_A
\]
This is \emph{not} trace-preserving on all of $\Hspace_{\Ideas}$, only on $\Hspace_A$. On full $\Hspace_{\Ideas}$, it is trace-nonincreasing.
\end{definition}

\begin{proposition}[ACBE Channel Properties]
\label{prop:acbe-properties}
The ACBE channel satisfies:
\begin{enumerate}
  \item On $\Hspace_A$: $\Trace(\mathcal{E}_{\mathrm{ACBE}}(\rho)) = \Trace(\Projector_A \rho)$ (trace-preserving on World A)
  \item \textbf{World-shifting}: $\mathcal{E}_{\mathrm{ACBE}}(\qket{An}\qbra{Am}) = \qket{Dn}\qbra{Dm}$
  \item \textbf{Coherence-preserving within A}: Off-diagonal coherences between $\qket{An}$ and $\qket{Am}$ are preserved but shifted to D
  \item \textbf{Annihilating on B, C, D}: $\mathcal{E}_{\mathrm{ACBE}}(\rho) = 0$ if $\rho$ is supported on $\Hspace_B \oplus \Hspace_C \oplus \Hspace_D$
\end{enumerate}
\end{proposition}

%% ----------------------------------------------------------------------------
\section{Noetic Channels}
\label{sec:noetic-channels}

We now develop the standard noetic channels that arise in TKS dynamics, specializing quantum information constructs to the 40-dimensional Idea space.

\subsection{Depolarizing Channel}

\begin{definition}[Noetic Depolarizing Channel]
\label{def:depolarizing}
The \textbf{noetic depolarizing channel} with parameter $p \in [0,1]$ is:
\[
\mathcal{D}_p(\rho) = (1-p) \rho + p \cdot \frac{\mathbf{I}}{40}
\]
This replaces the state with the maximally mixed state $\rho_{\mathrm{mix}} = \frac{1}{40}\mathbf{I}$ with probability $p$.
\end{definition}

\begin{proposition}[Depolarizing Channel Kraus Representation]
\label{prop:depolarizing-kraus}
The depolarizing channel $\mathcal{D}_p$ has Kraus representation:
\[
\mathcal{D}_p(\rho) = (1 - \tfrac{40p}{39}) \rho + \frac{p}{39} \sum_{I \neq J} \qket{I}\qbra{J} \rho \qket{J}\qbra{I}
\]
with Kraus operators:
\begin{itemize}
  \item $K_0 = \sqrt{1 - \frac{40p}{39}} \mathbf{I}$
  \item $K_{IJ} = \sqrt{\frac{p}{39 \cdot 40}} \qket{I}\qbra{J}$ for $I \neq J$
\end{itemize}
(Here we use $39 \cdot 40$ swap operators plus the identity.)
\end{proposition}

\begin{proposition}[Depolarizing Channel Properties]
\label{prop:depolarizing-properties}
The depolarizing channel satisfies:
\begin{enumerate}
  \item \textbf{Fixed point}: $\mathcal{D}_p(\rho_{\mathrm{mix}}) = \rho_{\mathrm{mix}}$ (maximally mixed state is fixed)
  \item \textbf{Contraction}: $\|\mathcal{D}_p(\rho) - \rho_{\mathrm{mix}}\|_1 = (1-p)\|\rho - \rho_{\mathrm{mix}}\|_1$
  \item \textbf{Entropy increase}: $S(\mathcal{D}_p(\rho)) \geq S(\rho)$ with equality iff $p = 0$ or $\rho = \rho_{\mathrm{mix}}$
  \item \textbf{Composition}: $\mathcal{D}_p \circ \mathcal{D}_q = \mathcal{D}_{p + q - pq}$
\end{enumerate}
\end{proposition}

\begin{proof}
(1) Direct: $\mathcal{D}_p(\rho_{\mathrm{mix}}) = (1-p)\rho_{\mathrm{mix}} + p \rho_{\mathrm{mix}} = \rho_{\mathrm{mix}}$.

(2) The channel is a convex combination with the maximally mixed state, which contracts distance to it.

(3) The von Neumann entropy $S(\rho) = -\Trace(\rho \log \rho)$ is concave, so mixing with the maximum-entropy state increases entropy.

(4) $\mathcal{D}_p(\mathcal{D}_q(\rho)) = (1-p)[(1-q)\rho + q\rho_{\mathrm{mix}}] + p\rho_{\mathrm{mix}} = (1-p)(1-q)\rho + [1-(1-p)(1-q)]\rho_{\mathrm{mix}}$.
\end{proof}

\subsection{Dephasing Channel}

\begin{definition}[Noetic Dephasing Channel]
\label{def:dephasing-channel}
The \textbf{noetic dephasing channel} (complete dephasing in the Idea basis) is:
\[
\mathcal{E}_{\mathrm{deph}}(\rho) = \sum_{I \in \Ideas} \qket{I}\qbra{I} \rho \qket{I}\qbra{I} = \sum_{I} \rho_{II} \qket{I}\qbra{I}
\]
This destroys all off-diagonal coherences, leaving only diagonal (classical) probabilities.
\end{definition}

\begin{proposition}[Dephasing Kraus Representation]
\label{prop:dephasing-kraus}
The dephasing channel has Kraus operators $\{K_I = \qket{I}\qbra{I}\}_{I \in \Ideas}$ (the 40 Idea projectors), satisfying:
\[
\sum_{I \in \Ideas} K_I^\dagger K_I = \sum_I \qket{I}\qbra{I} = \mathbf{I}
\]
\end{proposition}

\begin{definition}[Partial Dephasing Channel]
\label{def:partial-dephasing}
The \textbf{partial dephasing channel} with parameter $\gamma \in [0,1]$ is:
\[
\mathcal{E}_{\mathrm{deph}}^\gamma(\rho) = \gamma \mathcal{E}_{\mathrm{deph}}(\rho) + (1-\gamma) \rho
\]
This interpolates between identity ($\gamma = 0$) and complete dephasing ($\gamma = 1$).

In matrix form, the off-diagonal elements are damped:
\[
[\mathcal{E}_{\mathrm{deph}}^\gamma(\rho)]_{IJ} = \begin{cases}
\rho_{IJ} & \text{if } I = J \\
(1-\gamma) \rho_{IJ} & \text{if } I \neq J
\end{cases}
\]
\end{definition}

\begin{proposition}[Dephasing Properties]
\label{prop:dephasing-properties}
The dephasing channel satisfies:
\begin{enumerate}
  \item \textbf{Idempotent}: $\mathcal{E}_{\mathrm{deph}} \circ \mathcal{E}_{\mathrm{deph}} = \mathcal{E}_{\mathrm{deph}}$
  \item \textbf{Fixed points}: All classical states $\rho_{\mathrm{cl}}$ (diagonal in Idea basis) are fixed
  \item \textbf{Entropy}: $S(\mathcal{E}_{\mathrm{deph}}(\rho)) = H_{\mathrm{cl}}(\rho) \geq S(\rho)$ (Theorem~\ref{thm:entropy-info-correspondence})
  \item \textbf{Quantum-to-classical}: $\mathcal{E}_{\mathrm{deph}}$ maps $\mathcal{D}(\Hspace_{\Ideas})$ onto the simplex of classical distributions over $\Ideas$
\end{enumerate}
\end{proposition}

\subsection{World Dephasing}

\begin{definition}[World Dephasing Channel]
\label{def:world-dephasing}
The \textbf{world dephasing channel} destroys coherences between different Worlds but preserves coherences within each World:
\[
\mathcal{E}_{\mathrm{W-deph}}(\rho) = \sum_{W \in \{A,B,C,D\}} \Projector_W \rho \Projector_W
\]
This has Kraus operators $\{K_W = \Projector_W\}_{W \in \{A,B,C,D\}}$.
\end{definition}

\begin{proposition}[World Dephasing vs Full Dephasing]
\label{prop:world-vs-full-deph}
The world dephasing channel preserves more quantum information than full dephasing:
\begin{enumerate}
  \item $\mathcal{E}_{\mathrm{deph}} = \mathcal{E}_{\mathrm{deph}} \circ \mathcal{E}_{\mathrm{W-deph}}$ (full dephasing factors through world dephasing)
  \item $S(\rho) \leq S(\mathcal{E}_{\mathrm{W-deph}}(\rho)) \leq S(\mathcal{E}_{\mathrm{deph}}(\rho))$
  \item If $\rho = \sum_W \rho_W$ with each $\rho_W$ supported on $\Hspace_W$, then $\mathcal{E}_{\mathrm{W-deph}}(\rho) = \rho$
\end{enumerate}
\end{proposition}

\subsection{Amplitude Damping}

\begin{definition}[World Amplitude Damping]
\label{def:amplitude-damping}
The \textbf{world amplitude damping channel} with parameter $\gamma \in [0,1]$ models decay from higher Worlds to lower Worlds. For the $A \to D$ decay:
\[
\mathcal{E}_{\mathrm{damp}}^{A \to D}(\rho) = K_0 \rho K_0^\dagger + K_1 \rho K_1^\dagger
\]
with Kraus operators:
\begin{align}
K_0 &= \Projector_B + \Projector_C + \Projector_D + \sqrt{1-\gamma} \Projector_A \\
K_1 &= \sqrt{\gamma} \sum_{n=1}^{10} \qket{Dn}\qbra{An}
\end{align}
\end{definition}

\begin{proposition}[Amplitude Damping Properties]
\label{prop:amp-damp-properties}
The amplitude damping channel satisfies:
\begin{enumerate}
  \item \textbf{Population transfer}: Probability flows from World A to World D at rate $\gamma$
  \item \textbf{Coherence decay}: Off-diagonal coherences between A and other Worlds decay as $\sqrt{1-\gamma}$
  \item \textbf{Steady state}: World D states are fixed points
  \item \textbf{Connection to ACBE}: At $\gamma = 1$, this becomes the ACBE channel (restricted to A)
\end{enumerate}
\end{proposition}

\begin{proof}
For $\rho = \qket{An}\qbra{An}$:
\begin{align}
\mathcal{E}_{\mathrm{damp}}^{A \to D}(\rho) &= (1-\gamma)\qket{An}\qbra{An} + \gamma \qket{Dn}\qbra{Dn}
\end{align}
showing population transfer. For $\rho = \qket{An}\qbra{Bm}$:
\begin{align}
\mathcal{E}_{\mathrm{damp}}^{A \to D}(\rho) &= \sqrt{1-\gamma} \qket{An}\qbra{Bm}
\end{align}
showing coherence decay.
\end{proof}

\subsection{Fixed Points and Attractors}

\begin{definition}[Channel Fixed Point]
\label{def:channel-fixed-point}
A density operator $\rho_{\mathrm{fix}}$ is a \textbf{fixed point} of channel $\mathcal{E}$ if:
\[
\mathcal{E}(\rho_{\mathrm{fix}}) = \rho_{\mathrm{fix}}
\]
The set of fixed points is denoted $\mathrm{Fix}(\mathcal{E})$.
\end{definition}

\begin{theorem}[Channel Fixed Points and Classical Attractors]
\label{thm:channel-classical-attractor}
\textbf{(Main Theorem: Connecting Quantum Fixed Points to Classical Dynamics)}

Let $\mathcal{E}$ be a quantum channel on $\Hspace_{\Ideas}$ of the form:
\[
\mathcal{E}(\rho) = \sum_{k} p_k \, \qop{\nu}_k \rho \qop{\nu}_k^\dagger
\]
where $\{p_k\}$ is a probability distribution over noetics with bijective $\mathfrak{n}_k$. Let $\mu$ be a probability measure on $\Ideas$ and let $\rho_\mu = \sum_I \mu(\{I\}) \qket{I}\qbra{I}$ be the corresponding classical (diagonal) density matrix. Then:

\begin{enumerate}
  \item \textbf{Classical-Quantum Correspondence}: $\rho_\mu$ is a fixed point of $\mathcal{E}$ iff $\mu$ is invariant under the classical random dynamical system with transition kernel:
  \[
  P(I \to J) = \sum_{k : \mathfrak{n}_k(I) = J} p_k
  \]

  \item \textbf{Fixed-Point Existence}: $\mathrm{Fix}(\mathcal{E}) \neq \emptyset$. Every quantum channel on a finite-dimensional space has at least one fixed point.

  \item \textbf{Unique Classical Fixed Point}: If the classical Markov chain with kernel $P$ is irreducible and aperiodic, there is a unique classical fixed point $\rho_{\mu_*}$ corresponding to the stationary distribution $\mu_*$.

  \item \textbf{Quantum Fixed Points}: In general, $\mathrm{Fix}(\mathcal{E})$ may contain quantum (non-diagonal) fixed points in addition to classical ones. These correspond to ``dark states'' with coherences that are invariant under all $\qop{\nu}_k$ in the mixture.
\end{enumerate}
\end{theorem}

\begin{proof}
\textbf{(1)} For classical $\rho_\mu = \sum_I \mu_I \qket{I}\qbra{I}$:
\begin{align}
\mathcal{E}(\rho_\mu) &= \sum_k p_k \sum_I \mu_I \qket{\mathfrak{n}_k(I)}\qbra{\mathfrak{n}_k(I)} \\
&= \sum_J \left( \sum_k p_k \sum_{I : \mathfrak{n}_k(I) = J} \mu_I \right) \qket{J}\qbra{J} \\
&= \sum_J \left( \sum_I P(I \to J) \mu_I \right) \qket{J}\qbra{J}
\end{align}
This equals $\rho_\mu$ iff $\sum_I P(I \to J) \mu_I = \mu_J$ for all $J$, i.e., $\mu$ is stationary for $P$.

\textbf{(2)} By Brouwer's fixed-point theorem: $\mathcal{E}$ is a continuous map on the compact convex set $\mathcal{D}(\Hspace_{\Ideas})$, so it has a fixed point.

\textbf{(3)} Standard Markov chain theory: irreducibility and aperiodicity imply a unique stationary distribution.

\textbf{(4)} A quantum fixed point $\rho_{\mathrm{fix}}$ satisfies $\sum_k p_k \qop{\nu}_k \rho_{\mathrm{fix}} \qop{\nu}_k^\dagger = \rho_{\mathrm{fix}}$. If $\rho_{\mathrm{fix}}$ has off-diagonal coherence $\rho_{IJ} \neq 0$ for $I \neq J$, then we need $\sum_k p_k \rho_{\mathfrak{n}_k^{-1}(I), \mathfrak{n}_k^{-1}(J)} = \rho_{IJ}$, which can occur for special ``dark'' coherences.
\end{proof}

\begin{corollary}[Connection to v7.2 Invariant Measures]
\label{cor:invariant-measure-channel}
The classical fixed points of random noetic channels correspond exactly to the $\noe{k}$-invariant measures of Definition~\ref{def:invariant-measure} (v7.2), generalized to mixtures of noetics.
\end{corollary}

%% ----------------------------------------------------------------------------
\section{Channel Composition}
\label{sec:channel-composition}

\begin{definition}[Sequential Composition]
\label{def:sequential-composition}
For quantum channels $\mathcal{E}_1, \mathcal{E}_2$, the \textbf{sequential composition} is:
\[
(\mathcal{E}_2 \circ \mathcal{E}_1)(\rho) = \mathcal{E}_2(\mathcal{E}_1(\rho))
\]
If $\mathcal{E}_1$ has Kraus operators $\{K_i\}$ and $\mathcal{E}_2$ has Kraus operators $\{L_j\}$, then $\mathcal{E}_2 \circ \mathcal{E}_1$ has Kraus operators $\{L_j K_i\}_{i,j}$.
\end{definition}

\begin{proposition}[CPTP is Closed under Composition]
\label{prop:cptp-composition}
The set $\mathrm{CPTP}(\Hspace_{\Ideas})$ is closed under sequential composition: if $\mathcal{E}_1, \mathcal{E}_2$ are quantum channels, then so is $\mathcal{E}_2 \circ \mathcal{E}_1$.
\end{proposition}

\begin{proof}
CP: $(L_j K_i) \rho (L_j K_i)^\dagger = L_j (K_i \rho K_i^\dagger) L_j^\dagger$ preserves positivity.

TP: $\sum_{i,j} (L_j K_i)^\dagger (L_j K_i) = \sum_i K_i^\dagger (\sum_j L_j^\dagger L_j) K_i = \sum_i K_i^\dagger K_i = \mathbf{I}$.
\end{proof}

\begin{definition}[Convex Combination of Channels]
\label{def:convex-channels}
For channels $\{\mathcal{E}_k\}$ and probabilities $\{p_k\}$, the \textbf{convex combination}:
\[
\mathcal{E} = \sum_k p_k \mathcal{E}_k
\]
is also a quantum channel.
\end{definition}

\begin{proposition}[Noetic Cascade as Channel Composition]
\label{prop:noetic-cascade}
The ACBE cascade corresponds to sequential composition of single-step channels:
\[
\mathcal{E}_{\mathrm{ACBE}} = \mathcal{E}_{C \to D} \circ \mathcal{E}_{B \to C} \circ \mathcal{E}_{A \to B}
\]
where each $\mathcal{E}_{W \to W'}$ is a world-transition channel.
\end{proposition}

%% ----------------------------------------------------------------------------
\section{Lindblad Master Equation}
\label{sec:master-equations}

For continuous-time evolution beyond simple unitary dynamics, we develop the Lindblad (or GKSL) master equation framework.

\subsection{Quantum Dynamical Semigroups}

\begin{definition}[Quantum Dynamical Semigroup]
\label{def:qds}
A \textbf{quantum dynamical semigroup} is a family of quantum channels $\{\mathcal{E}_t\}_{t \geq 0}$ satisfying:
\begin{enumerate}
  \item \textbf{Identity}: $\mathcal{E}_0 = \mathrm{id}$
  \item \textbf{Semigroup}: $\mathcal{E}_{t+s} = \mathcal{E}_t \circ \mathcal{E}_s$ for all $t, s \geq 0$
  \item \textbf{Continuity}: $\lim_{t \to 0^+} \mathcal{E}_t(\rho) = \rho$ for all $\rho$
\end{enumerate}
\end{definition}

\begin{definition}[Generator of a Semigroup]
\label{def:generator}
The \textbf{generator} (or \textbf{Lindbladian}) of a quantum dynamical semigroup is:
\[
\mathcal{L}(\rho) = \lim_{t \to 0^+} \frac{\mathcal{E}_t(\rho) - \rho}{t}
\]
The evolution satisfies the \textbf{master equation}:
\[
\frac{d\rho}{dt} = \mathcal{L}(\rho)
\]
with solution $\rho(t) = \mathcal{E}_t(\rho(0)) = e^{t\mathcal{L}}(\rho(0))$.
\end{definition}

\subsection{Lindblad Form}

\begin{theorem}[Lindblad-Gorini-Kossakowski-Sudarshan (LGKS) Theorem]
\label{thm:lindblad}
A linear map $\mathcal{L}$ is the generator of a quantum dynamical semigroup iff it has the \textbf{Lindblad form}:
\[
\mathcal{L}(\rho) = -i[H, \rho] + \sum_{k=1}^{N} \left( L_k \rho L_k^\dagger - \frac{1}{2}\{L_k^\dagger L_k, \rho\} \right)
\]
where:
\begin{itemize}
  \item $H = H^\dagger$ is a Hermitian operator (the \textbf{effective Hamiltonian})
  \item $\{L_k\}$ are \textbf{Lindblad operators} (or \textbf{jump operators})
  \item $\{A, B\} = AB + BA$ is the anticommutator
  \item $N \leq 40^2 - 1 = 1599$ for $\Hspace_{\Ideas}$
\end{itemize}
\end{theorem}

\begin{proof}
(Sketch) The necessity comes from requiring $\mathcal{E}_t = e^{t\mathcal{L}}$ to be CPTP for all $t \geq 0$. Complete positivity of $\mathcal{E}_t$ requires the ``dissipator'' part $\sum_k (L_k \rho L_k^\dagger - \frac{1}{2}\{L_k^\dagger L_k, \rho\})$ to have this specific form. Trace preservation requires the coefficient of $\rho$ in the anticommutator term.
\end{proof}

\begin{definition}[Dissipator]
\label{def:dissipator}
The \textbf{dissipator} associated with Lindblad operators $\{L_k\}$ is:
\[
\mathcal{D}[\{L_k\}](\rho) = \sum_k \left( L_k \rho L_k^\dagger - \frac{1}{2}\{L_k^\dagger L_k, \rho\} \right)
\]
The Lindblad equation can be written as:
\[
\frac{d\rho}{dt} = -i[H, \rho] + \mathcal{D}[\{L_k\}](\rho)
\]
\end{definition}

\begin{remark}[Unitary vs Dissipative Dynamics]
\label{rem:unitary-vs-dissipative}
The Lindblad equation separates into:
\begin{enumerate}
  \item \textbf{Unitary part}: $-i[H, \rho]$ --- generates coherent, reversible evolution (the von Neumann equation of Section~\ref{sec:density-evolution})
  \item \textbf{Dissipative part}: $\mathcal{D}[\{L_k\}](\rho)$ --- generates irreversible evolution, entropy increase, decoherence
\end{enumerate}
Pure unitary evolution ($L_k = 0$) preserves purity; any non-trivial dissipator generically decreases purity.
\end{remark}

\subsection{Noetic Jump Operators}

\begin{definition}[Noetic Jump Operators]
\label{def:noetic-jump}
The \textbf{noetic jump operator} for transition from Idea $I$ to Idea $J$ is:
\[
L_{I \to J} = \sqrt{\gamma_{IJ}} \qket{J}\qbra{I}
\]
where $\gamma_{IJ} \geq 0$ is the transition rate. The corresponding dissipator is:
\[
\mathcal{D}[L_{I \to J}](\rho) = \gamma_{IJ} \left( \qket{J}\qbra{I} \rho \qket{I}\qbra{J} - \frac{1}{2}\{\qket{I}\qbra{I}, \rho\} \right)
\]
\end{definition}

\begin{proposition}[Effect of Noetic Jump]
\label{prop:jump-effect}
The noetic jump $L_{I \to J}$ causes:
\begin{enumerate}
  \item \textbf{Population transfer}: $\rho_{II} \to \rho_{II}(1 - \gamma_{IJ} dt)$ and $\rho_{JJ} \to \rho_{JJ} + \gamma_{IJ} \rho_{II} dt$
  \item \textbf{Coherence decay}: $\rho_{IK} \to \rho_{IK}(1 - \frac{\gamma_{IJ}}{2} dt)$ for $K \neq I$
  \item \textbf{Coherence creation}: $\rho_{JK} \to \rho_{JK} + \gamma_{IJ} \rho_{IK} dt \cdot \delta_{J \neq K}$ (coherence transfer)
\end{enumerate}
\end{proposition}

\begin{definition}[Noetic Master Equation]
\label{def:noetic-master-equation}
The general \textbf{noetic master equation} for TKS dynamics is:
\[
\frac{d\rho}{dt} = -i[\qop{H}_{\mathrm{eff}}, \rho] + \sum_{I,J} \mathcal{D}[L_{I \to J}](\rho)
\]
where $\qop{H}_{\mathrm{eff}}$ is the effective noetic Hamiltonian and the sum is over all Idea transitions.
\end{definition}

\begin{example}[ACBE Dissipative Dynamics]
\label{ex:acbe-dissipative}
The continuous-time ACBE dynamics with rate $\gamma$ has jump operators:
\[
L_n^{(A \to D)} = \sqrt{\gamma} \qket{Dn}\qbra{An} \quad \text{for } n = 1, \ldots, 10
\]
The master equation:
\[
\frac{d\rho}{dt} = \gamma \sum_{n=1}^{10} \left( \qket{Dn}\qbra{An} \rho \qket{An}\qbra{Dn} - \frac{1}{2}\{\qket{An}\qbra{An}, \rho\} \right)
\]
describes exponential decay from World A to World D with time constant $1/\gamma$.
\end{example}

\begin{proposition}[ACBE Steady State]
\label{prop:acbe-steady}
The steady state of the ACBE dissipative dynamics (Example~\ref{ex:acbe-dissipative}) starting from any initial state $\rho_0$ with support on Worlds A and D is:
\[
\rho_{\infty} = \Projector_D \rho_0 \Projector_D + \sum_n \rho_{An,An}(0) \qket{Dn}\qbra{Dn}
\]
All population initially in World A ends up in World D, while World D population is unchanged.
\end{proposition}

\subsection{Steady States of Lindblad Dynamics}

\begin{definition}[Lindblad Steady State]
\label{def:lindblad-steady}
A density operator $\rho_{\mathrm{ss}}$ is a \textbf{steady state} of the Lindblad dynamics with generator $\mathcal{L}$ if:
\[
\mathcal{L}(\rho_{\mathrm{ss}}) = 0
\]
\end{definition}

\begin{theorem}[Existence of Lindblad Steady States]
\label{thm:lindblad-steady-existence}
Every Lindblad generator on $\Hspace_{\Ideas}$ has at least one steady state $\rho_{\mathrm{ss}} \in \mathcal{D}(\Hspace_{\Ideas})$.
\end{theorem}

\begin{proof}
The quantum dynamical semigroup $\mathcal{E}_t = e^{t\mathcal{L}}$ maps $\mathcal{D}(\Hspace_{\Ideas})$ to itself. By compactness of $\mathcal{D}(\Hspace_{\Ideas})$ and Brouwer's theorem (or Cesaro averaging), a fixed point exists.
\end{proof}

\begin{definition}[Detailed Balance]
\label{def:detailed-balance}
A Lindblad generator satisfies \textbf{detailed balance} with respect to state $\rho_{\mathrm{eq}}$ if:
\[
\qbra{\phi}\mathcal{L}(\qket{\psi}\qbra{\psi})\qket{\phi} \cdot \qbra{\psi}\rho_{\mathrm{eq}}\qket{\psi} = \qbra{\psi}\mathcal{L}(\qket{\phi}\qbra{\phi})\qket{\psi} \cdot \qbra{\phi}\rho_{\mathrm{eq}}\qket{\phi}
\]
for all $\qket{\psi}, \qket{\phi}$. This is the quantum analogue of classical detailed balance.
\end{definition}

\begin{theorem}[Noetic Detailed Balance]
\label{thm:noetic-detailed-balance}
For noetic jump operators $\{L_{I \to J} = \sqrt{\gamma_{IJ}} \qket{J}\qbra{I}\}$, detailed balance with respect to $\rho_{\mathrm{eq}} = \sum_I p_I \qket{I}\qbra{I}$ holds iff:
\[
\gamma_{IJ} p_I = \gamma_{JI} p_J \quad \text{for all } I, J
\]
This is the classical detailed balance condition for the transition rates.
\end{theorem}

%% ----------------------------------------------------------------------------
%% HANDOFF NOTE
%% ----------------------------------------------------------------------------
\begin{tcolorbox}[colback=green!5!white,colframe=green!50!black,title=\textbf{Handoff: Next Steps},breakable]
This completes the quantum channels chapter (Chapter 4). The subsequent chapters should:
\begin{itemize}
  \item \textbf{Chapter 5 (Entanglement)}: Use channels to study entanglement dynamics; local operations and classical communication (LOCC); entanglement-breaking channels
  \item \textbf{Chapter 6 (Measurement)}: Channels as measurement back-action; instrument formalism; POVMs from Kraus operators
\end{itemize}

Key results established in this chapter:
\begin{enumerate}
  \item \textbf{CPTP characterization}: Quantum channels are completely positive trace-preserving maps (Def.~\ref{def:quantum-channel})
  \item \textbf{Choi-Jamio\l{}kowski isomorphism}: Channels $\leftrightarrow$ positive operators (Thm.~\ref{thm:choi-jamiolkowski})
  \item \textbf{Kraus representation}: $\mathcal{E}(\rho) = \sum_i K_i \rho K_i^\dagger$ with $\sum_i K_i^\dagger K_i = \mathbf{I}$ (Thm.~\ref{thm:kraus})
  \item \textbf{Standard noetic channels}: Depolarizing, dephasing, world dephasing, amplitude damping (Sec.~\ref{sec:noetic-channels})
  \item \textbf{Channel-Attractor Correspondence Theorem~\ref{thm:channel-classical-attractor}}: Fixed points of random noetic channels correspond to invariant measures of classical dynamics
  \item \textbf{Lindblad master equation}: $\frac{d\rho}{dt} = -i[H,\rho] + \sum_k (L_k \rho L_k^\dagger - \frac{1}{2}\{L_k^\dagger L_k, \rho\})$ (Thm.~\ref{thm:lindblad})
  \item \textbf{Noetic jump operators}: $L_{I \to J} = \sqrt{\gamma_{IJ}}\qket{J}\qbra{I}$ for Idea transitions (Def.~\ref{def:noetic-jump})
  \item \textbf{ACBE as channel}: Kraus operators $K_n = \qket{Dn}\qbra{An}$ (Def.~\ref{def:acbe-channel-v73})
  \item \textbf{Detailed balance}: Quantum-classical correspondence for equilibrium (Thm.~\ref{thm:noetic-detailed-balance})
\end{enumerate}

\textbf{Connection to v7.2}: The Channel-Attractor Correspondence Theorem~\ref{thm:channel-classical-attractor} provides the precise link between quantum fixed points and classical invariant measures (Definition~\ref{def:invariant-measure}). The ACBE channel formalism extends Definition~\ref{def:acbe-channel} from v7.2.
\end{tcolorbox}

%% ============================================================================
%% CHAPTER 5: NOETIC ENTANGLEMENT
%% ============================================================================
\chapter{Noetic Entanglement}
\label{ch:noetic-entanglement}

\begin{tcolorbox}[colback=blue!5!white,colframe=blue!75!black,title=\vnewthree]
This chapter develops entanglement theory for multi-Idea systems in TKS, extending v7.2 foundations with detailed quantification and operational interpretation.
\end{tcolorbox}

%% ----------------------------------------------------------------------------
\section{Bipartite Noetic States}
\label{sec:bipartite-states}

This section develops the theory of two-Idea systems, extending the v7.2 foundations (Definition~\ref{def:tensor-product}, \ref{def:entangled-state}) with rigorous mathematical structure.

\subsection{Tensor Product of Idea Spaces}

\begin{definition}[Bipartite Noetic Hilbert Space]
\label{def:bipartite-hilbert}
The \textbf{bipartite noetic Hilbert space} for two Ideas is the tensor product:
\[
\Hspace_{\Ideas}^{(2)} = \Hspace_{\Ideas} \otimes \Hspace_{\Ideas}
\]
with:
\begin{itemize}
  \item \textbf{Dimension}: $\dim(\Hspace_{\Ideas}^{(2)}) = 40 \times 40 = 1600$
  \item \textbf{Basis}: $\{\qket{I} \otimes \qket{J} \equiv \qket{I, J}\}_{I, J \in \Ideas}$
  \item \textbf{Inner product}: $\qbra{I, J}\ket{I', J'}\rangle = \delta_{II'} \delta_{JJ'}$
\end{itemize}
We denote the two subsystems as $\mathsf{A}$ (first Idea) and $\mathsf{B}$ (second Idea).
\end{definition}

\begin{remark}[Physical Interpretation]
\label{rem:bipartite-interpretation}
A bipartite noetic state represents two distinct Ideas that may be:
\begin{enumerate}
  \item \textbf{Two simultaneous Ideas}: Ideas held concurrently in mind
  \item \textbf{Sequential Ideas}: An Idea at time $t$ correlated with an Idea at time $t'$
  \item \textbf{Ideas of different agents}: Noetic states of two different thinkers
  \item \textbf{Idea-Acquisition pair}: An Idea correlated with its acquisition state (as in RPM)
\end{enumerate}
\end{remark}

\begin{definition}[General Bipartite State]
\label{def:general-bipartite}
A general pure state in $\Hspace_{\Ideas}^{(2)}$ has the form:
\[
\qket{\Psi} = \sum_{I, J \in \Ideas} c_{IJ} \qket{I, J}
\]
where $c_{IJ} \in \mathbb{C}$ satisfy $\sum_{I,J} |c_{IJ}|^2 = 1$. The coefficient matrix $C = (c_{IJ})$ is a $40 \times 40$ complex matrix with $\|C\|_F = 1$ (Frobenius norm).
\end{definition}

\begin{definition}[Product State]
\label{def:product-state-v73}
A \textbf{product state} (or \textbf{separable pure state}) has the form:
\[
\qket{\Psi_{\mathrm{prod}}} = \qket{\psi} \otimes \qket{\phi} = \qket{\psi, \phi}
\]
for some $\qket{\psi}, \qket{\phi} \in \Hspace_{\Ideas}$. Equivalently, the coefficient matrix has rank 1: $C = \vec{a} \vec{b}^T$ for column vectors $\vec{a}, \vec{b}$.
\end{definition}

\begin{definition}[Entangled Pure State]
\label{def:entangled-pure}
A pure state $\qket{\Psi} \in \Hspace_{\Ideas}^{(2)}$ is \textbf{entangled} if it is not a product state, i.e., $\mathrm{rank}(C) > 1$ where $C$ is the coefficient matrix.
\end{definition}

\begin{proposition}[Entanglement Criterion for Pure States]
\label{prop:entanglement-criterion}
For a pure state $\qket{\Psi}$ with coefficient matrix $C$:
\begin{enumerate}
  \item $\qket{\Psi}$ is a product state iff $\mathrm{rank}(C) = 1$
  \item $\qket{\Psi}$ is entangled iff $\mathrm{rank}(C) \geq 2$
  \item The \textbf{Schmidt rank} of $\qket{\Psi}$ equals $\mathrm{rank}(C)$
\end{enumerate}
\end{proposition}

\subsection{Schmidt Decomposition}

The Schmidt decomposition provides a canonical form for bipartite pure states, revealing the structure of entanglement.

\begin{theorem}[Schmidt Decomposition]
\label{thm:schmidt}
Every pure state $\qket{\Psi} \in \Hspace_{\Ideas}^{(2)}$ can be written uniquely (up to phase) as:
\[
\qket{\Psi} = \sum_{k=1}^{r} \lambda_k \qket{e_k} \otimes \qket{f_k}
\]
where:
\begin{itemize}
  \item $r = \mathrm{rank}(C) \leq 40$ is the \textbf{Schmidt rank}
  \item $\lambda_1 \geq \lambda_2 \geq \cdots \geq \lambda_r > 0$ are the \textbf{Schmidt coefficients}
  \item $\sum_{k=1}^{r} \lambda_k^2 = 1$ (normalization)
  \item $\{\qket{e_k}\}$ is an orthonormal set in $\Hspace_{\Ideas}$ (first subsystem)
  \item $\{\qket{f_k}\}$ is an orthonormal set in $\Hspace_{\Ideas}$ (second subsystem)
\end{itemize}
\end{theorem}

\begin{proof}
Apply singular value decomposition (SVD) to the coefficient matrix $C$: $C = U \Sigma V^\dagger$ where $U, V$ are unitary and $\Sigma = \mathrm{diag}(\lambda_1, \ldots, \lambda_r, 0, \ldots, 0)$. Define $\qket{e_k} = \sum_I U_{Ik} \qket{I}$ and $\qket{f_k} = \sum_J V_{Jk}^* \qket{J}$. Then:
\[
\qket{\Psi} = \sum_{I,J} c_{IJ} \qket{I,J} = \sum_{I,J} \sum_k U_{Ik} \lambda_k V_{Jk}^* \qket{I,J} = \sum_k \lambda_k \qket{e_k} \otimes \qket{f_k}
\]
Orthonormality of $\{\qket{e_k}\}$ and $\{\qket{f_k}\}$ follows from unitarity of $U$ and $V$.
\end{proof}

\begin{corollary}[Schmidt Decomposition Properties]
\label{cor:schmidt-properties}
For a state with Schmidt decomposition $\qket{\Psi} = \sum_k \lambda_k \qket{e_k, f_k}$:
\begin{enumerate}
  \item \textbf{Product state}: $r = 1$ (single term)
  \item \textbf{Entangled}: $r \geq 2$
  \item \textbf{Maximally entangled}: $r = 40$ and $\lambda_k = \frac{1}{\sqrt{40}}$ for all $k$
  \item \textbf{Reduced density matrices}:
  \begin{align}
  \rho_\mathsf{A} &= \Trace_\mathsf{B}(\qket{\Psi}\qbra{\Psi}) = \sum_k \lambda_k^2 \qket{e_k}\qbra{e_k} \\
  \rho_\mathsf{B} &= \Trace_\mathsf{A}(\qket{\Psi}\qbra{\Psi}) = \sum_k \lambda_k^2 \qket{f_k}\qbra{f_k}
  \end{align}
  \item \textbf{Eigenvalue matching}: $\rho_\mathsf{A}$ and $\rho_\mathsf{B}$ have identical nonzero eigenvalues $\{\lambda_k^2\}$
\end{enumerate}
\end{corollary}

\begin{definition}[Partial Trace]
\label{def:partial-trace}
The \textbf{partial trace} over subsystem $\mathsf{B}$ is the linear map $\Trace_\mathsf{B} : \BoundedOp(\Hspace_{\Ideas}^{(2)}) \to \BoundedOp(\Hspace_{\Ideas})$ defined on basis elements by:
\[
\Trace_\mathsf{B}(\qket{I, J}\qbra{I', J'}) = \delta_{JJ'} \qket{I}\qbra{I'}
\]
Similarly, $\Trace_\mathsf{A}(\qket{I, J}\qbra{I', J'}) = \delta_{II'} \qket{J}\qbra{J'}$.
\end{definition}

\begin{proposition}[Partial Trace Properties]
\label{prop:partial-trace}
For any bipartite operator $\rho \in \BoundedOp(\Hspace_{\Ideas}^{(2)})$:
\begin{enumerate}
  \item $\Trace(\Trace_\mathsf{B}(\rho)) = \Trace(\rho)$
  \item If $\rho \geq 0$, then $\Trace_\mathsf{B}(\rho) \geq 0$
  \item If $\rho$ is a density operator, so is $\Trace_\mathsf{B}(\rho)$
  \item $\Trace_\mathsf{B}(\rho_\mathsf{A} \otimes \rho_\mathsf{B}) = \rho_\mathsf{A} \cdot \Trace(\rho_\mathsf{B})$
\end{enumerate}
\end{proposition}

\subsection{Separable and Entangled Mixed States}

\begin{definition}[Separable Mixed State]
\label{def:separable-mixed}
A bipartite density operator $\rho \in \mathcal{D}(\Hspace_{\Ideas}^{(2)})$ is \textbf{separable} if it can be written as a convex combination of product states:
\[
\rho_{\mathrm{sep}} = \sum_i p_i \, \rho_\mathsf{A}^{(i)} \otimes \rho_\mathsf{B}^{(i)}
\]
where $p_i \geq 0$, $\sum_i p_i = 1$, and each $\rho_\mathsf{A}^{(i)}, \rho_\mathsf{B}^{(i)}$ is a density operator on the respective subsystem.
\end{definition}

\begin{definition}[Entangled Mixed State]
\label{def:entangled-mixed}
A bipartite density operator $\rho$ is \textbf{entangled} if it is not separable.
\end{definition}

\begin{remark}[Separability is Hard to Determine]
\label{rem:separability-hard}
Unlike pure states, determining whether a mixed state is separable is computationally difficult (NP-hard in general). We will develop operational criteria in Section~\ref{sec:entanglement-measures}.
\end{remark}

\begin{definition}[World-World Subspace]
\label{def:world-world-subspace}
For Worlds $W, W' \in \{A, B, C, D\}$, the \textbf{$(W, W')$-subspace} is:
\[
\Hspace_{WW'} = \Hspace_W \otimes \Hspace_{W'} \subset \Hspace_{\Ideas}^{(2)}
\]
with dimension $10 \times 10 = 100$. The full bipartite space decomposes as:
\[
\Hspace_{\Ideas}^{(2)} = \bigoplus_{W, W' \in \{A,B,C,D\}} \Hspace_{WW'}
\]
into 16 World-World sectors.
\end{definition}

\begin{proposition}[World-Localized Entanglement]
\label{prop:world-localized}
A state $\qket{\Psi}$ supported on $\Hspace_{WW'}$ (i.e., both Ideas in fixed Worlds $W$ and $W'$) has Schmidt rank at most 10. The maximum entanglement within a single World-World sector is:
\[
S_{\max}^{(WW')} = \log(10) \approx 2.30
\]
\end{proposition}

%% ----------------------------------------------------------------------------
\section{Entanglement Measures}
\label{sec:entanglement-measures}

This section develops quantitative measures of entanglement for TKS, extending v7.2 Definition~\ref{def:entanglement-entropy}.

\subsection{Entanglement Entropy}

\begin{definition}[Entanglement Entropy (Pure States)]
\label{def:entanglement-entropy-v73}
For a pure bipartite state $\qket{\Psi} \in \Hspace_{\Ideas}^{(2)}$, the \textbf{entanglement entropy} is:
\[
E(\qket{\Psi}) = S(\rho_\mathsf{A}) = S(\rho_\mathsf{B}) = -\sum_{k=1}^{r} \lambda_k^2 \log(\lambda_k^2)
\]
where $\{\lambda_k\}$ are the Schmidt coefficients and $S(\rho) = -\Trace(\rho \log \rho)$ is the von Neumann entropy.
\end{definition}

\begin{proposition}[Entanglement Entropy Properties]
\label{prop:entanglement-entropy-props}
The entanglement entropy satisfies:
\begin{enumerate}
  \item \textbf{Non-negativity}: $E(\qket{\Psi}) \geq 0$
  \item \textbf{Product states}: $E(\qket{\psi} \otimes \qket{\phi}) = 0$
  \item \textbf{Entangled states}: $E(\qket{\Psi}) > 0$ iff $\qket{\Psi}$ is entangled
  \item \textbf{Upper bound}: $E(\qket{\Psi}) \leq \log(\min(\dim \Hspace_\mathsf{A}, \dim \Hspace_\mathsf{B})) = \log(40)$
  \item \textbf{Invariance}: $E((U_\mathsf{A} \otimes U_\mathsf{B})\qket{\Psi}) = E(\qket{\Psi})$ for unitaries $U_\mathsf{A}, U_\mathsf{B}$
\end{enumerate}
\end{proposition}

\begin{theorem}[Maximum Entanglement in TKS]
\label{thm:max-entanglement-v73}
The maximum entanglement entropy for TKS bipartite states is:
\[
E_{\max} = \log(40) \approx 3.69 \text{ (in natural log units)}
\]
achieved by the maximally entangled state:
\[
\qket{\Phi^+} = \frac{1}{\sqrt{40}} \sum_{I \in \Ideas} \qket{I, I}
\]
\end{theorem}

\begin{proof}
By Corollary~\ref{cor:schmidt-properties}, maximum entanglement entropy occurs when all 40 Schmidt coefficients equal $\frac{1}{\sqrt{40}}$. Then:
\[
E = -\sum_{k=1}^{40} \frac{1}{40} \log\left(\frac{1}{40}\right) = \log(40)
\]
The state $\qket{\Phi^+}$ achieves this: its coefficient matrix is $C = \frac{1}{\sqrt{40}} \mathbf{I}_{40}$, which has 40 singular values all equal to $\frac{1}{\sqrt{40}}$.
\end{proof}

\subsection{Entanglement of Formation}

For mixed states, entanglement entropy does not directly apply. We define the entanglement of formation.

\begin{definition}[Entanglement of Formation]
\label{def:eof}
For a mixed state $\rho \in \mathcal{D}(\Hspace_{\Ideas}^{(2)})$, the \textbf{entanglement of formation} is:
\[
E_F(\rho) = \min_{\{p_i, \qket{\psi_i}\}} \sum_i p_i E(\qket{\psi_i})
\]
where the minimum is over all pure-state decompositions $\rho = \sum_i p_i \qket{\psi_i}\qbra{\psi_i}$.
\end{definition}

\begin{proposition}[Entanglement of Formation Properties]
\label{prop:eof-props}
The entanglement of formation satisfies:
\begin{enumerate}
  \item $E_F(\rho) \geq 0$
  \item $E_F(\rho) = 0$ iff $\rho$ is separable
  \item $E_F(\qket{\psi}\qbra{\psi}) = E(\qket{\psi})$ for pure states
  \item $E_F$ is convex: $E_F(\sum_i p_i \rho_i) \leq \sum_i p_i E_F(\rho_i)$
  \item $E_F$ is monotonic under LOCC (local operations and classical communication)
\end{enumerate}
\end{proposition}

\subsection{Negativity and PPT Criterion}

\begin{definition}[Partial Transpose]
\label{def:partial-transpose}
The \textbf{partial transpose} with respect to subsystem $\mathsf{B}$ is defined on basis elements by:
\[
(\qket{I, J}\qbra{I', J'})^{T_\mathsf{B}} = \qket{I, J'}\qbra{I', J}
\]
Extended linearly to all operators. The partial transpose $\rho^{T_\mathsf{B}}$ swaps indices in the second subsystem.
\end{definition}

\begin{theorem}[Peres-Horodecki (PPT) Criterion]
\label{thm:ppt}
If $\rho$ is separable, then $\rho^{T_\mathsf{B}} \geq 0$ (positive partial transpose, PPT).

Equivalently: if $\rho^{T_\mathsf{B}}$ has any negative eigenvalue, then $\rho$ is entangled.
\end{theorem}

\begin{proof}
For a product state $\rho = \rho_\mathsf{A} \otimes \rho_\mathsf{B}$: $\rho^{T_\mathsf{B}} = \rho_\mathsf{A} \otimes \rho_\mathsf{B}^T$. Since $\rho_\mathsf{B}^T$ has the same eigenvalues as $\rho_\mathsf{B} \geq 0$, we have $\rho^{T_\mathsf{B}} \geq 0$. By convexity, all separable states satisfy PPT.
\end{proof}

\begin{remark}[PPT is Necessary but Not Sufficient]
\label{rem:ppt-not-sufficient}
In dimensions $2 \times 2$ and $2 \times 3$, PPT is also sufficient for separability. In $\Hspace_{\Ideas}^{(2)}$ ($40 \times 40$), there exist ``PPT entangled states'' (also called bound entangled states) that are entangled but have positive partial transpose.
\end{remark}

\begin{definition}[Negativity]
\label{def:negativity}
The \textbf{negativity} of a bipartite state $\rho$ is:
\[
\mathcal{N}(\rho) = \frac{\|\rho^{T_\mathsf{B}}\|_1 - 1}{2}
\]
where $\|A\|_1 = \Trace(\sqrt{A^\dagger A})$ is the trace norm. Equivalently:
\[
\mathcal{N}(\rho) = \sum_{\lambda_i < 0} |\lambda_i|
\]
where $\{\lambda_i\}$ are eigenvalues of $\rho^{T_\mathsf{B}}$.
\end{definition}

\begin{definition}[Logarithmic Negativity]
\label{def:log-negativity}
The \textbf{logarithmic negativity} is:
\[
E_\mathcal{N}(\rho) = \log(\|\rho^{T_\mathsf{B}}\|_1) = \log(1 + 2\mathcal{N}(\rho))
\]
This is an entanglement monotone (non-increasing under LOCC).
\end{definition}

\begin{proposition}[Negativity Properties]
\label{prop:negativity-props}
The negativity satisfies:
\begin{enumerate}
  \item $\mathcal{N}(\rho) \geq 0$
  \item $\mathcal{N}(\rho) = 0$ for all PPT states (including all separable states)
  \item $\mathcal{N}(\rho) > 0$ implies $\rho$ is entangled
  \item For a pure state $\qket{\Psi}$ with Schmidt coefficients $\{\lambda_k\}$:
  \[
  \mathcal{N}(\qket{\Psi}\qbra{\Psi}) = \frac{1}{2}\left[\left(\sum_k \lambda_k\right)^2 - 1\right]
  \]
\end{enumerate}
\end{proposition}

\subsection{Concurrence}

For the analogy with two-qubit systems, we adapt the concurrence measure to TKS.

\begin{definition}[Concurrence for Pure States]
\label{def:concurrence-pure}
For a pure state $\qket{\Psi}$ with Schmidt coefficients $\{\lambda_k\}_{k=1}^{r}$, the \textbf{concurrence} is:
\[
C(\qket{\Psi}) = \sqrt{2(1 - \Trace(\rho_\mathsf{A}^2))} = \sqrt{2\left(1 - \sum_k \lambda_k^4\right)}
\]
\end{definition}

\begin{proposition}[Concurrence Properties]
\label{prop:concurrence-props}
The concurrence satisfies:
\begin{enumerate}
  \item $0 \leq C(\qket{\Psi}) \leq \sqrt{2(1 - 1/40)} = \sqrt{78/40} \approx 1.396$
  \item $C(\qket{\Psi}) = 0$ iff $\qket{\Psi}$ is a product state
  \item Maximum concurrence is achieved by maximally entangled states
  \item \textbf{Relation to purity}: $C^2 = 2(1 - \mathrm{Pur}(\rho_\mathsf{A}))$ where $\mathrm{Pur}(\rho) = \Trace(\rho^2)$
\end{enumerate}
\end{proposition}

\begin{definition}[Concurrence for Two-Element Subspaces]
\label{def:concurrence-two-element}
When restricting to a two-dimensional subspace of each subsystem (e.g., two Ideas $I_1, I_2$ on each side), the standard two-qubit concurrence applies. For $\qket{\Psi} = a\qket{I_1, I_1} + b\qket{I_1, I_2} + c\qket{I_2, I_1} + d\qket{I_2, I_2}$:
\[
C_{\{I_1, I_2\}}(\qket{\Psi}) = 2|ad - bc|
\]
This is the Wootters concurrence for the two-qubit subsystem.
\end{definition}

%% ----------------------------------------------------------------------------
\section{Noetic Bell States and Maximal Entanglement}
\label{sec:maximally-entangled}

This section develops the TKS analogues of Bell states---maximally entangled states that serve as fundamental resources for quantum protocols.

\subsection{Bell Basis for TKS}

\begin{definition}[Noetic Bell States]
\label{def:noetic-bell}
The \textbf{noetic Bell states} are the maximally entangled basis states for $\Hspace_{\Ideas}^{(2)}$. The primary Bell state is:
\[
\qket{\Phi^+} = \frac{1}{\sqrt{40}} \sum_{I \in \Ideas} \qket{I, I} = \frac{1}{\sqrt{40}} \sum_{W \in \{A,B,C,D\}} \sum_{n=1}^{10} \qket{Wn, Wn}
\]
The \textbf{generalized Bell basis} consists of $40^2 = 1600$ states:
\[
\qket{\Phi_{IJ}} = (\mathbf{I} \otimes X_J Z_I) \qket{\Phi^+}
\]
where $X_J$ and $Z_I$ are generalized Pauli operators on $\Hspace_{\Ideas}$ (see below).
\end{definition}

\begin{definition}[Generalized Pauli Operators]
\label{def:generalized-pauli}
The \textbf{generalized Pauli operators} (Weyl operators) on $\Hspace_{\Ideas}$ are:
\begin{align}
X_J &= \sum_{I \in \Ideas} \qket{I \oplus J}\qbra{I} \quad \text{(shift by $J$)} \\
Z_I &= \sum_{K \in \Ideas} \omega^{\langle I, K \rangle} \qket{K}\qbra{K} \quad \text{(phase)}
\end{align}
where $\oplus$ is addition in $\mathbb{Z}_{40}$ (identifying $\Ideas \cong \mathbb{Z}_{40}$), $\omega = e^{2\pi i / 40}$, and $\langle I, K \rangle$ is the inner product in $\mathbb{Z}_{40}$.
\end{definition}

\begin{proposition}[Bell Basis Properties]
\label{prop:bell-basis}
The generalized Bell states satisfy:
\begin{enumerate}
  \item \textbf{Orthonormality}: $\qbra{\Phi_{IJ}}\Phi_{I'J'}\rangle = \delta_{II'} \delta_{JJ'}$
  \item \textbf{Completeness}: $\sum_{I, J} \qket{\Phi_{IJ}}\qbra{\Phi_{IJ}} = \mathbf{I}^{(2)}$
  \item \textbf{Maximal entanglement}: Each $\qket{\Phi_{IJ}}$ has entanglement entropy $\log(40)$
  \item \textbf{Local unitary equivalence}: All $\qket{\Phi_{IJ}}$ are related by local unitaries
\end{enumerate}
\end{proposition}

\subsection{World-Restricted Bell States}

\begin{definition}[World Bell States]
\label{def:world-bell}
For Worlds $W, W' \in \{A, B, C, D\}$, the \textbf{$(W, W')$-Bell state} is:
\[
\qket{\Phi^+_{WW'}} = \frac{1}{\sqrt{10}} \sum_{n=1}^{10} \qket{Wn, W'n}
\]
This has support only in the $(W, W')$-sector and entanglement entropy $\log(10)$.
\end{definition}

\begin{example}[World-Entangled Bell States]
\label{ex:world-bell-states}
The 16 World Bell states include:
\begin{align}
\qket{\Phi^+_{AA}} &= \frac{1}{\sqrt{10}} \sum_{n=1}^{10} \qket{An, An} & \text{(intra-World A)} \\
\qket{\Phi^+_{AD}} &= \frac{1}{\sqrt{10}} \sum_{n=1}^{10} \qket{An, Dn} & \text{(Spiritual-Physical)} \\
\qket{\Phi^+_{BC}} &= \frac{1}{\sqrt{10}} \sum_{n=1}^{10} \qket{Bn, Cn} & \text{(Psychic-Phenomenal)}
\end{align}
The $\qket{\Phi^+_{AD}}$ state represents maximal correlation between Spiritual (A) and Physical (D) Ideas---the quantum analogue of the ACBE descent.
\end{example}

\begin{proposition}[Decomposition of Maximal Bell State]
\label{prop:bell-decomposition}
The maximally entangled state decomposes as:
\[
\qket{\Phi^+} = \frac{1}{2} \sum_{W \in \{A,B,C,D\}} \qket{\Phi^+_{WW}}
\]
This shows $\qket{\Phi^+}$ as an equal superposition over intra-World correlations.
\end{proposition}

\subsection{Bell Inequalities in TKS}

\begin{definition}[Noetic CHSH Observable]
\label{def:chsh-observable}
For dichotomic observables $\qop{A}_1, \qop{A}_2$ on subsystem $\mathsf{A}$ and $\qop{B}_1, \qop{B}_2$ on subsystem $\mathsf{B}$ (each with eigenvalues $\pm 1$), the \textbf{CHSH operator} is:
\[
\qop{C}_{\mathrm{CHSH}} = \qop{A}_1 \otimes \qop{B}_1 + \qop{A}_1 \otimes \qop{B}_2 + \qop{A}_2 \otimes \qop{B}_1 - \qop{A}_2 \otimes \qop{B}_2
\]
\end{definition}

\begin{theorem}[CHSH Inequality]
\label{thm:chsh}
For any local hidden variable model:
\[
|\langle \qop{C}_{\mathrm{CHSH}} \rangle_{\mathrm{LHV}}| \leq 2
\]
Quantum mechanics (and TKS) allows:
\[
|\langle \qop{C}_{\mathrm{CHSH}} \rangle_{\mathrm{QM}}| \leq 2\sqrt{2} \approx 2.83
\]
with the maximum achieved by maximally entangled states.
\end{theorem}

\begin{definition}[Noetic Bell Observables]
\label{def:noetic-bell-observables}
For TKS, natural dichotomic observables include:
\begin{itemize}
  \item \textbf{World parity}: $\qop{O}_W = \Projector_{A \cup B} - \Projector_{C \cup D}$ (Spiritual+Psychic vs Phenomenal+Physical)
  \item \textbf{Element parity}: $\qop{O}_{\mathrm{odd}} = \sum_{n \text{ odd}} \Projector_n - \sum_{n \text{ even}} \Projector_n$
  \item \textbf{Noetic observables}: $\qop{O}_k = 2\qop{\nu}_k^\dagger \qop{\nu}_k - \mathbf{I}$ for noetic $k$
\end{itemize}
\end{definition}

\begin{theorem}[Bell Violation in TKS]
\label{thm:bell-violation-tks}
The noetic Bell state $\qket{\Phi^+}$ violates the CHSH inequality with appropriate noetic observables. Specifically, for observables constructed from noetic operators corresponding to non-commuting noetics:
\[
\qbra{\Phi^+}\qop{C}_{\mathrm{CHSH}}\qket{\Phi^+} = 2\sqrt{2}
\]
achieving the Tsirelson bound.
\end{theorem}

\begin{proof}
(Sketch) The maximally entangled state $\qket{\Phi^+}$ satisfies $(\qop{A} \otimes \mathbf{I})\qket{\Phi^+} = (\mathbf{I} \otimes \qop{A}^T)\qket{\Phi^+}$ for any operator $\qop{A}$. Choosing $\qop{A}_1, \qop{A}_2$ such that they anti-commute and similarly for $\qop{B}_1, \qop{B}_2$, the expectation value achieves $2\sqrt{2}$. The noetic structure provides natural non-commuting operators via the non-abelian noetic algebra.
\end{proof}

\begin{corollary}[Nonlocal Noetic Correlations]
\label{cor:nonlocal-correlations}
Entangled noetic states exhibit correlations that cannot be explained by classical (local realistic) TKS models. This has implications for:
\begin{enumerate}
  \item \textbf{Quantum RPM}: Correlated acquisition cannot be simulated classically
  \item \textbf{Noetic communication}: Entanglement enables superdense coding of Idea information
  \item \textbf{Multi-agent TKS}: Shared entangled Ideas between agents produce nonlocal correlations
\end{enumerate}
\end{corollary}

%% ----------------------------------------------------------------------------
\section{World-World Entanglement}
\label{sec:entanglement-acquisition}

This section studies entanglement between different Worlds, connecting to the ACBE descent and classical TKS correlations.

\subsection{Inter-World Entanglement Structure}

\begin{definition}[Inter-World Entanglement]
\label{def:inter-world-entanglement}
For a bipartite state $\rho$ on $\Hspace_{\Ideas}^{(2)}$, the \textbf{$(W, W')$-entanglement} is:
\[
E_{WW'}(\rho) = E_F(\Pi_{WW'} \rho \Pi_{WW'} / p_{WW'})
\]
where $\Pi_{WW'} = \Projector_W \otimes \Projector_{W'}$ projects onto the $(W, W')$-sector and $p_{WW'} = \Trace(\Pi_{WW'} \rho)$ is the probability of that sector.
\end{definition}

\begin{proposition}[World Entanglement Decomposition]
\label{prop:world-entanglement-decomp}
For a state $\rho$ on $\Hspace_{\Ideas}^{(2)}$:
\[
E_F(\rho) \geq \sum_{W, W'} p_{WW'} E_{WW'}(\rho)
\]
with equality when $\rho$ has no coherences between different World-World sectors.
\end{proposition}

\begin{definition}[ACBE-Correlated State]
\label{def:acbe-correlated}
An \textbf{ACBE-correlated state} has support only on the ACBE descent pairs:
\[
\rho_{\mathrm{ACBE}} \in \mathcal{D}(\Hspace_{AD} \oplus \Hspace_{BC} \oplus \Hspace_{CB} \oplus \Hspace_{DA})
\]
representing correlations that respect the Spiritual $\to$ Physical descent structure.
\end{definition}

\subsection{ACBE and Entanglement Dynamics}

\begin{definition}[Local ACBE Channel]
\label{def:local-acbe}
The \textbf{local ACBE channel} on the first subsystem is:
\[
(\mathcal{E}_{\mathrm{ACBE}} \otimes \mathbf{I})(\rho) = \sum_{n=1}^{10} (K_n \otimes \mathbf{I}) \rho (K_n^\dagger \otimes \mathbf{I})
\]
where $K_n = \qket{Dn}\qbra{An}$ are the ACBE Kraus operators.
\end{definition}

\begin{theorem}[ACBE Entanglement Transformation]
\label{thm:acbe-entanglement}
\textbf{(Main Theorem: Entanglement under Noetic Operations)}

Let $\qket{\Psi} \in \Hspace_{\Ideas}^{(2)}$ be a bipartite pure state and $\mathcal{E}_{\mathrm{ACBE}}$ the ACBE channel applied to subsystem $\mathsf{A}$. Then:

\begin{enumerate}
  \item \textbf{World shift}: $(\mathcal{E}_{\mathrm{ACBE}} \otimes \mathbf{I})$ maps $(A, W')$-sector to $(D, W')$-sector

  \item \textbf{Entanglement preservation within A}: For $\qket{\Psi}$ supported on $\Hspace_{AW'}$:
  \[
  E((\mathcal{E}_{\mathrm{ACBE}} \otimes \mathbf{I})(\qket{\Psi}\qbra{\Psi})) = E(\qket{\Psi})
  \]
  The ACBE channel preserves entanglement entropy when acting locally.

  \item \textbf{Bell state transformation}:
  \[
  (\mathcal{E}_{\mathrm{ACBE}} \otimes \mathbf{I})(\qket{\Phi^+_{AW'}}\qbra{\Phi^+_{AW'}}) = \qket{\Phi^+_{DW'}}\qbra{\Phi^+_{DW'}}
  \]

  \item \textbf{Cross-World coherence}: Coherences between $(A, W')$ and $(A, W'')$ sectors are preserved and shifted to $(D, W')$ and $(D, W'')$.

  \item \textbf{Annihilation}: States with no support on World A (in subsystem $\mathsf{A}$) are annihilated.
\end{enumerate}
\end{theorem}

\begin{proof}
\textbf{(1)}: The Kraus operators $K_n = \qket{Dn}\qbra{An}$ map $\qket{An} \to \qket{Dn}$.

\textbf{(2)}: For $\qket{\Psi} = \sum_{n,m} c_{nm} \qket{An, W'm}$:
\begin{align}
(\mathcal{E}_{\mathrm{ACBE}} \otimes \mathbf{I})(\qket{\Psi}\qbra{\Psi}) &= \sum_{n} (K_n \otimes \mathbf{I}) \qket{\Psi}\qbra{\Psi} (K_n^\dagger \otimes \mathbf{I}) \\
&= \sum_{n,m,m'} c_{nm} c_{nm'}^* \qket{Dn, W'm}\qbra{Dn, W'm'}
\end{align}
This is pure (rank 1) with the same coefficient structure, hence the same entanglement.

\textbf{(3)}: Direct calculation: $(\mathcal{E}_{\mathrm{ACBE}} \otimes \mathbf{I})(\frac{1}{10}\sum_n \qket{An, W'n}\qbra{An, W'n}) = \frac{1}{10}\sum_n \qket{Dn, W'n}\qbra{Dn, W'n}$.

\textbf{(4)}: Coherences $\qket{An, W'm}\qbra{An, W''m'}$ map to $\qket{Dn, W'm}\qbra{Dn, W''m'}$.

\textbf{(5)}: $K_n \qket{Wn'} = 0$ for $W \neq A$.
\end{proof}

\begin{corollary}[ACBE Creates No Entanglement]
\label{cor:acbe-no-create}
The local ACBE channel $\mathcal{E}_{\mathrm{ACBE}} \otimes \mathbf{I}$ cannot create entanglement from a product state:
\[
\qket{\psi} \otimes \qket{\phi} \xmapsto{\mathcal{E}_{\mathrm{ACBE}} \otimes \mathbf{I}} \mathcal{E}_{\mathrm{ACBE}}(\qket{\psi}\qbra{\psi}) \otimes \qket{\phi}\qbra{\phi}
\]
which is separable. This is a general property of local operations.
\end{corollary}

\begin{corollary}[Entanglement-ACBE Commutativity]
\label{cor:entanglement-acbe-commute}
For the A-D Bell state: applying ACBE to both subsystems produces a D-D Bell state:
\[
(\mathcal{E}_{\mathrm{ACBE}} \otimes \mathcal{E}_{\mathrm{ACBE}})(\qket{\Phi^+_{AA}}\qbra{\Phi^+_{AA}}) = \qket{\Phi^+_{DD}}\qbra{\Phi^+_{DD}}
\]
The full ACBE descent preserves maximal entanglement within each World-World sector.
\end{corollary}

\subsection{Monogamy of Entanglement}

\begin{theorem}[Monogamy of Noetic Entanglement]
\label{thm:monogamy}
For a tripartite pure state $\qket{\Psi_{ABC}} \in \Hspace_{\Ideas}^{(3)} = \Hspace_{\Ideas}^{\otimes 3}$, the squared concurrences satisfy:
\[
C^2(\rho_{AB}) + C^2(\rho_{AC}) \leq C^2(\rho_{A(BC)})
\]
where $\rho_{AB} = \Trace_C(\qket{\Psi}\qbra{\Psi})$, etc., and $\rho_{A(BC)}$ is the state with $BC$ treated as a single system.

In TKS terms: if Idea A is maximally entangled with Idea B, it cannot be entangled with Idea C.
\end{theorem}

\begin{corollary}[Monogamy Constraints on Noetic Networks]
\label{cor:monogamy-networks}
In a network of $n$ Ideas with pairwise entanglement $E_{ij}$, the monogamy inequality constrains:
\[
\sum_{j \neq i} E_{ij} \leq E_{i,\mathrm{rest}}
\]
This limits how entanglement can be distributed across multiple Ideas simultaneously.
\end{corollary}

\begin{definition}[Entanglement of Assistance]
\label{def:eoa}
The \textbf{entanglement of assistance} for a mixed state $\rho_{AB}$ is:
\[
E_A(\rho_{AB}) = \max_{\{p_i, \qket{\psi_i}\}} \sum_i p_i E(\qket{\psi_i})
\]
the maximum average entanglement over pure-state decompositions.
\end{definition}

\begin{proposition}[Monogamy via Entanglement of Assistance]
\label{prop:monogamy-eoa}
For tripartite $\qket{\Psi_{ABC}}$:
\[
E_F(\rho_{AB}) + E_F(\rho_{AC}) \leq E(\qket{\Psi_{A(BC)}})
\]
where $E_F$ is entanglement of formation.
\end{proposition}

%% ----------------------------------------------------------------------------
\section{Multipartite Entanglement}
\label{sec:multipartite-entanglement}

This section extends to systems of three or more Ideas.

\subsection{Tripartite Hilbert Space}

\begin{definition}[Tripartite Noetic Space]
\label{def:tripartite}
The \textbf{tripartite noetic Hilbert space} is:
\[
\Hspace_{\Ideas}^{(3)} = \Hspace_{\Ideas} \otimes \Hspace_{\Ideas} \otimes \Hspace_{\Ideas}
\]
with dimension $40^3 = 64000$ and basis $\{\qket{I, J, K}\}_{I, J, K \in \Ideas}$.
\end{definition}

\begin{definition}[Fully Separable State]
\label{def:fully-separable}
A tripartite state $\rho$ is \textbf{fully separable} if:
\[
\rho = \sum_i p_i \, \rho_\mathsf{A}^{(i)} \otimes \rho_\mathsf{B}^{(i)} \otimes \rho_\mathsf{C}^{(i)}
\]
\end{definition}

\begin{definition}[Biseparable State]
\label{def:biseparable}
A tripartite state is \textbf{biseparable} if it is separable with respect to some bipartition, e.g.:
\[
\rho = \sum_i p_i \, \rho_\mathsf{A}^{(i)} \otimes \rho_{\mathsf{BC}}^{(i)}
\]
where $\rho_{\mathsf{BC}}^{(i)}$ may be entangled.
\end{definition}

\begin{definition}[Genuinely Multipartite Entangled]
\label{def:gme}
A state is \textbf{genuinely multipartite entangled (GME)} if it is not biseparable with respect to any bipartition.
\end{definition}

\subsection{GHZ and W States}

\begin{definition}[Noetic GHZ State]
\label{def:ghz}
The \textbf{noetic GHZ (Greenberger-Horne-Zeilinger) state} is:
\[
\qket{\mathrm{GHZ}} = \frac{1}{\sqrt{40}} \sum_{I \in \Ideas} \qket{I, I, I}
\]
This exhibits maximal correlation: measuring any two subsystems in the Idea basis yields perfectly correlated outcomes.
\end{definition}

\begin{proposition}[GHZ Properties]
\label{prop:ghz-props}
The noetic GHZ state satisfies:
\begin{enumerate}
  \item \textbf{GME}: $\qket{\mathrm{GHZ}}$ is genuinely tripartite entangled
  \item \textbf{Reduced states}: $\rho_\mathsf{A} = \rho_\mathsf{B} = \rho_\mathsf{C} = \frac{1}{40}\mathbf{I}$ (maximally mixed)
  \item \textbf{Bipartite entanglement}: $E(\rho_{AB}) = 0$ (no pairwise entanglement after tracing out C)
  \item \textbf{Fragility}: Losing any one subsystem destroys all entanglement
\end{enumerate}
\end{proposition}

\begin{definition}[Noetic W State]
\label{def:w-state}
The \textbf{noetic W state} for a fixed Idea $I_0$ is:
\[
\qket{\mathrm{W}_{I_0}} = \frac{1}{\sqrt{3}} \left( \qket{I_0, I_1, I_1} + \qket{I_1, I_0, I_1} + \qket{I_1, I_1, I_0} \right)
\]
where $I_1 \neq I_0$ is another fixed Idea. More generally, for the full W state:
\[
\qket{\mathrm{W}} = \frac{1}{\sqrt{3 \cdot 40 \cdot 39}} \sum_{I \neq J} \left( \qket{I, J, J} + \qket{J, I, J} + \qket{J, J, I} \right)
\]
\end{definition}

\begin{proposition}[W State Properties]
\label{prop:w-props}
The noetic W state satisfies:
\begin{enumerate}
  \item \textbf{GME}: $\qket{\mathrm{W}}$ is genuinely tripartite entangled
  \item \textbf{Robustness}: Tracing out any subsystem leaves an entangled state
  \item \textbf{Pairwise entanglement}: $E(\rho_{AB}) > 0$
  \item \textbf{Different entanglement class}: $\qket{\mathrm{W}}$ and $\qket{\mathrm{GHZ}}$ cannot be transformed into each other by LOCC
\end{enumerate}
\end{proposition}

\begin{theorem}[GHZ vs W Classification]
\label{thm:ghz-w-classes}
Genuinely tripartite entangled pure states in TKS fall into two inequivalent SLOCC (stochastic LOCC) classes:
\begin{enumerate}
  \item \textbf{GHZ class}: States locally equivalent to $\qket{\mathrm{GHZ}}$
  \item \textbf{W class}: States locally equivalent to $\qket{\mathrm{W}}$
\end{enumerate}
Plus the separable and biseparable classes.
\end{theorem}

\subsection{Noetic Graph States}

\begin{definition}[Noetic Graph State]
\label{def:graph-state}
For a graph $G = (V, E)$ with $|V| = n$ vertices (each corresponding to an Idea), the \textbf{noetic graph state} is:
\[
\qket{G} = \prod_{(i,j) \in E} \mathrm{CZ}_{ij} \qket{+}^{\otimes n}
\]
where $\qket{+} = \frac{1}{\sqrt{40}}\sum_I \qket{I}$ and $\mathrm{CZ}_{ij}$ is the controlled-Z gate:
\[
\mathrm{CZ}_{ij} = \sum_{I, J} \omega^{\langle I, J \rangle} \qket{I, J}\qbra{I, J}
\]
with $\omega = e^{2\pi i / 40}$.
\end{definition}

\begin{example}[Linear Cluster State]
\label{ex:cluster}
For a linear graph with $n$ Ideas, the cluster state enables measurement-based quantum computation on noetic states.
\end{example}

%% ----------------------------------------------------------------------------
%% HANDOFF NOTE
%% ----------------------------------------------------------------------------
\begin{tcolorbox}[colback=green!5!white,colframe=green!50!black,title=\textbf{Handoff: Next Steps},breakable]
This completes Chapter 5 on Noetic Entanglement. The subsequent chapters should:
\begin{itemize}
  \item \textbf{Chapter 6 (Measurement)}: Measurement-induced entanglement changes; quantum state tomography for TKS; measurement in entangled basis
  \item \textbf{Chapter 7 (Quantum Fractals)}: Entanglement generation by fractal operators; entanglement entropy of fractal attractors
\end{itemize}

Key results established in this chapter:
\begin{enumerate}
  \item \textbf{Bipartite structure}: $\Hspace_{\Ideas}^{(2)} = \Hspace_{\Ideas} \otimes \Hspace_{\Ideas}$ with $\dim = 1600$ (Sec.~\ref{sec:bipartite-states})
  \item \textbf{Schmidt decomposition}: Canonical form for pure bipartite states (Thm.~\ref{thm:schmidt})
  \item \textbf{Entanglement measures}: Entropy, formation, negativity, concurrence (Sec.~\ref{sec:entanglement-measures})
  \item \textbf{Maximum entanglement}: $E_{\max} = \log(40)$ achieved by $\qket{\Phi^+}$ (Thm.~\ref{thm:max-entanglement-v73})
  \item \textbf{Noetic Bell states}: $\qket{\Phi^+} = \frac{1}{\sqrt{40}}\sum_I \qket{I,I}$ and World Bell states $\qket{\Phi^+_{WW'}}$ (Sec.~\ref{sec:maximally-entangled})
  \item \textbf{Bell inequality violation}: TKS achieves Tsirelson bound $2\sqrt{2}$ (Thm.~\ref{thm:bell-violation-tks})
  \item \textbf{ACBE Entanglement Theorem~\ref{thm:acbe-entanglement}}: Local ACBE preserves entanglement entropy, shifts World sectors
  \item \textbf{Monogamy}: Constraints on entanglement distribution (Thm.~\ref{thm:monogamy})
  \item \textbf{Multipartite}: GHZ and W states, graph states (Sec.~\ref{sec:multipartite-entanglement})
\end{enumerate}

\textbf{Connection to v7.2}: Extends Definition~\ref{def:tensor-product}, \ref{def:entangled-state}, \ref{def:entanglement-entropy} from v7.2 with full mathematical treatment. The ACBE entanglement dynamics (Theorem~\ref{thm:acbe-entanglement}) provides the quantum refinement of the classical ACBE descent.

\textbf{Connection to Classical TKS}: Bell inequality violations (Theorem~\ref{thm:bell-violation-tks}) demonstrate that entangled noetic states produce correlations impossible in classical (local realistic) TKS models. The monogamy constraints (Theorem~\ref{thm:monogamy}) limit how Ideas can be correlated across a network.
\end{tcolorbox}

%% ============================================================================
%% CHAPTER 6: QUANTUM MEASUREMENT IN TKS
%% ============================================================================
\chapter{Quantum Measurement in TKS}
\label{ch:quantum-measurement}

\begin{tcolorbox}[colback=blue!5!white,colframe=blue!75!black,title=\vnewthree]
This chapter develops measurement theory for TKS, describing how noetic states collapse upon observation.
\end{tcolorbox}

%% ----------------------------------------------------------------------------
\section{Projective Measurements}
\label{sec:projective-measurements}

This section develops the theory of projective (von Neumann) measurements on noetic states, extending v7.2 foundations (Definition~\ref{def:world-measurement}, \ref{def:element-measurement}).

\subsection{Projection-Valued Measures}

\begin{definition}[Projection-Valued Measure]
\label{def:pvm}
A \textbf{projection-valued measure (PVM)} on $\Hspace_{\Ideas}$ is a collection $\{\Projector_m\}_{m \in \mathcal{M}}$ of orthogonal projectors indexed by a measurable space $\mathcal{M}$, satisfying:
\begin{enumerate}
  \item \textbf{Orthogonality}: $\Projector_m \Projector_{m'} = \delta_{mm'} \Projector_m$
  \item \textbf{Completeness}: $\sum_{m \in \mathcal{M}} \Projector_m = \mathbf{I}$
  \item \textbf{Hermiticity}: $\Projector_m^\dagger = \Projector_m$
  \item \textbf{Idempotence}: $\Projector_m^2 = \Projector_m$
\end{enumerate}
\end{definition}

\begin{definition}[Projective Measurement]
\label{def:projective-measurement}
A \textbf{projective measurement} (or \textbf{von Neumann measurement}) associated with PVM $\{\Projector_m\}$ on a pure state $\qket{\psi}$ consists of:
\begin{enumerate}
  \item \textbf{Born rule}: Outcome $m$ occurs with probability
  \[
  p(m|\psi) = \qbra{\psi}\Projector_m\qket{\psi} = \|\Projector_m\qket{\psi}\|^2
  \]

  \item \textbf{State collapse}: Conditional on outcome $m$, the post-measurement state is
  \[
  \qket{\psi_m} = \frac{\Projector_m\qket{\psi}}{\|\Projector_m\qket{\psi}\|} = \frac{\Projector_m\qket{\psi}}{\sqrt{p(m|\psi)}}
  \]
\end{enumerate}
\end{definition}

\begin{proposition}[Born Rule Properties]
\label{prop:born-rule}
The Born rule satisfies:
\begin{enumerate}
  \item \textbf{Normalization}: $\sum_m p(m|\psi) = 1$
  \item \textbf{Positivity}: $p(m|\psi) \geq 0$
  \item \textbf{Certainty}: $p(m|\psi) = 1$ iff $\qket{\psi} \in \mathrm{Im}(\Projector_m)$
  \item \textbf{Exclusion}: $p(m|\psi) = 0$ iff $\qket{\psi} \perp \mathrm{Im}(\Projector_m)$
\end{enumerate}
\end{proposition}

\begin{proof}
\textbf{(1)}: $\sum_m p(m|\psi) = \sum_m \qbra{\psi}\Projector_m\qket{\psi} = \qbra{\psi}(\sum_m \Projector_m)\qket{\psi} = \qbra{\psi}\mathbf{I}\qket{\psi} = 1$.

\textbf{(2)}: $p(m|\psi) = \|\Projector_m\qket{\psi}\|^2 \geq 0$.

\textbf{(3)}: $p(m|\psi) = 1$ iff $\Projector_m\qket{\psi} = \qket{\psi}$ (since $\sum_{m' \neq m} p(m'|\psi) = 0$).

\textbf{(4)}: $p(m|\psi) = 0$ iff $\Projector_m\qket{\psi} = 0$ iff $\qket{\psi} \in \ker(\Projector_m) = \mathrm{Im}(\Projector_m)^\perp$.
\end{proof}

\subsection{Measurement in the Idea Basis}

\begin{definition}[Complete Idea Measurement]
\label{def:idea-measurement}
The \textbf{complete Idea measurement} uses the PVM $\{\Projector_I\}_{I \in \Ideas}$ where:
\[
\Projector_I = \qket{I}\qbra{I}
\]
For state $\qket{\psi} = \sum_{I \in \Ideas} c_I \qket{I}$:
\begin{itemize}
  \item \textbf{Probability}: $p(I|\psi) = |c_I|^2$
  \item \textbf{Post-measurement}: $\qket{\psi_I} = \qket{I}$
\end{itemize}
This is the finest-grained measurement, resolving both World and Element.
\end{definition}

\begin{definition}[World Measurement]
\label{def:world-measurement-v73}
The \textbf{World measurement} uses the coarse-grained PVM $\{\Projector_W\}_{W \in \{A,B,C,D\}}$ where:
\[
\Projector_W = \sum_{n=1}^{10} \qket{Wn}\qbra{Wn}
\]
For state $\qket{\psi} = \sum_{W,n} c_{Wn} \qket{Wn}$:
\begin{itemize}
  \item \textbf{Probability}: $p(W|\psi) = \sum_{n=1}^{10} |c_{Wn}|^2$
  \item \textbf{Post-measurement}: $\qket{\psi_W} = \frac{1}{\sqrt{p(W|\psi)}} \sum_{n=1}^{10} c_{Wn} \qket{Wn}$
\end{itemize}
This determines the World but preserves Element superposition within that World.
\end{definition}

\begin{definition}[Element-Only Measurement]
\label{def:element-only-measurement}
The \textbf{Element-only measurement} (ignoring World) uses the PVM $\{\Projector^{(n)}\}_{n=1}^{10}$ where:
\[
\Projector^{(n)} = \sum_{W \in \{A,B,C,D\}} \qket{Wn}\qbra{Wn}
\]
This determines which Element (1-10) but preserves World superposition.
\end{definition}

\begin{example}[Measuring a Superposition]
\label{ex:measuring-superposition}
Consider the equal superposition over World A:
\[
\qket{\psi} = \frac{1}{\sqrt{10}} \sum_{n=1}^{10} \qket{An}
\]
\begin{itemize}
  \item \textbf{World measurement}: $p(A|\psi) = 1$, $p(B|\psi) = p(C|\psi) = p(D|\psi) = 0$. Post-measurement: $\qket{\psi_A} = \qket{\psi}$ (unchanged).

  \item \textbf{Idea measurement}: $p(An|\psi) = \frac{1}{10}$ for each $n$. Post-measurement: $\qket{An}$ for outcome $An$.

  \item \textbf{Element-only measurement}: $p(n|\psi) = \frac{1}{10}$ for each $n$. Post-measurement: $\qket{An}$ (since only A has support).
\end{itemize}
\end{example}

\subsection{Post-Measurement State Collapse}

\begin{theorem}[L\"uders Rule]
\label{thm:luders}
For a projective measurement with PVM $\{\Projector_m\}$ on a density operator $\rho$:
\begin{enumerate}
  \item \textbf{Outcome probability}: $p(m|\rho) = \Trace(\Projector_m \rho)$

  \item \textbf{Post-measurement state} (selective, outcome $m$ observed):
  \[
  \rho_m = \frac{\Projector_m \rho \Projector_m}{\Trace(\Projector_m \rho)}
  \]

  \item \textbf{Post-measurement state} (non-selective, outcome unknown):
  \[
  \rho' = \sum_m \Projector_m \rho \Projector_m
  \]
\end{enumerate}
\end{theorem}

\begin{proof}
\textbf{(1)}: For $\rho = \sum_i p_i \qket{\psi_i}\qbra{\psi_i}$:
\[
p(m|\rho) = \sum_i p_i p(m|\psi_i) = \sum_i p_i \qbra{\psi_i}\Projector_m\qket{\psi_i} = \Trace(\Projector_m \rho)
\]

\textbf{(2)}: The L\"uders rule is the unique update rule satisfying repeatability (measuring again gives same outcome) and minimal disturbance (does not change states already in the eigenspace).

\textbf{(3)}: Summing over all outcomes weighted by probability gives the non-selective update:
\[
\rho' = \sum_m p(m|\rho) \cdot \rho_m = \sum_m \Projector_m \rho \Projector_m
\]
\end{proof}

\begin{corollary}[Measurement as Channel]
\label{cor:measurement-channel}
Non-selective projective measurement defines a quantum channel:
\[
\mathcal{M}(\rho) = \sum_m \Projector_m \rho \Projector_m
\]
with Kraus operators $K_m = \Projector_m$. This is a \textbf{completely dephasing channel} in the measurement basis.
\end{corollary}

\begin{proposition}[Measurement Destroys Coherence]
\label{prop:measurement-coherence}
For non-selective measurement with PVM $\{\Projector_m\}$:
\[
\qbra{m}\rho'\qket{m'} = \delta_{mm'} \qbra{m}\rho\qket{m}
\]
All off-diagonal elements (coherences) between different measurement outcomes are destroyed.
\end{proposition}

\begin{definition}[Repeatability]
\label{def:repeatability}
A measurement is \textbf{repeatable} if measuring again immediately yields the same outcome with certainty. Projective measurements are repeatable:
\[
p(m|\rho_m) = \Trace(\Projector_m \rho_m) = \Trace\left(\Projector_m \frac{\Projector_m \rho \Projector_m}{\Trace(\Projector_m \rho)}\right) = 1
\]
\end{definition}

\subsection{Measurement on Bipartite States}

\begin{definition}[Local Measurement]
\label{def:local-measurement}
A \textbf{local measurement} on subsystem $\mathsf{A}$ of a bipartite state $\rho_{AB}$ uses PVM $\{\Projector_m \otimes \mathbf{I}_\mathsf{B}\}$:
\begin{itemize}
  \item \textbf{Probability}: $p(m) = \Trace((\Projector_m \otimes \mathbf{I}) \rho_{AB})$
  \item \textbf{Post-measurement}: $\rho_{AB}^{(m)} = \frac{(\Projector_m \otimes \mathbf{I}) \rho_{AB} (\Projector_m \otimes \mathbf{I})}{p(m)}$
\end{itemize}
\end{definition}

\begin{theorem}[Entanglement and Measurement Correlations]
\label{thm:entanglement-correlations}
For the Bell state $\qket{\Phi^+} = \frac{1}{\sqrt{40}}\sum_I \qket{I, I}$:
\begin{enumerate}
  \item Local Idea measurement on $\mathsf{A}$ yielding outcome $I$ leaves $\mathsf{B}$ in state $\qket{I}$
  \item The joint probability for outcomes $(I, J)$ under local measurements on both subsystems is:
  \[
  p(I, J) = \frac{1}{40} \delta_{IJ}
  \]
  \item Outcomes are perfectly correlated: $\mathsf{A}$ and $\mathsf{B}$ always yield the same Idea
\end{enumerate}
\end{theorem}

\begin{proof}
\textbf{(1)}: $(\Projector_I \otimes \mathbf{I})\qket{\Phi^+} = \frac{1}{\sqrt{40}}\qket{I, I}$. After normalization: $\qket{I, I}$.

\textbf{(2)}: $p(I, J) = |\qbra{I, J}\Phi^+\rangle|^2 = \frac{1}{40}\delta_{IJ}$.

\textbf{(3)}: Follows from (2): $p(I \neq J) = 0$.
\end{proof}

%% ----------------------------------------------------------------------------
\section{POVM Measurements}
\label{sec:povm-measurements}

This section develops the general theory of positive operator-valued measures, extending v7.2 Definition~\ref{def:povm}.

\subsection{POVM Definition and Properties}

\begin{definition}[Positive Operator-Valued Measure]
\label{def:povm-v73}
A \textbf{positive operator-valued measure (POVM)} on $\Hspace_{\Ideas}$ is a collection $\{E_m\}_{m \in \mathcal{M}}$ of positive semidefinite operators satisfying:
\begin{enumerate}
  \item \textbf{Positivity}: $E_m \geq 0$ (i.e., $\qbra{\psi}E_m\qket{\psi} \geq 0$ for all $\qket{\psi}$)
  \item \textbf{Completeness}: $\sum_{m \in \mathcal{M}} E_m = \mathbf{I}$
\end{enumerate}
The operators $E_m$ are called \textbf{POVM elements} or \textbf{effects}.
\end{definition}

\begin{remark}[PVM as Special POVM]
\label{rem:pvm-povm}
Every PVM is a POVM (projectors are positive). POVMs generalize PVMs by allowing:
\begin{itemize}
  \item Non-orthogonal elements: $E_m E_{m'} \neq 0$ for $m \neq m'$
  \item Non-idempotent elements: $E_m^2 \neq E_m$
  \item More outcomes than dimension: $|\mathcal{M}| > \dim(\Hspace_{\Ideas}) = 40$
\end{itemize}
\end{remark}

\begin{definition}[POVM Measurement]
\label{def:povm-measurement}
A \textbf{POVM measurement} with effects $\{E_m\}$ on state $\rho$ yields:
\begin{itemize}
  \item \textbf{Outcome probability}: $p(m|\rho) = \Trace(E_m \rho)$
\end{itemize}
Note: POVMs specify only outcome probabilities, not post-measurement states. For state update, we need additional structure (Kraus operators).
\end{definition}

\begin{proposition}[POVM Probability Properties]
\label{prop:povm-probs}
For any POVM $\{E_m\}$ and density operator $\rho$:
\begin{enumerate}
  \item $p(m|\rho) \geq 0$ for all $m$
  \item $\sum_m p(m|\rho) = 1$
  \item $p(m|\rho) \leq 1$ for all $m$
\end{enumerate}
\end{proposition}

\begin{proof}
\textbf{(1)}: $p(m|\rho) = \Trace(E_m \rho) \geq 0$ since $E_m \geq 0$ and $\rho \geq 0$.

\textbf{(2)}: $\sum_m p(m|\rho) = \sum_m \Trace(E_m \rho) = \Trace((\sum_m E_m) \rho) = \Trace(\rho) = 1$.

\textbf{(3)}: Follows from (1) and (2).
\end{proof}

\subsection{POVM from Kraus Operators}

\begin{theorem}[Kraus-POVM Correspondence]
\label{thm:kraus-povm}
\textbf{(Neumark's Theorem, Operational Form)}

Every POVM $\{E_m\}$ can be realized by a measurement channel with Kraus operators $\{K_{m,\alpha}\}$ where:
\[
E_m = \sum_\alpha K_{m,\alpha}^\dagger K_{m,\alpha}
\]
The post-measurement state conditional on outcome $m$ is:
\[
\rho_m = \frac{\sum_\alpha K_{m,\alpha} \rho K_{m,\alpha}^\dagger}{\Trace(E_m \rho)}
\]
\end{theorem}

\begin{definition}[Minimal Kraus Realization]
\label{def:minimal-kraus}
A \textbf{minimal Kraus realization} of POVM element $E_m$ uses a single Kraus operator:
\[
K_m = \sqrt{E_m}
\]
where $\sqrt{E_m}$ is the unique positive square root of $E_m \geq 0$. Then $E_m = K_m^\dagger K_m = K_m^2$.
\end{definition}

\begin{corollary}[POVM State Update]
\label{cor:povm-update}
For a POVM with minimal Kraus realization $K_m = \sqrt{E_m}$, the post-measurement state is:
\[
\rho_m = \frac{\sqrt{E_m} \rho \sqrt{E_m}}{\Trace(E_m \rho)}
\]
\end{corollary}

\subsection{Important POVMs for TKS}

\begin{definition}[World POVM]
\label{def:world-povm-v73}
The \textbf{World POVM} $\{E_A, E_B, E_C, E_D\}$ with $E_W = \Projector_W$ is actually a PVM. It determines which World the Idea belongs to.
\end{definition}

\begin{definition}[Fuzzy Idea POVM]
\label{def:fuzzy-idea-povm}
The \textbf{$\epsilon$-fuzzy Idea POVM} models imperfect Idea discrimination:
\[
E_I^\epsilon = (1 - \epsilon)\qket{I}\qbra{I} + \frac{\epsilon}{39} \sum_{J \neq I} \qket{J}\qbra{J}
\]
where $\epsilon \in [0, 1)$ is the error parameter. Properties:
\begin{itemize}
  \item $E_I^\epsilon \geq 0$ for $\epsilon \in [0, 1]$
  \item $\sum_I E_I^\epsilon = \mathbf{I}$
  \item $\epsilon = 0$: Perfect discrimination (PVM)
  \item $\epsilon \to 1$: No discrimination ($E_I^\epsilon \to \frac{1}{40}\mathbf{I}$)
\end{itemize}
\end{definition}

\begin{proposition}[Fuzzy POVM Probabilities]
\label{prop:fuzzy-probs}
For state $\qket{\psi} = \sum_I c_I \qket{I}$ and fuzzy POVM:
\[
p(I|\psi, \epsilon) = (1 - \epsilon)|c_I|^2 + \frac{\epsilon}{39}(1 - |c_I|^2) = (1 - \frac{40\epsilon}{39})|c_I|^2 + \frac{\epsilon}{39}
\]
As $\epsilon$ increases, probabilities approach uniform $\frac{1}{40}$.
\end{proposition}

\begin{definition}[Symmetric Informationally Complete POVM]
\label{def:sic-povm}
A \textbf{SIC-POVM} on $\Hspace_{\Ideas}$ consists of $40^2 = 1600$ effects:
\[
E_{ij} = \frac{1}{40} \qket{\phi_{ij}}\qbra{\phi_{ij}}
\]
where $\{\qket{\phi_{ij}}\}$ are $1600$ vectors satisfying:
\[
|\langle\phi_{ij}|\phi_{kl}\rangle|^2 = \frac{1}{41} \quad \text{for } (i,j) \neq (k,l)
\]
SIC-POVMs are optimal for quantum state tomography.
\end{definition}

\begin{definition}[Coarse-Grained POVM]
\label{def:coarse-povm}
Given a fine-grained POVM $\{E_m\}$ and a partition $\mathcal{M} = \bigsqcup_k \mathcal{M}_k$, the \textbf{coarse-grained POVM} is:
\[
\tilde{E}_k = \sum_{m \in \mathcal{M}_k} E_m
\]
Example: The World POVM is the coarse-graining of the Idea PVM by World.
\end{definition}

\subsection{Neumark's Dilation Theorem}

\begin{theorem}[Neumark's Theorem]
\label{thm:neumark}
Every POVM $\{E_m\}_{m=1}^M$ on $\Hspace_{\Ideas}$ can be realized as a projective measurement on an extended Hilbert space:
\[
\Hspace_{\mathrm{ext}} = \Hspace_{\Ideas} \otimes \Hspace_{\mathrm{anc}}
\]
where $\Hspace_{\mathrm{anc}}$ is an ancilla space. Specifically, there exists:
\begin{enumerate}
  \item An isometry $V: \Hspace_{\Ideas} \to \Hspace_{\mathrm{ext}}$
  \item A PVM $\{\Projector_m\}$ on $\Hspace_{\mathrm{ext}}$
\end{enumerate}
such that $E_m = V^\dagger \Projector_m V$.
\end{theorem}

\begin{corollary}[Measurement Implementation]
\label{cor:measurement-implementation}
Any generalized measurement on a noetic state can be physically implemented by:
\begin{enumerate}
  \item Coupling to an ancilla system (e.g., a measurement apparatus)
  \item Performing a projective measurement on the joint system
  \item Recording the outcome
\end{enumerate}
\end{corollary}

%% ----------------------------------------------------------------------------
\section{Noetic Observables}
\label{sec:noetic-observables}

This section develops the theory of observables (self-adjoint operators) that correspond to measurable properties in TKS.

\subsection{General Observable Theory}

\begin{definition}[Observable]
\label{def:observable-v73}
A \textbf{noetic observable} is a self-adjoint (Hermitian) operator $\qop{O} = \qop{O}^\dagger$ on $\Hspace_{\Ideas}$. By the spectral theorem:
\[
\qop{O} = \sum_{\lambda \in \Spectrum(\qop{O})} \lambda \, \Projector_\lambda
\]
where $\Spectrum(\qop{O})$ is the spectrum (set of eigenvalues) and $\Projector_\lambda$ projects onto the $\lambda$-eigenspace.
\end{definition}

\begin{definition}[Expectation Value]
\label{def:expectation}
The \textbf{expectation value} of observable $\qop{O}$ in state $\rho$ is:
\[
\langle \qop{O} \rangle_\rho = \Trace(\qop{O} \rho)
\]
For pure state $\qket{\psi}$: $\langle \qop{O} \rangle_\psi = \qbra{\psi}\qop{O}\qket{\psi}$.
\end{definition}

\begin{definition}[Variance and Standard Deviation]
\label{def:variance}
The \textbf{variance} of observable $\qop{O}$ in state $\rho$ is:
\[
(\Delta O)^2_\rho = \langle \qop{O}^2 \rangle_\rho - \langle \qop{O} \rangle_\rho^2 = \Trace(\qop{O}^2 \rho) - (\Trace(\qop{O} \rho))^2
\]
The \textbf{standard deviation} is $\Delta O = \sqrt{(\Delta O)^2}$.
\end{definition}

\begin{proposition}[Observable Statistics]
\label{prop:observable-stats}
For observable $\qop{O} = \sum_\lambda \lambda \Projector_\lambda$ measured on state $\rho$:
\begin{enumerate}
  \item \textbf{Outcome probability}: $p(\lambda) = \Trace(\Projector_\lambda \rho)$
  \item \textbf{Expectation}: $\langle \qop{O} \rangle = \sum_\lambda \lambda \, p(\lambda)$
  \item \textbf{Variance}: $(\Delta O)^2 = \sum_\lambda (\lambda - \langle \qop{O} \rangle)^2 p(\lambda)$
  \item \textbf{Eigenstate certainty}: $(\Delta O)^2 = 0$ iff $\rho$ is supported on a single eigenspace
\end{enumerate}
\end{proposition}

\subsection{World Observable}

\begin{definition}[World Number Observable]
\label{def:world-observable}
The \textbf{World number observable} assigns numerical values to Worlds:
\[
\qop{W} = 1 \cdot \Projector_A + 2 \cdot \Projector_B + 3 \cdot \Projector_C + 4 \cdot \Projector_D = \sum_{W \in \{A,B,C,D\}} w(W) \Projector_W
\]
where $w(A) = 1$, $w(B) = 2$, $w(C) = 3$, $w(D) = 4$.
\end{definition}

\begin{proposition}[World Observable Properties]
\label{prop:world-observable}
The World observable $\qop{W}$ satisfies:
\begin{enumerate}
  \item $\Spectrum(\qop{W}) = \{1, 2, 3, 4\}$ (four eigenvalues)
  \item Each eigenspace has dimension 10 (10 Elements per World)
  \item For state $\qket{\psi} = \sum_{W,n} c_{Wn}\qket{Wn}$:
  \[
  \langle \qop{W} \rangle_\psi = \sum_W w(W) \cdot p(W|\psi) = \sum_W w(W) \sum_{n=1}^{10} |c_{Wn}|^2
  \]
\end{enumerate}
\end{proposition}

\begin{example}[World Expectation Values]
\label{ex:world-expectations}
\begin{itemize}
  \item Pure World A state: $\langle \qop{W} \rangle = 1$ (Spiritual)
  \item Pure World D state: $\langle \qop{W} \rangle = 4$ (Physical)
  \item Equal superposition $\qket{\Phi^+_{\mathrm{all}}} = \frac{1}{\sqrt{40}}\sum_I \qket{I}$: $\langle \qop{W} \rangle = \frac{1+2+3+4}{4} = 2.5$
  \item ACBE descent interpretation: $\langle \qop{W} \rangle$ increases as Ideas descend from Spiritual to Physical
\end{itemize}
\end{example}

\begin{definition}[World Parity Observable]
\label{def:world-parity}
The \textbf{World parity observable} distinguishes upper (A, B) from lower (C, D) Worlds:
\[
\qop{W}_{\mathrm{parity}} = \Projector_A + \Projector_B - \Projector_C - \Projector_D
\]
with eigenvalues $\pm 1$. Measures whether the Idea is ``higher'' (Spiritual/Psychic) or ``lower'' (Phenomenal/Physical).
\end{definition}

\subsection{Element Observable}

\begin{definition}[Element Number Observable]
\label{def:element-observable}
The \textbf{Element number observable} gives the Element index within any World:
\[
\qop{N} = \sum_{W \in \{A,B,C,D\}} \sum_{n=1}^{10} n \cdot \qket{Wn}\qbra{Wn}
\]
This has spectrum $\{1, 2, \ldots, 10\}$, each with multiplicity 4 (one per World).
\end{definition}

\begin{proposition}[Element Observable Properties]
\label{prop:element-observable}
The Element observable $\qop{N}$ satisfies:
\begin{enumerate}
  \item $[\qop{W}, \qop{N}] = 0$ (World and Element commute)
  \item Simultaneous eigenstates are exactly $\{\qket{Wn}\}_{W,n}$
  \item Joint eigenvalues $(w, n)$ uniquely identify each Idea
\end{enumerate}
\end{proposition}

\begin{proof}
\textbf{(1)}: Both $\qop{W}$ and $\qop{N}$ are diagonal in the Idea basis, hence commute.

\textbf{(2)}: $\qop{W}\qket{Wn} = w(W)\qket{Wn}$ and $\qop{N}\qket{Wn} = n\qket{Wn}$.

\textbf{(3)}: The pair $(w(W), n)$ uniquely specifies $Wn$ since $w$ is injective on Worlds.
\end{proof}

\begin{definition}[Element Parity Observable]
\label{def:element-parity}
The \textbf{Element parity observable}:
\[
\qop{N}_{\mathrm{parity}} = \sum_{W,n} (-1)^n \qket{Wn}\qbra{Wn}
\]
has eigenvalues $\pm 1$ distinguishing odd and even Elements.
\end{definition}

\subsection{Noetic-Induced Observables}

\begin{definition}[Noetic Occupation Observable]
\label{def:noetic-occupation}
For noetic $k \in \{1, \ldots, 22\}$, the \textbf{noetic $k$ occupation observable} is:
\[
\qop{O}_k = \qop{\nu}_k^\dagger \qop{\nu}_k
\]
This measures whether the state has support in the image of noetic $\nu_k$.
\end{definition}

\begin{proposition}[Noetic Occupation Properties]
\label{prop:noetic-occupation}
For noetic occupation observable $\qop{O}_k = \qop{\nu}_k^\dagger \qop{\nu}_k$:
\begin{enumerate}
  \item $\qop{O}_k$ is positive semidefinite: $\qop{O}_k \geq 0$
  \item If $\qop{\nu}_k$ is a partial isometry (idempotent noetic), then $\qop{O}_k$ is a projector
  \item $\langle \qop{O}_k \rangle_\psi = \|\qop{\nu}_k\qket{\psi}\|^2$: measures ``how much'' $\qket{\psi}$ is affected by $\nu_k$
  \item $\Trace(\qop{O}_k) = \mathrm{rank}(\qop{\nu}_k) = |\mathrm{Im}(\nu_k)|$
\end{enumerate}
\end{proposition}

\begin{definition}[Noetic Activity Observable]
\label{def:noetic-activity}
The \textbf{total noetic activity observable} sums over all noetics:
\[
\qop{A}_{\mathrm{total}} = \sum_{k=1}^{22} \qop{\nu}_k^\dagger \qop{\nu}_k
\]
High expectation value indicates the state is ``active'' under many noetics.
\end{definition}

\begin{definition}[Noetic Transition Observable]
\label{def:noetic-transition}
For noetics $j, k$, the \textbf{transition observable} from $\nu_j$ to $\nu_k$ is:
\[
\qop{T}_{jk} = \qop{\nu}_k^\dagger \qop{\nu}_j + \qop{\nu}_j^\dagger \qop{\nu}_k
\]
This is Hermitian and measures coherence between the two noetic transformations.
\end{definition}

\subsection{Compatible and Incompatible Observables}

\begin{definition}[Compatible Observables]
\label{def:compatible}
Observables $\qop{O}_1$ and $\qop{O}_2$ are \textbf{compatible} (simultaneously measurable) if they commute:
\[
[\qop{O}_1, \qop{O}_2] = \qop{O}_1 \qop{O}_2 - \qop{O}_2 \qop{O}_1 = 0
\]
Compatible observables share a common eigenbasis.
\end{definition}

\begin{proposition}[TKS Compatible Pairs]
\label{prop:tks-compatible}
The following pairs of observables are compatible:
\begin{enumerate}
  \item $(\qop{W}, \qop{N})$: World and Element
  \item $(\qop{W}, \qop{W}_{\mathrm{parity}})$: World number and World parity
  \item $(\qop{N}, \qop{N}_{\mathrm{parity}})$: Element number and Element parity
  \item $(\qop{O}_k, \qop{O}_k)$: Any observable with itself
\end{enumerate}
\end{proposition}

\begin{theorem}[Robertson Uncertainty Relation]
\label{thm:robertson}
For any observables $\qop{O}_1, \qop{O}_2$ and state $\rho$:
\[
\Delta O_1 \cdot \Delta O_2 \geq \frac{1}{2} |\langle [\qop{O}_1, \qop{O}_2] \rangle_\rho|
\]
If $[\qop{O}_1, \qop{O}_2] \neq 0$, there exist states with nonzero uncertainty product.
\end{theorem}

\begin{corollary}[Noetic Uncertainty Relations]
\label{cor:noetic-uncertainty}
For non-commuting noetic operators $\qop{\nu}_j, \qop{\nu}_k$:
\[
\Delta(\qop{\nu}_j^\dagger \qop{\nu}_j) \cdot \Delta(\qop{\nu}_k^\dagger \qop{\nu}_k) \geq \frac{1}{2} |\langle [\qop{\nu}_j^\dagger \qop{\nu}_j, \qop{\nu}_k^\dagger \qop{\nu}_k] \rangle|
\]
Certain pairs of noetic occupations cannot both be sharply defined.
\end{corollary}

%% ----------------------------------------------------------------------------
\section{Measurement Disturbance}
\label{sec:measurement-disturbance}

This section analyzes how measurement disturbs the noetic state.

\subsection{State Disturbance}

\begin{definition}[Measurement Disturbance]
\label{def:disturbance}
The \textbf{disturbance} caused by measurement $\mathcal{M}$ on state $\rho$ is quantified by:
\[
D(\rho, \mathcal{M}) = \frac{1}{2}\|\rho - \mathcal{M}(\rho)\|_1
\]
where $\mathcal{M}(\rho)$ is the non-selective post-measurement state and $\|\cdot\|_1$ is trace norm.
\end{definition}

\begin{proposition}[Disturbance Bounds]
\label{prop:disturbance-bounds}
\begin{enumerate}
  \item $0 \leq D(\rho, \mathcal{M}) \leq 1$
  \item $D(\rho, \mathcal{M}) = 0$ iff $\rho$ is unchanged by measurement (e.g., $\rho$ diagonal in measurement basis)
  \item For projective measurement in basis $\{|m\rangle\}$: $D(\rho, \mathcal{M}) = \sum_{m \neq m'} |\rho_{mm'}|$ (sum of off-diagonal magnitudes)
\end{enumerate}
\end{proposition}

\begin{theorem}[Information-Disturbance Tradeoff]
\label{thm:info-disturbance}
For any measurement on a quantum state, there is a fundamental tradeoff between:
\begin{itemize}
  \item \textbf{Information gain}: How much the measurement outcome tells us about the state
  \item \textbf{State disturbance}: How much the measurement changes the state
\end{itemize}
Perfect information (complete measurement) maximizes disturbance for non-eigenstate inputs.
\end{theorem}

\subsection{Measurement Back-Action on Entangled States}

\begin{theorem}[Measurement-Induced Collapse of Entanglement]
\label{thm:measurement-collapse}
For a maximally entangled state $\qket{\Phi^+} \in \Hspace_{\Ideas}^{(2)}$:
\begin{enumerate}
  \item \textbf{Before local measurement}: $E(\qket{\Phi^+}) = \log(40)$ (maximal entanglement)
  \item \textbf{After local Idea measurement}: The post-measurement state is a product state with $E = 0$
  \item \textbf{After local World measurement}: Entanglement reduced to $E \leq \log(10)$
\end{enumerate}
Local measurement generically decreases entanglement.
\end{theorem}

\begin{proof}
\textbf{(2)}: If Idea $I$ is measured on subsystem A, the post-measurement state is $\qket{I, I}$, a product state.

\textbf{(3)}: If World $W$ is measured on A, the state collapses to $\qket{\Phi^+_{WW}}$, which has entanglement $\log(10)$.
\end{proof}

\subsection{Quantum Zeno Effect}

\begin{theorem}[Quantum Zeno Effect in TKS]
\label{thm:zeno}
Frequent repeated measurements can freeze the evolution of a noetic state. If measurement $\mathcal{M}$ is performed $n$ times during time $T$ with evolution $\qop{U}(T/n)$ between measurements:
\[
\lim_{n \to \infty} [\mathcal{M} \circ \qop{U}(T/n)]^n (\rho) = \mathcal{M}(\rho)
\]
The state is ``frozen'' in the measurement subspace.
\end{theorem}

\begin{corollary}[Zeno Effect for Noetic Evolution]
\label{cor:zeno-noetic}
Continuously measuring the World observable prevents transitions between Worlds under noetic evolution. An Idea in World A stays in World A despite noetic operators that would otherwise cause descent.
\end{corollary}

%% ----------------------------------------------------------------------------
\section{Weak Measurements}
\label{sec:weak-measurements}

This section develops weak measurement theory, which extracts partial information with minimal disturbance.

\subsection{Weak Measurement Formalism}

\begin{definition}[Weak Measurement]
\label{def:weak-measurement}
A \textbf{weak measurement} of observable $\qop{O}$ on system state $\qket{\psi}$ is implemented by:
\begin{enumerate}
  \item Coupling to a pointer (ancilla) initially in state $\qket{\phi_0}$
  \item Interaction Hamiltonian $H_{\mathrm{int}} = g \qop{O} \otimes \qop{p}$ for short time, where $\qop{p}$ is the pointer momentum
  \item Measuring the pointer position $\qop{x}$
\end{enumerate}
In the weak coupling limit ($g \to 0$), the pointer shift gives information about $\langle \qop{O} \rangle$ with minimal disturbance to $\qket{\psi}$.
\end{definition}

\begin{definition}[Weak Value]
\label{def:weak-value}
The \textbf{weak value} of observable $\qop{O}$ with pre-selection $\qket{\psi_i}$ and post-selection $\qket{\psi_f}$ is:
\[
O_w = \frac{\qbra{\psi_f}\qop{O}\qket{\psi_i}}{\langle\psi_f|\psi_i\rangle}
\]
The weak value is complex-valued in general and can lie outside the spectrum of $\qop{O}$.
\end{definition}

\begin{proposition}[Weak Value Properties]
\label{prop:weak-value-props}
Weak values satisfy:
\begin{enumerate}
  \item \textbf{Reduces to eigenvalue}: If $\qket{\psi_i} = \qket{\psi_f} = \qket{\lambda}$ (eigenstate of $\qop{O}$), then $O_w = \lambda$
  \item \textbf{Reduces to expectation}: If $\qket{\psi_f} = \qket{\psi_i}$, then $O_w = \langle \qop{O} \rangle_{\psi_i}$
  \item \textbf{Anomalous values}: When $\langle\psi_f|\psi_i\rangle \approx 0$, $O_w$ can be arbitrarily large (even outside eigenvalue range)
  \item \textbf{Linearity}: $(a\qop{O}_1 + b\qop{O}_2)_w = a(\qop{O}_1)_w + b(\qop{O}_2)_w$
\end{enumerate}
\end{proposition}

\begin{example}[Anomalous Weak Value in TKS]
\label{ex:anomalous-weak}
Consider pre-selection $\qket{\psi_i} = \frac{1}{\sqrt{2}}(\qket{A1} + \qket{D1})$ and post-selection $\qket{\psi_f} = \frac{1}{\sqrt{2}}(\qket{A1} - \qket{D1})$. For the World observable:
\[
W_w = \frac{\qbra{\psi_f}\qop{W}\qket{\psi_i}}{\langle\psi_f|\psi_i\rangle} = \frac{\frac{1}{2}(1 - 4)}{\frac{1}{2}(1 - 1)} = \frac{-3/2}{0}
\]
The denominator is zero (orthogonal states), indicating this post-selection never occurs. Near-orthogonal cases yield very large weak values.

For $\qket{\psi_f} = \frac{1}{\sqrt{2}}(\qket{A1} - e^{i\epsilon}\qket{D1})$ with small $\epsilon$:
\[
W_w \approx \frac{-3/2}{i\epsilon/2} = \frac{3}{i\epsilon}
\]
yielding a large imaginary weak value.
\end{example}

\subsection{Weak Values in TKS Context}

\begin{definition}[Weak World Number]
\label{def:weak-world}
The \textbf{weak World number} for pre-selection $\qket{\psi_i}$ and post-selection $\qket{\psi_f}$ is:
\[
W_w = \frac{\qbra{\psi_f}\qop{W}\qket{\psi_i}}{\langle\psi_f|\psi_i\rangle} = \frac{\sum_W w(W) \langle\psi_f|\Projector_W|\psi_i\rangle}{\langle\psi_f|\psi_i\rangle}
\]
This can indicate an ``effective World'' between pre- and post-selected states that lies between or outside $\{1, 2, 3, 4\}$.
\end{definition}

\begin{definition}[Weak Noetic Value]
\label{def:weak-noetic}
For noetic operator $\qop{\nu}_k$, the \textbf{weak noetic value} is:
\[
(\nu_k)_w = \frac{\qbra{\psi_f}\qop{\nu}_k\qket{\psi_i}}{\langle\psi_f|\psi_i\rangle}
\]
This measures the ``effective noetic transformation'' between pre- and post-selected states.
\end{definition}

\begin{proposition}[Weak Value Decomposition]
\label{prop:weak-decomposition}
Any weak value can be decomposed into real and imaginary parts:
\[
O_w = \mathrm{Re}(O_w) + i \, \mathrm{Im}(O_w)
\]
\begin{itemize}
  \item $\mathrm{Re}(O_w)$: Related to shift in pointer position
  \item $\mathrm{Im}(O_w)$: Related to shift in pointer momentum (back-action)
\end{itemize}
\end{proposition}

\subsection{Weak Measurement and Classical TKS}

\begin{theorem}[Weak Measurement Classical Correspondence]
\label{thm:weak-classical}
\textbf{(Main Theorem: Measurement-Classical Connection)}

In the appropriate classical limit, weak measurements on noetic states correspond to classical TKS observations:
\begin{enumerate}
  \item \textbf{World weak value}: For states with support on multiple Worlds, the real part of $W_w$ gives a weighted average World position, corresponding to the classical ``location'' of an Idea in the World hierarchy.

  \item \textbf{Noetic weak value}: For states evolving under noetic $\nu_k$, the weak value $(\nu_k)_w$ with appropriate pre/post-selection gives the classical noetic transition amplitude.

  \item \textbf{Classical limit}: When $\qket{\psi_i}$ and $\qket{\psi_f}$ are both localized on single Ideas $I$ and $J$:
  \[
  (\nu_k)_w = \frac{\langle J|\qop{\nu}_k|I\rangle}{\langle J|I\rangle} = \begin{cases} 1 & \text{if } \nu_k(I) = J \\ 0 & \text{if } \nu_k(I) \neq J, I = J \\ \text{undefined} & \text{if } I \neq J \end{cases}
  \]
  This recovers the classical noetic function: $\nu_k(I) = J$ iff the weak value is 1.

  \item \textbf{Superposition information}: For superposition states, the weak value encodes interference between classical paths, providing quantum corrections to classical TKS dynamics.
\end{enumerate}
\end{theorem}

\begin{proof}
\textbf{(1)}: For $\qket{\psi_i} = \sum_{W,n} c_{Wn}^{(i)}\qket{Wn}$ and $\qket{\psi_f} = \sum_{W,n} c_{Wn}^{(f)}\qket{Wn}$:
\[
W_w = \frac{\sum_W w(W) \sum_n (c_{Wn}^{(f)})^* c_{Wn}^{(i)}}{\sum_{W,n} (c_{Wn}^{(f)})^* c_{Wn}^{(i)}}
\]
When both states have similar World support, $\mathrm{Re}(W_w)$ gives a weighted average.

\textbf{(3)}: For $\qket{\psi_i} = \qket{I}$ and $\qket{\psi_f} = \qket{J}$:
\[
(\nu_k)_w = \frac{\langle J|\qop{\nu}_k|I\rangle}{\langle J|I\rangle}
\]
If $I = J$: denominator is 1, numerator is $\langle I|\qop{\nu}_k|I\rangle = \delta_{\nu_k(I), I}$.
If $I \neq J$: denominator is 0 (post-selection impossible without evolution).

When $\nu_k(I) = J$ and we use appropriate unitary evolution connecting $\qket{I}$ to $\qket{J}$, the weak value becomes 1, confirming the classical transition.

\textbf{(4)}: For superpositions $\qket{\psi} = \sum_I c_I\qket{I}$, the weak value involves interference terms $(c_J^{(f)})^* c_I^{(i)}$ that have no classical analog---these are the quantum corrections.
\end{proof}

\begin{corollary}[Observation Without Collapse]
\label{cor:observation-no-collapse}
Weak measurements allow extraction of information about noetic states with arbitrarily small disturbance:
\begin{enumerate}
  \item \textbf{Non-invasive probing}: In the weak limit, the state is nearly unchanged
  \item \textbf{Statistical reconstruction}: Many weak measurements on identically prepared states can reconstruct expectation values
  \item \textbf{Classical correspondence}: This is the quantum analog of classical observation, which does not disturb the observed system
\end{enumerate}
\end{corollary}

\begin{definition}[Sequential Weak Measurements]
\label{def:sequential-weak}
A sequence of weak measurements of observables $\qop{O}_1, \qop{O}_2, \ldots, \qop{O}_n$ with final post-selection yields:
\[
(O_n \cdots O_2 O_1)_w = \frac{\qbra{\psi_f}\qop{O}_n \cdots \qop{O}_2 \qop{O}_1\qket{\psi_i}}{\langle\psi_f|\psi_i\rangle}
\]
This corresponds to a noetic fractal: a sequence of noetic operations $\nu_{k_n} \circ \cdots \circ \nu_{k_1}$.
\end{definition}

\begin{theorem}[Weak Fractal Value]
\label{thm:weak-fractal}
For a noetic fractal $\Phi = \nu_{k_n} \circ \cdots \circ \nu_{k_1}$ with corresponding quantum operator $\qop{\Phi} = \qop{\nu}_{k_n} \cdots \qop{\nu}_{k_1}$:
\[
\Phi_w = \frac{\qbra{\psi_f}\qop{\Phi}\qket{\psi_i}}{\langle\psi_f|\psi_i\rangle}
\]
gives the weak value of the entire fractal process. When $\qket{\psi_i} = \qket{I}$ and $\qket{\psi_f} = \qket{\Phi(I)}$:
\[
\Phi_w = 1
\]
confirming the classical fractal transition.
\end{theorem}

%% ----------------------------------------------------------------------------
%% HANDOFF NOTE
%% ----------------------------------------------------------------------------
\begin{tcolorbox}[colback=green!5!white,colframe=green!50!black,title=\textbf{Handoff: Next Steps},breakable]
This completes Chapter 6 on Quantum Measurement in TKS. The subsequent chapters should:
\begin{itemize}
  \item \textbf{Chapter 7 (Quantum Fractals)}: Use measurement theory to analyze fractal fixed points; measurement-based computation on noetic states
  \item \textbf{Chapter 8 (Applications)}: Apply weak measurements to RPM goal achievement; POVM design for Idea discrimination
\end{itemize}

Key results established in this chapter:
\begin{enumerate}
  \item \textbf{Projective measurement}: PVM, Born rule, L\"uders update (Sec.~\ref{sec:projective-measurements})
  \item \textbf{Measurement types}: Idea, World, Element measurements with distinct post-measurement states
  \item \textbf{POVM formalism}: Generalized measurements, Kraus realization, Neumark dilation (Sec.~\ref{sec:povm-measurements})
  \item \textbf{Noetic observables}: World $\qop{W}$, Element $\qop{N}$, noetic occupation $\qop{O}_k$ (Sec.~\ref{sec:noetic-observables})
  \item \textbf{Compatibility}: $[\qop{W}, \qop{N}] = 0$; uncertainty relations for non-commuting noetics
  \item \textbf{Measurement disturbance}: Information-disturbance tradeoff, Zeno effect (Sec.~\ref{sec:measurement-disturbance})
  \item \textbf{Weak measurements}: Weak values, anomalous values (Sec.~\ref{sec:weak-measurements})
  \item \textbf{Classical correspondence}: Theorem~\ref{thm:weak-classical} connects weak measurements to classical TKS observations
  \item \textbf{Weak fractal values}: Theorem~\ref{thm:weak-fractal} gives weak values for fractal sequences
\end{enumerate}

\textbf{Connection to v7.2}: Extends Definition~\ref{def:world-measurement}, \ref{def:element-measurement}, \ref{def:noetic-observable}, \ref{def:povm} from v7.2 Chapter 10 with complete mathematical treatment.

\textbf{Connection to Classical TKS}: Theorem~\ref{thm:weak-classical} establishes that weak measurements in the classical limit recover classical TKS observations. The weak noetic value $(\nu_k)_w = 1$ iff the classical noetic transition $\nu_k(I) = J$ occurs.

\textbf{Connection to Entanglement (Chapter 5)}: Theorem~\ref{thm:entanglement-correlations} shows Bell state correlations under local measurement. Theorem~\ref{thm:measurement-collapse} shows measurement reduces entanglement.
\end{tcolorbox}

%% ============================================================================
%% CHAPTER 7: QUANTUM FRACTALS
%% ============================================================================
\chapter{Quantum Noetic Fractals}
\label{ch:quantum-fractals}

\begin{tcolorbox}[colback=blue!5!white,colframe=blue!75!black,title=\vnewthree]
This chapter extends fractal calculus to quantum operators, developing quantum versions of fractal dynamics and attractors.
\end{tcolorbox}

%% ----------------------------------------------------------------------------
\section{Quantum Fractal Operators}
\label{sec:quantum-fractal-operators}

This section lifts classical noetic fractals (v7.0--v7.2) to quantum operators on the Hilbert space $\Hspace_{\Ideas}$.

\subsection{Finite Quantum Fractals}

\begin{definition}[Quantum Fractal Operator]
\label{def:quantum-fractal-op}
For a classical fractal $\Phi = \langle k_1 : k_2 : \cdots : k_n \rangle$ of length $n$, the \textbf{quantum fractal operator} is the composition:
\[
\qop{\Phi} = \qop{\nu}_{k_n} \qop{\nu}_{k_{n-1}} \cdots \qop{\nu}_{k_2} \qop{\nu}_{k_1} \in \BoundedOp(\Hspace_{\Ideas})
\]
where $\qop{\nu}_k$ are the quantum noetic operators from Chapter 3. The operators are applied right-to-left, matching the classical evaluation order.
\end{definition}

\begin{proposition}[Quantum-Classical Consistency]
\label{prop:qc-fractal-consistency}
The quantum fractal operator is consistent with classical fractal evaluation:
\[
\qop{\Phi}\qket{I} = \qket{\Phi(I)}
\]
for any Idea basis state $\qket{I}$, where $\Phi(I) = \mathfrak{n}_{k_n}(\cdots \mathfrak{n}_{k_2}(\mathfrak{n}_{k_1}(I))\cdots)$ is the classical evaluation.
\end{proposition}

\begin{proof}
By induction on $n$. For $n = 1$: $\qop{\nu}_{k_1}\qket{I} = \qket{\nu_{k_1}(I)}$ by Definition~\ref{def:quantum-noetic}. For $n + 1$:
\[
\qop{\Phi}_{n+1}\qket{I} = \qop{\nu}_{k_{n+1}}(\qop{\Phi}_n\qket{I}) = \qop{\nu}_{k_{n+1}}\qket{\Phi_n(I)} = \qket{\nu_{k_{n+1}}(\Phi_n(I))} = \qket{\Phi_{n+1}(I)}
\]
\end{proof}

\begin{definition}[Quantum Fractal on Superpositions]
\label{def:qfractal-superposition}
For a superposition state $\qket{\psi} = \sum_I c_I \qket{I}$:
\[
\qop{\Phi}\qket{\psi} = \sum_I c_I \qket{\Phi(I)}
\]
The quantum fractal acts linearly, preserving interference patterns.
\end{definition}

\begin{example}[Fractal on Equal Superposition]
\label{ex:fractal-superposition}
Let $\qket{\psi} = \frac{1}{\sqrt{10}}\sum_{n=1}^{10}\qket{An}$ (equal superposition over World A) and $\Phi = \langle 5 \rangle$ a single-step descent fractal. Then:
\[
\qop{\Phi}\qket{\psi} = \frac{1}{\sqrt{10}}\sum_{n=1}^{10}\qket{\nu_5(An)}
\]
If $\nu_5$ maps World A to World B, this becomes a superposition over World B, demonstrating coherent descent.
\end{example}

\subsection{Spectral Properties of Finite Quantum Fractals}

\begin{theorem}[Spectral Structure of Quantum Fractals]
\label{thm:fractal-spectrum}
For quantum fractal operator $\qop{\Phi}$:
\begin{enumerate}
  \item $\Spectrum(\qop{\Phi}) \subseteq \{0, 1\}$ when all $\qop{\nu}_k$ are idempotent
  \item $\|\qop{\Phi}\| \leq 1$ (contraction or isometry)
  \item $\qop{\Phi}$ has at least one eigenvalue $\lambda$ with $|\lambda| = 1$ if it has any fixed points
\end{enumerate}
\end{theorem}

\begin{proof}
\textbf{(1)}: If each $\qop{\nu}_k$ is idempotent ($\qop{\nu}_k^2 = \qop{\nu}_k$), then $\qop{\nu}_k$ is a projector with spectrum $\{0, 1\}$. The composition of projectors has spectrum in $\{0, 1\}$.

\textbf{(2)}: Each $\qop{\nu}_k$ satisfies $\|\qop{\nu}_k\| \leq 1$ by Proposition~\ref{prop:quantum-noetic-norm}. Hence $\|\qop{\Phi}\| \leq \prod_j \|\qop{\nu}_{k_j}\| \leq 1$.

\textbf{(3)}: If $\qket{I}$ is a classical fixed point ($\Phi(I) = I$), then $\qop{\Phi}\qket{I} = \qket{I}$, so $\lambda = 1$ is an eigenvalue.
\end{proof}

\begin{definition}[Fractal Rank]
\label{def:fractal-rank}
The \textbf{rank} of quantum fractal $\qop{\Phi}$ is:
\[
\mathrm{rank}(\qop{\Phi}) = \dim(\mathrm{Im}(\qop{\Phi})) = |\{\Phi(I) : I \in \Ideas\}|
\]
This counts the number of distinct output Ideas.
\end{definition}

\begin{proposition}[Rank Decrease]
\label{prop:rank-decrease}
For fractal composition:
\[
\mathrm{rank}(\qop{\Phi} \circ \qop{\Psi}) \leq \min(\mathrm{rank}(\qop{\Phi}), \mathrm{rank}(\qop{\Psi}))
\]
Longer fractals generically have lower rank (more contraction).
\end{proposition}

\subsection{Quantum Fractal Algebra}

\begin{theorem}[Quantum Fractal Monoid]
\label{thm:qfractal-monoid}
The quantum fractal operators form a monoid $(\mathfrak{Q}\FracSet, \cdot, \mathbf{I})$ where:
\begin{enumerate}
  \item Composition: $\qop{\Phi} \cdot \qop{\Psi} = \qop{\Phi \circ \Psi}$
  \item Identity: $\mathbf{I}$ (empty fractal)
  \item Associativity: $(\qop{\Phi} \cdot \qop{\Psi}) \cdot \qop{\Xi} = \qop{\Phi} \cdot (\qop{\Psi} \cdot \qop{\Xi})$
\end{enumerate}
This is the quantum lift of the classical fractal monoid (v7.0).
\end{theorem}

\begin{definition}[Adjoint Fractal]
\label{def:adjoint-fractal}
The \textbf{adjoint} of quantum fractal $\qop{\Phi} = \qop{\nu}_{k_n} \cdots \qop{\nu}_{k_1}$ is:
\[
\qop{\Phi}^\dagger = \qop{\nu}_{k_1}^\dagger \cdots \qop{\nu}_{k_n}^\dagger
\]
Note the reversed order. The adjoint fractal ``reverses'' the classical fractal path.
\end{definition}

\begin{proposition}[Self-Adjoint Fractals]
\label{prop:self-adjoint-fractal}
$\qop{\Phi}$ is self-adjoint ($\qop{\Phi} = \qop{\Phi}^\dagger$) iff:
\begin{enumerate}
  \item The fractal is a palindrome: $k_j = k_{n+1-j}$ for all $j$, AND
  \item Each $\qop{\nu}_{k_j}$ is self-adjoint
\end{enumerate}
Self-adjoint fractals correspond to reversible noetic processes.
\end{proposition}

%% ----------------------------------------------------------------------------
\section{Transfinite Quantum Fractals}
\label{sec:transfinite-quantum-fractals}

This section extends quantum fractals to transfinite ordinal indices, quantizing the v7.2 transfinite fractal theory.

\subsection{Ordinal-Indexed Quantum Operators}

\begin{definition}[Transfinite Quantum Fractal]
\label{def:transfinite-qfractal}
For a transfinite fractal $\TransFrac{\alpha}$ indexed by ordinal $\alpha$, the \textbf{transfinite quantum fractal operator} $\qop{\Phi}^{(\alpha)}$ is defined by transfinite recursion:
\begin{enumerate}
  \item \textbf{Base}: $\qop{\Phi}^{(0)} = \mathbf{I}$ (identity operator)

  \item \textbf{Successor}: For $\alpha = \beta + 1$ with $\TransFrac{\alpha}(\beta) = k$:
  \[
  \qop{\Phi}^{(\alpha)} = \qop{\nu}_k \cdot \qop{\Phi}^{(\beta)}
  \]

  \item \textbf{Limit}: For limit ordinal $\lambda$:
  \[
  \qop{\Phi}^{(\lambda)} = \mathrm{SOT}\text{-}\lim_{\beta \to \lambda} \qop{\Phi}^{(\beta)}
  \]
  where SOT denotes the strong operator topology.
\end{enumerate}
\end{definition}

\begin{theorem}[Transfinite Limit Existence]
\label{thm:transfinite-limit}
For any coherent system of transfinite fractals $(\TransFrac{\beta})_{\beta < \lambda}$, the strong operator limit $\qop{\Phi}^{(\lambda)}$ exists and is a bounded operator on $\Hspace_{\Ideas}$.
\end{theorem}

\begin{proof}
Consider the sequence of operators $(\qop{\Phi}^{(\beta)})_{\beta < \lambda}$. For each basis vector $\qket{I}$:
\begin{enumerate}
  \item The sequence $\qop{\Phi}^{(\beta)}\qket{I} = \qket{\TransFrac{\beta}(I)}$ takes values in the finite set $\Ideas$
  \item By the classical Stabilization Theorem (v7.2 Theorem~\ref{thm:stabilization}), there exists $\gamma_I < \lambda$ such that $\TransFrac{\beta}(I) = \TransFrac{\gamma_I}(I)$ for all $\beta \geq \gamma_I$
  \item Thus $\qop{\Phi}^{(\beta)}\qket{I}$ stabilizes for each $I$
\end{enumerate}
By linearity and the finite dimension of $\Hspace_{\Ideas}$, the strong limit exists:
\[
\qop{\Phi}^{(\lambda)}\qket{I} = \lim_{\beta \to \lambda} \qop{\Phi}^{(\beta)}\qket{I} = \qket{\TransFrac{\lambda}(I)}
\]
\end{proof}

\begin{definition}[$\omega$-Quantum Fractal]
\label{def:omega-qfractal}
The \textbf{$\omega$-quantum fractal} for noetic $k$ is:
\[
\qop{\Phi}^{(\omega)}_k = \mathrm{SOT}\text{-}\lim_{n \to \infty} \qop{\nu}_k^n
\]
This is the quantum analogue of the $\omega$-eigenfractal $\TransFrac{\omega}_k$ from v7.2.
\end{definition}

\begin{theorem}[Omega-Fractal Convergence]
\label{thm:omega-convergence}
For any noetic $k$:
\begin{enumerate}
  \item $\qop{\Phi}^{(\omega)}_k$ exists and is a projector
  \item $\qop{\Phi}^{(\omega)}_k = \qop{\Phi}^{(\omega)}_k \cdot \qop{\Phi}^{(\omega)}_k$ (idempotent)
  \item $\mathrm{Im}(\qop{\Phi}^{(\omega)}_k) = \mathrm{span}\{\qket{I} : \nu_k(I) = I\}$ (fixed point subspace)
\end{enumerate}
\end{theorem}

\begin{proof}
\textbf{(1)}: By finite dimensionality and the classical idempotence of noetics, each sequence $\nu_k^n(I)$ stabilizes.

\textbf{(2)}: $\qop{\Phi}^{(\omega)}_k \cdot \qop{\Phi}^{(\omega)}_k = \mathrm{SOT}\text{-}\lim_{n,m} \qop{\nu}_k^n \qop{\nu}_k^m = \mathrm{SOT}\text{-}\lim_{n} \qop{\nu}_k^{2n} = \qop{\Phi}^{(\omega)}_k$.

\textbf{(3)}: For $\qket{I}$ with $\nu_k(I) = I$: $\qop{\Phi}^{(\omega)}_k\qket{I} = \lim_n \qop{\nu}_k^n\qket{I} = \qket{I}$.
For $\qket{I}$ with $\nu_k(I) \neq I$: $\qop{\Phi}^{(\omega)}_k\qket{I} = \qket{J}$ where $J = \nu_k^\omega(I)$ is the classical fixed point.
\end{proof}

\subsection{Transfinite Operator Composition}

\begin{definition}[Ordinal Operator Composition]
\label{def:ordinal-op-composition}
For transfinite quantum fractals $\qop{\Phi}^{(\alpha)}$ and $\qop{\Phi}^{(\beta)}$:
\[
\qop{\Phi}^{(\alpha)} \circ_\Ord \qop{\Phi}^{(\beta)} = \qop{\Phi}^{(\alpha+\beta)}
\]
where $\TransFrac{\alpha+\beta} = \TransFrac{\alpha} \circ_\Ord \TransFrac{\beta}$ is the classical ordinal composition.
\end{definition}

\begin{theorem}[Transfinite Quantum Monoid]
\label{thm:transfinite-qmonoid}
The transfinite quantum fractals form a monoid under ordinal composition:
\begin{enumerate}
  \item \textbf{Associativity}: $(\qop{\Phi}^{(\alpha)} \circ_\Ord \qop{\Phi}^{(\beta)}) \circ_\Ord \qop{\Phi}^{(\gamma)} = \qop{\Phi}^{(\alpha)} \circ_\Ord (\qop{\Phi}^{(\beta)} \circ_\Ord \qop{\Phi}^{(\gamma)})$
  \item \textbf{Identity}: $\qop{\Phi}^{(0)} = \mathbf{I}$
\end{enumerate}
This quantizes the classical transfinite fractal monoid (v7.2 Corollary~\ref{cor:transfinite-monoid}).
\end{theorem}

\begin{corollary}[Transfinite Absorption]
\label{cor:transfinite-absorption}
For any finite quantum fractal $\qop{\Phi}^{(n)}$ and $\omega$-quantum fractal $\qop{\Phi}^{(\omega)}_k$:
\[
\qop{\Phi}^{(\omega)}_k \circ_\Ord \qop{\Phi}^{(n)} \sim \qop{\Phi}^{(\omega)}_k
\]
where $\sim$ denotes operational equivalence (same action on all states). The infinite fractal ``absorbs'' the finite one.
\end{corollary}

%% ----------------------------------------------------------------------------
\section{Quantum Eigenfractals}
\label{sec:quantum-eigenfractals}

This section develops the eigenstate theory for quantum fractal operators.

\subsection{Eigenstates of Finite Fractals}

\begin{definition}[Quantum Eigenfractal State]
\label{def:quantum-eigenfractal}
A state $\qket{\psi} \in \Hspace_{\Ideas}$ is a \textbf{quantum eigenfractal} of $\qop{\Phi}$ with eigenvalue $\lambda$ if:
\[
\qop{\Phi}\qket{\psi} = \lambda \qket{\psi}
\]
The set of all eigenfractal states forms the \textbf{eigenfractal subspace} $\mathcal{E}_\lambda(\qop{\Phi})$.
\end{definition}

\begin{theorem}[Classical Fixed Points as Eigenstates]
\label{thm:fixed-eigenstates}
If $I \in \Ideas$ is a classical fixed point of $\Phi$ (i.e., $\Phi(I) = I$), then $\qket{I}$ is an eigenfractal of $\qop{\Phi}$ with eigenvalue $\lambda = 1$.
\end{theorem}

\begin{proof}
$\qop{\Phi}\qket{I} = \qket{\Phi(I)} = \qket{I} = 1 \cdot \qket{I}$.
\end{proof}

\begin{definition}[Quantum Fixed Point Projector]
\label{def:fixed-projector}
The \textbf{fixed point projector} for quantum fractal $\qop{\Phi}$ is:
\[
\Projector_{\mathrm{fix}}(\qop{\Phi}) = \sum_{I : \Phi(I) = I} \qket{I}\qbra{I}
\]
This projects onto the subspace spanned by classical fixed points.
\end{definition}

\begin{proposition}[Eigenfractal Superpositions]
\label{prop:eigenfractal-super}
Any superposition of classical fixed points is an eigenfractal with eigenvalue 1:
\[
\qket{\psi_{\mathrm{fix}}} = \sum_{I : \Phi(I) = I} c_I \qket{I} \implies \qop{\Phi}\qket{\psi_{\mathrm{fix}}} = \qket{\psi_{\mathrm{fix}}}
\]
The eigenfractal subspace $\mathcal{E}_1(\qop{\Phi})$ has dimension equal to the number of classical fixed points.
\end{proposition}

\subsection{Non-Classical Eigenfractals}

\begin{theorem}[Cycle Eigenfractals]
\label{thm:cycle-eigenfractals}
Suppose $\Phi$ has a cycle of length $m$: $I_1 \to I_2 \to \cdots \to I_m \to I_1$. Then the states:
\[
\qket{\psi_r} = \frac{1}{\sqrt{m}} \sum_{j=1}^{m} e^{2\pi i r j / m} \qket{I_j}, \quad r = 0, 1, \ldots, m-1
\]
are eigenfractals of $\qop{\Phi}^m$ with eigenvalue 1, and eigenfractals of $\qop{\Phi}$ with eigenvalue $e^{2\pi i r / m}$.
\end{theorem}

\begin{proof}
For single application:
\begin{align*}
\qop{\Phi}\qket{\psi_r} &= \frac{1}{\sqrt{m}} \sum_{j=1}^{m} e^{2\pi i r j / m} \qket{\Phi(I_j)} \\
&= \frac{1}{\sqrt{m}} \sum_{j=1}^{m} e^{2\pi i r j / m} \qket{I_{j+1 \mod m}} \\
&= \frac{1}{\sqrt{m}} \sum_{k=1}^{m} e^{2\pi i r (k-1) / m} \qket{I_k} \quad (\text{substituting } k = j+1) \\
&= e^{-2\pi i r / m} \qket{\psi_r}
\end{align*}
For $\qop{\Phi}^m$: $(\qop{\Phi}^m)\qket{\psi_r} = (e^{-2\pi i r / m})^m \qket{\psi_r} = \qket{\psi_r}$.
\end{proof}

\begin{corollary}[Spectrum from Cycles]
\label{cor:spectrum-cycles}
The spectrum of $\qop{\Phi}$ includes:
\begin{itemize}
  \item $\lambda = 1$ with multiplicity $\geq$ (number of fixed points)
  \item $\lambda = e^{2\pi i r / m}$ for $r = 0, \ldots, m-1$ for each $m$-cycle
  \item $\lambda = 0$ with multiplicity $= 40 - \mathrm{rank}(\qop{\Phi})$
\end{itemize}
\end{corollary}

\begin{definition}[Coherent Cycle State]
\label{def:coherent-cycle}
The \textbf{coherent cycle state} for an $m$-cycle is:
\[
\qket{\psi_{\mathrm{cycle}}} = \frac{1}{\sqrt{m}} \sum_{j=1}^{m} \qket{I_j}
\]
This is the $r = 0$ eigenfractal with eigenvalue 1 under $\qop{\Phi}^m$.
\end{definition}

\subsection{Spectral Decomposition of Quantum Fractals}

\begin{theorem}[Quantum Fractal Spectral Theorem]
\label{thm:qfractal-spectral}
Every quantum fractal operator $\qop{\Phi}$ admits a spectral decomposition:
\[
\qop{\Phi} = \sum_{\lambda \in \Spectrum(\qop{\Phi})} \lambda \, \Projector_\lambda^{(\Phi)}
\]
where $\Projector_\lambda^{(\Phi)}$ is the projector onto the generalized eigenspace for eigenvalue $\lambda$.

For the case where $\qop{\Phi}$ is a normal operator ($[\qop{\Phi}, \qop{\Phi}^\dagger] = 0$):
\[
\qop{\Phi} = \sum_{\lambda} \lambda \, \Projector_\lambda, \quad \Projector_\lambda \Projector_{\lambda'} = \delta_{\lambda\lambda'} \Projector_\lambda
\]
\end{theorem}

\begin{definition}[Fractal Spectral Measure]
\label{def:fractal-spectral-measure}
The \textbf{spectral measure} of state $\qket{\psi}$ with respect to quantum fractal $\qop{\Phi}$ is:
\[
\mu_\psi(\lambda) = \|\Projector_\lambda^{(\Phi)}\qket{\psi}\|^2
\]
This gives the probability of ``collapsing'' to eigenspace $\lambda$ upon fractal measurement.
\end{definition}

%% ----------------------------------------------------------------------------
\section{Quantum Fractal Attractors}
\label{sec:quantum-attractors}

This section develops the quantum analogue of IFS (Iterated Function System) attractors.

\subsection{Quantum Iterated Function Systems}

\begin{definition}[Quantum IFS]
\label{def:quantum-ifs}
A \textbf{quantum iterated function system (QIFS)} on $\Hspace_{\Ideas}$ is a pair $(\{T_k\}_{k=1}^K, \{p_k\}_{k=1}^K)$ where:
\begin{enumerate}
  \item $T_k : \Hspace_{\Ideas} \to \Hspace_{\Ideas}$ are quantum operations (completely positive maps)
  \item $p_k \geq 0$ with $\sum_k p_k = 1$ are selection probabilities
\end{enumerate}
The associated \textbf{QIFS channel} is:
\[
\mathcal{T}(\rho) = \sum_{k=1}^{K} p_k \, T_k(\rho)
\]
\end{definition}

\begin{definition}[Noetic QIFS]
\label{def:noetic-qifs}
The \textbf{noetic QIFS} uses quantum noetic operators:
\[
\mathcal{T}_{\mathrm{noetic}}(\rho) = \sum_{k=1}^{22} p_k \, \qop{\nu}_k \rho \qop{\nu}_k^\dagger
\]
with probabilities $p_k$ over the 22 noetics.
\end{definition}

\begin{theorem}[QIFS Channel Properties]
\label{thm:qifs-channel}
The noetic QIFS channel $\mathcal{T}_{\mathrm{noetic}}$ is:
\begin{enumerate}
  \item \textbf{Completely positive}: Maps positive operators to positive operators
  \item \textbf{Trace non-increasing}: $\Trace(\mathcal{T}(\rho)) \leq \Trace(\rho)$
  \item \textbf{Trace-preserving} iff $\sum_k p_k \qop{\nu}_k^\dagger \qop{\nu}_k = \mathbf{I}$
\end{enumerate}
\end{theorem}

\subsection{Fixed-Point Density Matrices}

\begin{definition}[Quantum Attractor]
\label{def:quantum-attractor}
A \textbf{quantum attractor} for QIFS channel $\mathcal{T}$ is a density matrix $\rho_*$ satisfying:
\[
\mathcal{T}(\rho_*) = \rho_*
\]
This is the quantum analogue of the classical IFS attractor.
\end{definition}

\begin{theorem}[Quantum Attractor Existence]
\label{thm:quantum-attractor-existence}
Every trace-preserving noetic QIFS has at least one quantum attractor $\rho_*$.
\end{theorem}

\begin{proof}
The set of density matrices $\mathcal{D}(\Hspace_{\Ideas})$ is compact and convex. The channel $\mathcal{T}$ is a continuous affine map. By the Brouwer fixed point theorem (or Schauder's theorem), a fixed point exists.
\end{proof}

\begin{definition}[Iterated Attractor Convergence]
\label{def:attractor-convergence}
The \textbf{iterated attractor sequence} starting from $\rho_0$ is:
\[
\rho_n = \mathcal{T}^n(\rho_0) = \underbrace{\mathcal{T} \circ \mathcal{T} \circ \cdots \circ \mathcal{T}}_{n \text{ times}}(\rho_0)
\]
\end{definition}

\begin{theorem}[Attractor Convergence]
\label{thm:attractor-convergence}
If the noetic QIFS is \textbf{strictly contractive} (i.e., $\exists \gamma < 1$ such that $\|\mathcal{T}(\rho) - \mathcal{T}(\sigma)\|_1 \leq \gamma \|\rho - \sigma\|_1$), then:
\begin{enumerate}
  \item The quantum attractor $\rho_*$ is unique
  \item For any initial state $\rho_0$: $\lim_{n \to \infty} \mathcal{T}^n(\rho_0) = \rho_*$
  \item Convergence is exponential: $\|\rho_n - \rho_*\|_1 \leq \gamma^n \|\rho_0 - \rho_*\|_1$
\end{enumerate}
\end{theorem}

\begin{proof}
This is the Banach contraction mapping theorem applied to the complete metric space $(\mathcal{D}(\Hspace_{\Ideas}), \|\cdot\|_1)$.
\end{proof}

\subsection{Structure of Quantum Attractors}

\begin{proposition}[Attractor Purity]
\label{prop:attractor-purity}
The quantum attractor $\rho_*$ is:
\begin{enumerate}
  \item \textbf{Pure} ($\rho_* = \qket{\psi_*}\qbra{\psi_*}$) iff all noetic operators share a common eigenvector
  \item \textbf{Mixed} in general, even when starting from a pure state
\end{enumerate}
\end{proposition}

\begin{definition}[Attractor Support]
\label{def:attractor-support}
The \textbf{support} of quantum attractor $\rho_*$ is:
\[
\mathrm{supp}(\rho_*) = \mathrm{Im}(\rho_*) = \{\qket{\psi} : \qbra{\psi}\rho_*\qket{\psi} > 0\}
\]
This is the quantum analogue of the classical attractor set.
\end{definition}

\begin{theorem}[Attractor Support and Classical Attractor]
\label{thm:attractor-support-classical}
Let $A_{\mathrm{cl}}$ be the classical IFS attractor (the set of Ideas visited infinitely often). Then:
\[
\mathrm{supp}(\rho_*) = \mathrm{span}\{\qket{I} : I \in A_{\mathrm{cl}}\}
\]
The quantum attractor support is exactly the span of classical attractor Ideas.
\end{theorem}

\begin{proof}
\textbf{($\supseteq$)}: If $I \in A_{\mathrm{cl}}$, then for any $\epsilon > 0$, infinitely many iterations reach $I$. Thus $\qbra{I}\rho_*\qket{I} > 0$.

\textbf{($\subseteq$)}: If $I \notin A_{\mathrm{cl}}$, then only finitely many iterations can reach $I$. As $n \to \infty$, $\qbra{I}\rho_n\qket{I} \to 0$.
\end{proof}

\begin{example}[World Attractor]
\label{ex:world-attractor}
Consider a QIFS with noetics that all map to World D (Physical). The quantum attractor is:
\[
\rho_* = \frac{1}{10}\sum_{n=1}^{10} \qket{Dn}\qbra{Dn}
\]
the maximally mixed state over World D. This represents complete ``descent'' to the physical realm.
\end{example}

\begin{definition}[Entanglement in Attractor]
\label{def:attractor-entanglement}
For a bipartite QIFS on $\Hspace_{\Ideas}^{(2)}$, the \textbf{attractor entanglement} is:
\[
E_* = E(\rho_*)
\]
where $E$ is an entanglement measure (e.g., von Neumann entropy of reduced state).
\end{definition}

\begin{theorem}[Attractor Entanglement Bounds]
\label{thm:attractor-entanglement}
For a bipartite noetic QIFS with local operations (no entangling gates):
\[
E_* = 0
\]
The attractor is separable. Non-zero attractor entanglement requires non-local QIFS operations.
\end{theorem}

%% ----------------------------------------------------------------------------
\section{Decoherence and Fractal Dynamics}
\label{sec:decoherence-fractals}

This section analyzes how environmental decoherence affects quantum fractal evolution.

\subsection{Decoherence Channels}

\begin{definition}[Decoherence Channel]
\label{def:decoherence-channel}
A \textbf{decoherence channel} in the Idea basis is:
\[
\mathcal{D}(\rho) = \sum_{I \in \Ideas} \qket{I}\qbra{I} \rho \qket{I}\qbra{I} = \sum_I \rho_{II} \qket{I}\qbra{I}
\]
This eliminates all off-diagonal coherences, leaving only populations.
\end{definition}

\begin{definition}[Partial Decoherence]
\label{def:partial-decoherence}
The \textbf{$\gamma$-partial decoherence channel} interpolates between coherent and decohered:
\[
\mathcal{D}_\gamma(\rho) = (1 - \gamma)\rho + \gamma \mathcal{D}(\rho), \quad \gamma \in [0, 1]
\]
\begin{itemize}
  \item $\gamma = 0$: No decoherence (coherent evolution)
  \item $\gamma = 1$: Full decoherence (classical limit)
\end{itemize}
\end{definition}

\begin{proposition}[Decoherence Properties]
\label{prop:decoherence-props}
The decoherence channel $\mathcal{D}$ satisfies:
\begin{enumerate}
  \item \textbf{Idempotent}: $\mathcal{D}^2 = \mathcal{D}$
  \item \textbf{Trace-preserving}: $\Trace(\mathcal{D}(\rho)) = \Trace(\rho)$
  \item \textbf{Completely positive}: Kraus operators are $K_I = \qket{I}\qbra{I}$
  \item \textbf{Projection}: $\mathcal{D}$ projects onto the diagonal subalgebra
\end{enumerate}
\end{proposition}

\subsection{Decohered Quantum Fractals}

\begin{definition}[Decohered Quantum Fractal]
\label{def:decohered-qfractal}
The \textbf{decohered quantum fractal channel} combines quantum fractal evolution with decoherence:
\[
\mathcal{F}_\gamma(\rho) = \mathcal{D}_\gamma(\qop{\Phi} \rho \qop{\Phi}^\dagger)
\]
This models a quantum fractal in a noisy environment.
\end{definition}

\begin{theorem}[Classical Limit of Quantum Fractals]
\label{thm:quantum-frac-classical-limit}
In the full decoherence limit ($\gamma = 1$), the decohered quantum fractal reduces to classical fractal dynamics:
\[
\mathcal{F}_1(\rho) = \sum_I \left(\sum_J |\qbra{I}\qop{\Phi}\qket{J}|^2 \rho_{JJ}\right) \qket{I}\qbra{I}
\]
For a classical initial state $\rho = \qket{I}\qbra{I}$:
\[
\mathcal{F}_1(\qket{I}\qbra{I}) = \qket{\Phi(I)}\qbra{\Phi(I)}
\]
recovering the classical fractal map $I \mapsto \Phi(I)$.
\end{theorem}

\begin{proof}
With full decoherence:
\begin{align*}
\mathcal{F}_1(\rho) &= \mathcal{D}(\qop{\Phi} \rho \qop{\Phi}^\dagger) \\
&= \sum_I \qket{I}\qbra{I} \qop{\Phi} \rho \qop{\Phi}^\dagger \qket{I}\qbra{I} \\
&= \sum_I \qbra{I}\qop{\Phi} \rho \qop{\Phi}^\dagger\qket{I} \cdot \qket{I}\qbra{I}
\end{align*}
For $\rho = \qket{J}\qbra{J}$:
\[
\qbra{I}\qop{\Phi}\qket{J}\qbra{J}\qop{\Phi}^\dagger\qket{I} = |\qbra{I}\qop{\Phi}\qket{J}|^2 = \delta_{I, \Phi(J)}
\]
since $\qop{\Phi}\qket{J} = \qket{\Phi(J)}$.
\end{proof}

\begin{corollary}[Decoherence Kills Quantum Interference]
\label{cor:decoherence-interference}
For an initial superposition $\qket{\psi} = \frac{1}{\sqrt{2}}(\qket{I} + \qket{J})$ with $\Phi(I) = \Phi(J) = K$:
\begin{itemize}
  \item \textbf{Coherent} ($\gamma = 0$): $\qop{\Phi}\qket{\psi} = \sqrt{2}\qket{K}$ (constructive interference)
  \item \textbf{Decohered} ($\gamma = 1$): $\mathcal{F}_1(\qket{\psi}\qbra{\psi}) = \qket{K}\qbra{K}$ (same final state, but no interference in dynamics)
\end{itemize}
\end{corollary}

\subsection{Environment-Induced Superselection}

\begin{definition}[Superselection Sector]
\label{def:superselection}
A \textbf{superselection sector} is a subspace $\mathcal{H}_s \subset \Hspace_{\Ideas}$ such that:
\begin{enumerate}
  \item Superpositions within $\mathcal{H}_s$ are allowed
  \item Superpositions between different sectors $\mathcal{H}_s$ and $\mathcal{H}_{s'}$ are forbidden (rapidly decohere)
\end{enumerate}
\end{definition}

\begin{theorem}[World Superselection]
\label{thm:world-superselection}
Under strong World-basis decoherence, the Worlds $\{A, B, C, D\}$ become superselection sectors:
\[
\Hspace_{\Ideas} = \Hspace_A \oplus \Hspace_B \oplus \Hspace_C \oplus \Hspace_D
\]
\begin{enumerate}
  \item Superpositions within a World (e.g., $\alpha\qket{A1} + \beta\qket{A2}$) are stable
  \item Cross-World superpositions (e.g., $\alpha\qket{A1} + \beta\qket{B1}$) rapidly decohere
\end{enumerate}
\end{theorem}

\begin{proof}
Define the World decoherence channel:
\[
\mathcal{D}_W(\rho) = \sum_{W \in \{A,B,C,D\}} \Projector_W \rho \Projector_W
\]
This preserves within-World coherences but eliminates between-World coherences. Under repeated application or continuous coupling to an environment that measures World, cross-World coherences decay exponentially.
\end{proof}

\begin{definition}[Einselection]
\label{def:einselection}
\textbf{Environment-induced superselection (einselection)} occurs when the environment preferentially measures certain observables, selecting \textbf{pointer states} that are robust against decoherence.
\end{definition}

\begin{proposition}[Pointer States for Quantum Fractals]
\label{prop:pointer-states}
For a quantum fractal $\qop{\Phi}$ coupled to an environment that measures the output:
\begin{enumerate}
  \item \textbf{Pointer states}: Classical Idea basis states $\{\qket{I}\}$
  \item \textbf{Stable fractals}: Those whose output is a classical mixture of Ideas
  \item \textbf{Unstable fractals}: Those requiring coherent superpositions
\end{enumerate}
\end{proposition}

\subsection{Decoherence Time Scales}

\begin{definition}[Decoherence Time]
\label{def:decoherence-time}
The \textbf{decoherence time} $\tau_D$ for a quantum fractal is the time scale over which off-diagonal elements decay:
\[
\rho_{IJ}(t) \approx \rho_{IJ}(0) \cdot e^{-t/\tau_D} \quad \text{for } I \neq J
\]
\end{definition}

\begin{theorem}[Fractal-Decoherence Competition]
\label{thm:fractal-decoherence}
Let $\tau_F$ be the fractal application time (time to apply one noetic operation). The quantum-classical transition occurs when:
\begin{itemize}
  \item $\tau_D \gg \tau_F$: \textbf{Quantum regime}---coherent fractal evolution
  \item $\tau_D \ll \tau_F$: \textbf{Classical regime}---decohered (classical) fractal
  \item $\tau_D \sim \tau_F$: \textbf{Intermediate regime}---partially coherent
\end{itemize}
\end{theorem}

%% ----------------------------------------------------------------------------
\section{Quantum-Classical Correspondence}
\label{sec:quantum-classical-correspondence}

This section establishes the rigorous connection between quantum and classical fractal dynamics.

\subsection{Ehrenfest Theorem for Noetic Fractals}

\begin{theorem}[Noetic Ehrenfest Theorem]
\label{thm:noetic-ehrenfest}
For a quantum state $\rho$ evolving under quantum fractal $\qop{\Phi}$, the expectation value of World observable $\qop{W}$ satisfies:
\[
\langle \qop{W} \rangle_{\qop{\Phi}\rho\qop{\Phi}^\dagger} = \langle \qop{\Phi}^\dagger \qop{W} \qop{\Phi} \rangle_\rho
\]
In the classical limit (diagonal $\rho$ in Idea basis), this reduces to:
\[
\langle \qop{W} \rangle_{\Phi(\rho)} = \sum_I p_I \cdot w(\mathrm{World}(\Phi(I)))
\]
the classical expectation of World after fractal application.
\end{theorem}

\begin{proof}
For quantum evolution:
\[
\langle \qop{W} \rangle_{\qop{\Phi}\rho\qop{\Phi}^\dagger} = \Trace(\qop{W} \qop{\Phi} \rho \qop{\Phi}^\dagger) = \Trace(\qop{\Phi}^\dagger \qop{W} \qop{\Phi} \rho) = \langle \qop{\Phi}^\dagger \qop{W} \qop{\Phi} \rangle_\rho
\]
For classical state $\rho = \sum_I p_I \qket{I}\qbra{I}$:
\begin{align*}
\langle \qop{W} \rangle_{\qop{\Phi}\rho\qop{\Phi}^\dagger} &= \sum_I p_I \qbra{I}\qop{\Phi}^\dagger \qop{W} \qop{\Phi}\qket{I} \\
&= \sum_I p_I \qbra{\Phi(I)}\qop{W}\qket{\Phi(I)} \\
&= \sum_I p_I \cdot w(\mathrm{World}(\Phi(I)))
\end{align*}
\end{proof}

\begin{corollary}[World Descent Under Fractals]
\label{cor:world-descent}
If fractal $\Phi$ always maps to a World $\geq$ the input World (in ACBE order), then:
\[
\langle \qop{W} \rangle_{\qop{\Phi}\rho\qop{\Phi}^\dagger} \geq \langle \qop{W} \rangle_\rho
\]
Quantum fractals preserve the classical ACBE descent property in expectation.
\end{corollary}

\subsection{Correspondence Conditions}

\begin{definition}[Classical State]
\label{def:classical-state}
A quantum state $\rho$ is \textbf{classical} (with respect to the Idea basis) if:
\[
\rho = \sum_I p_I \qket{I}\qbra{I}
\]
i.e., $\rho$ is diagonal in the Idea basis with no coherences.
\end{definition}

\begin{definition}[Classicality Measure]
\label{def:classicality}
The \textbf{classicality} of state $\rho$ is:
\[
C(\rho) = 1 - \frac{\|\rho - \mathcal{D}(\rho)\|_1}{\|\rho\|_1}
\]
where $\mathcal{D}$ is the Idea-basis decoherence channel. $C(\rho) = 1$ for classical states, $C(\rho) < 1$ for quantum states.
\end{definition}

\begin{proposition}[Classical State Preservation]
\label{prop:classical-preservation}
Quantum fractal $\qop{\Phi}$ preserves classical states:
\[
\rho \text{ classical} \implies \qop{\Phi}\rho\qop{\Phi}^\dagger \text{ classical}
\]
The quantum fractal acts as a classical fractal on classical inputs.
\end{proposition}

\begin{proof}
For $\rho = \sum_I p_I \qket{I}\qbra{I}$:
\[
\qop{\Phi}\rho\qop{\Phi}^\dagger = \sum_I p_I \qket{\Phi(I)}\qbra{\Phi(I)}
\]
which is diagonal in the Idea basis (classical).
\end{proof}

\subsection{Main Correspondence Theorem}

\begin{theorem}[Quantum-Classical Fractal Correspondence]
\label{thm:qc-correspondence}
\textbf{(Main Theorem)}

Let $\qop{\Phi}$ be the quantum fractal operator for classical fractal $\Phi$. The following conditions are equivalent:
\begin{enumerate}
  \item \textbf{Quantum dynamics reduces to classical}: For all classical states $\rho$, the quantum evolution $\qop{\Phi}\rho\qop{\Phi}^\dagger$ equals the classical evolution $\Phi_*(\rho)$

  \item \textbf{Basis preservation}: $\qop{\Phi}\qket{I} = \qket{\Phi(I)}$ for all Idea basis states $\qket{I}$

  \item \textbf{Expectation correspondence}: For any observable $\qop{O}$ diagonal in the Idea basis and any classical state $\rho$:
  \[
  \langle \qop{O} \rangle_{\qop{\Phi}\rho\qop{\Phi}^\dagger} = \langle O \circ \Phi \rangle_\rho
  \]
  where $(O \circ \Phi)(I) = O(\Phi(I))$ is the classical composition

  \item \textbf{Fixed point correspondence}: $I$ is a quantum eigenfractal (with eigenvalue 1) iff $I$ is a classical fixed point:
  \[
  \qop{\Phi}\qket{I} = \qket{I} \iff \Phi(I) = I
  \]

  \item \textbf{Attractor correspondence}: The quantum attractor support equals the span of the classical attractor:
  \[
  \mathrm{supp}(\rho_*) = \mathrm{span}(A_{\mathrm{cl}})
  \]
\end{enumerate}

Moreover, quantum effects (interference, entanglement) arise precisely when:
\begin{enumerate}
  \item[(Q1)] The initial state $\rho$ has coherences (off-diagonal elements)
  \item[(Q2)] The fractal is applied to entangled multi-system states
  \item[(Q3)] Measurements are performed in non-Idea bases
\end{enumerate}
\end{theorem}

\begin{proof}
\textbf{(1) $\Leftrightarrow$ (2)}: By linearity, if $\qop{\Phi}\qket{I} = \qket{\Phi(I)}$ for all $I$, then for classical $\rho = \sum_I p_I \qket{I}\qbra{I}$:
\[
\qop{\Phi}\rho\qop{\Phi}^\dagger = \sum_I p_I \qket{\Phi(I)}\qbra{\Phi(I)} = \Phi_*(\rho)
\]
Conversely, if (1) holds, apply to $\rho = \qket{I}\qbra{I}$ to get (2).

\textbf{(2) $\Leftrightarrow$ (3)}: For diagonal observable $\qop{O} = \sum_I o_I \qket{I}\qbra{I}$ and classical $\rho$:
\begin{align*}
\langle \qop{O} \rangle_{\qop{\Phi}\rho\qop{\Phi}^\dagger} &= \sum_I p_I o_{\Phi(I)} = \sum_I p_I (O \circ \Phi)(I) = \langle O \circ \Phi \rangle_\rho
\end{align*}

\textbf{(2) $\Leftrightarrow$ (4)}: $\qop{\Phi}\qket{I} = \qket{I}$ iff $\qket{\Phi(I)} = \qket{I}$ iff $\Phi(I) = I$.

\textbf{(2) $\Rightarrow$ (5)}: Follows from Theorem~\ref{thm:attractor-support-classical}.

\textbf{(Q1)--(Q3)}: These are the only ways quantum effects enter:
\begin{itemize}
  \item (Q1): Coherences allow interference in the fractal output
  \item (Q2): Entanglement creates correlations not possible classically
  \item (Q3): Non-Idea-basis measurements reveal coherences/entanglement
\end{itemize}
\end{proof}

\begin{corollary}[Classical TKS as Decoherence Limit]
\label{cor:classical-tks}
Classical TKS fractal theory (v7.0--v7.2) is recovered from quantum TKS in the limit of:
\begin{enumerate}
  \item Complete decoherence in the Idea basis ($\gamma = 1$)
  \item No entanglement between systems
  \item Measurements only in the Idea basis
\end{enumerate}
\end{corollary}

\begin{theorem}[Quantum Correction Terms]
\label{thm:quantum-corrections}
For a state with small coherences $\rho = \rho_{\mathrm{cl}} + \epsilon \rho_{\mathrm{coh}}$ where $\rho_{\mathrm{cl}}$ is classical and $\rho_{\mathrm{coh}}$ contains only off-diagonal terms:
\[
\qop{\Phi}\rho\qop{\Phi}^\dagger = \Phi_*(\rho_{\mathrm{cl}}) + \epsilon \cdot \Delta_\Phi(\rho_{\mathrm{coh}}) + O(\epsilon^2)
\]
where $\Delta_\Phi(\rho_{\mathrm{coh}}) = \qop{\Phi}\rho_{\mathrm{coh}}\qop{\Phi}^\dagger$ is the \textbf{quantum correction} capturing interference effects.
\end{theorem}

\begin{definition}[Quantum Fractal Interference]
\label{def:quantum-interference}
The \textbf{quantum fractal interference} for state $\qket{\psi} = \sum_I c_I \qket{I}$ is:
\[
\mathcal{I}_\Phi(\psi) = \|\qop{\Phi}\qket{\psi}\|^2 - \sum_I |c_I|^2
\]
This measures the deviation from classical probability due to interference.
\end{definition}

\begin{proposition}[Interference Bounds]
\label{prop:interference-bounds}
The quantum fractal interference satisfies:
\begin{enumerate}
  \item $\mathcal{I}_\Phi(\psi) = 0$ if $\qket{\psi}$ is a basis state or if all $\Phi(I)$ are distinct
  \item $|\mathcal{I}_\Phi(\psi)| \leq 2\sum_{I \neq J} |c_I||c_J|$ (coherence bound)
  \item Maximum interference occurs when multiple inputs map to the same output with appropriate phases
\end{enumerate}
\end{proposition}

%% ----------------------------------------------------------------------------
%% HANDOFF NOTE
%% ----------------------------------------------------------------------------
\begin{tcolorbox}[colback=green!5!white,colframe=green!50!black,title=\textbf{Handoff: Next Steps},breakable]
This completes Chapter 7 on Quantum Noetic Fractals. The subsequent chapters should:
\begin{itemize}
  \item \textbf{Chapter 8 (Applications)}: Apply quantum fractals to RPM goal achievement; quantum error correction for noetic states; quantum algorithms for fractal fixed points
  \item \textbf{Chapter 9 (Advanced Topics)}: Topological phases in quantum TKS; quantum groups and Hopf algebra structure
\end{itemize}

Key results established in this chapter:
\begin{enumerate}
  \item \textbf{Quantum fractal operators}: $\qop{\Phi} = \qop{\nu}_{k_n} \cdots \qop{\nu}_{k_1}$ (Sec.~\ref{sec:quantum-fractal-operators})
  \item \textbf{Transfinite quantum fractals}: $\qop{\Phi}^{(\alpha)}$ for ordinal $\alpha$, SOT limits at limit ordinals (Sec.~\ref{sec:transfinite-quantum-fractals})
  \item \textbf{$\omega$-quantum fractals}: $\qop{\Phi}^{(\omega)}_k$ as projector onto fixed point subspace (Theorem~\ref{thm:omega-convergence})
  \item \textbf{Quantum eigenfractals}: Eigenstates from fixed points and cycles (Sec.~\ref{sec:quantum-eigenfractals})
  \item \textbf{Cycle eigenfractals}: $m$-cycle yields eigenvalues $e^{2\pi i r/m}$ (Theorem~\ref{thm:cycle-eigenfractals})
  \item \textbf{Quantum IFS attractors}: Fixed-point density matrices (Sec.~\ref{sec:quantum-attractors})
  \item \textbf{Attractor-classical correspondence}: $\mathrm{supp}(\rho_*) = \mathrm{span}(A_{\mathrm{cl}})$ (Theorem~\ref{thm:attractor-support-classical})
  \item \textbf{Decoherence}: Classical limit as $\gamma \to 1$ (Sec.~\ref{sec:decoherence-fractals})
  \item \textbf{World superselection}: Worlds as superselection sectors under decoherence (Theorem~\ref{thm:world-superselection})
  \item \textbf{Main Correspondence Theorem}: Theorem~\ref{thm:qc-correspondence} unifies quantum and classical fractal semantics
\end{enumerate}

\textbf{Connection to v7.2}: Quantizes Chapter 1 (Transfinite Fractals), extending Definitions~\ref{def:transfinite-fractal}, \ref{def:omega-eigenfractal}, \ref{def:transfinite-eval}, Theorems~\ref{thm:omega-eigen-idem}, \ref{thm:transfinite-fixed}, \ref{thm:stabilization}.

\textbf{Connection to Math Track}: Quantum fractals provide operator-theoretic realization of fractal calculus; spectral theory connects to measure theory (v7.2 Chapter 2).

\textbf{Connection to Chapters 1--6}: Builds on Hilbert spaces (Ch.~1), quantum noetic operators (Ch.~3), density matrices (Ch.~2), channels (Ch.~4), entanglement (Ch.~5), measurement (Ch.~6).

\textbf{Main Unifying Result}: Theorem~\ref{thm:qc-correspondence} establishes that quantum TKS reduces to classical TKS precisely when: (1) states are classical (diagonal), (2) no entanglement, (3) Idea-basis measurements. Quantum effects arise from coherences, entanglement, and non-classical measurements.
\end{tcolorbox}

